% concatenated from main.tex
\documentclass[11pt,a4paper]{article}
\usepackage[utf8]{inputenc}
\usepackage[T1]{fontenc}
\usepackage[english]{babel}
\usepackage{lmodern}
\usepackage{amsmath,amssymb,amsthm,mathtools,mathrsfs}
\usepackage[a4paper,left=2.05cm,right=2.05cm,top=1.9cm,bottom=2.1cm]{geometry}
\usepackage{microtype}
\microtypesetup{expansion=false}
\usepackage{booktabs,array,longtable}
\usepackage{enumitem}
\usepackage{xcolor}
\usepackage{graphicx}
\usepackage[hidelinks,bookmarks=true,bookmarksopen=true,pdfencoding=auto]{hyperref}
\usepackage{aliascnt}
\usepackage{needspace}
\usepackage{tikz}
\usetikzlibrary{arrows.meta,positioning,calc,fit}
\allowdisplaybreaks
\emergencystretch=3em

\newtheorem{theorem}{Theorem}[section]
\newaliascnt{proposition}{theorem}
\newtheorem{proposition}[proposition]{Proposition}
\aliascntresetthe{proposition}
\newaliascnt{lemma}{theorem}
\newtheorem{lemma}[lemma]{Lemma}
\aliascntresetthe{lemma}
\newaliascnt{corollary}{theorem}
\newtheorem{corollary}[corollary]{Corollary}
\aliascntresetthe{corollary}
\newaliascnt{definition}{theorem}
\newtheorem{definition}[definition]{Definition}
\aliascntresetthe{definition}
\newaliascnt{remark}{theorem}
\newtheorem{remark}[remark]{Remark}
\aliascntresetthe{remark}
\newaliascnt{example}{theorem}
\newtheorem{example}[example]{Example}
\aliascntresetthe{example}
\usepackage[nameinlink,noabbrev]{cleveref}
\crefname{theorem}{Theorem}{Theorems}\Crefname{theorem}{Theorem}{Theorems}
\crefname{proposition}{Proposition}{Propositions}\Crefname{proposition}{Proposition}{Propositions}
\crefname{lemma}{Lemma}{Lemmas}\Crefname{lemma}{Lemma}{Lemmas}
\crefname{corollary}{Corollary}{Corollaries}\Crefname{corollary}{Corollary}{Corollaries}
\crefname{definition}{Definition}{Definitions}\Crefname{definition}{Definition}{Definitions}
\crefname{remark}{Remark}{Remarks}\Crefname{remark}{Remark}{Remarks}
\crefname{example}{Example}{Examples}\Crefname{example}{Example}{Examples}

\newcommand{\RR}{\mathbb R}
\newcommand{\CC}{\mathbb C}
\newcommand{\HH}{\mathbb H}
\newcommand{\OO}{\mathbb O}
\newcommand{\Gtwo}{G_2}
\newcommand{\gtwo}{\mathfrak g_2}
\newcommand{\SO}{\mathrm{SO}}
\newcommand{\SU}{\mathrm{SU}}
\newcommand{\Sp}{\mathrm{Sp}}
\newcommand{\Aut}{\operatorname{Aut}}
\newcommand{\Der}{\operatorname{Der}}
\newcommand{\Stab}{\operatorname{Stab}}
\newcommand{\Lie}{\operatorname{Lie}}
\newcommand{\imO}{\operatorname{Im}\OO}
\newcommand{\Ad}{\operatorname{Ad}}
\newcommand{\Span}{\operatorname{span}}
\newcommand{\braket}[2]{\langle #1,#2\rangle}

\hypersetup{pdftitle={Quaternionic Reflections and Lie Generation in g2},pdfauthor={Santiago Pineda Montoya, Johan H. Rua Munoz, Borut Jurcic Zlobec},pdfkeywords={octonions, G2, Lie generation, quaternionic subalgebras, compact symmetric pairs, nongenerating locus}}

\title{\textbf{Quaternionic Reflections and Lie Generation in $\mathfrak g_2$}}
\author{Santiago Pineda Montoya\thanks{Universidad Nacional de Colombia, Medell\'{i}n, Colombia. Corresponding author: \texttt{sapinedamo@unal.edu.co}.}
\and Johan H. R\'{u}a Mu\~{n}oz\thanks{Instituto de F\'{i}sica, Facultad de Ciencias Exactas y Naturales, Universidad de Antioquia, Medell\'{i}n, Colombia. Email: \texttt{heriberto.rua@udea.edu.co}.}
\and Borut Jur\v{c}i\v{c} Zlobec\thanks{Independent researcher, Ljubljana, Slovenia. Email: \texttt{borut.jurcic.zlobec@gmail.com}.}}
\date{September 2026}

\begin{document}
\maketitle

\begin{abstract}
Let $\OO$ be the real octonion division algebra and $\mathfrak g_2=\Der(\OO)$.
A quaternionic reflection $T\in G_2=\Aut(\OO)$ fixes a quaternionic subalgebra pointwise and negates its orthogonal complement.
For the associated compact symmetric pair
$\mathfrak g_2=\mathfrak k_T\oplus\mathfrak m_T$, we give an explicit factored polynomial criterion for generation by an independently chosen even element and odd element.
Equivalently, for the reflected pair $(Z,TZT)$ with $Z\in\mathfrak g_2$, the criterion is
$\mathcal U_{38}=\mathcal R_{30}\mathcal H_8>0$.
The degree-$30$ Gram determinant detects reducibility on $\operatorname{Im}\OO$, while the degree-$8$ factor detects generated algebras of dimension at most three.
A principal $\mathfrak{su}(2)$ is the remaining irreducible proper possibility.
Such a reflected principal closure is realizable for a nonzero $Z$ if and only if its positive rotation frequencies have ratio $1:2:3$; for each such $Z$, all realizing reflections are parametrized by a two-torus and an open interval.

A different restriction arises when $Z$ must lie in the stabilizer algebra $\mathfrak k_A\cong\mathfrak{so}(4)$ of a fixed quaternionic subalgebra $H_A$.
For every reflection moving $H_A$, a fixed list of six elements of $\mathfrak k_A$ contains a generating choice; reflections preserving $H_A$ admit none.
For one explicit relative configuration $(\mathfrak k_A,S)$, we describe the entire nongenerating set by four geometric families and by six irredundant systems of scalar equations.
We determine its real dimension and all generated algebras of dimension at most three; a homogeneous degree-$34$ polynomial packages the six tests.
The degree is not asserted to be canonical or minimal.
Two explicit local coordinate maps use short words in a finite-time pulse and one reflection, whose generated subgroup is dense in $G_2$.
\end{abstract}

\noindent\textbf{Keywords:} octonions; exceptional Lie algebra $\mathfrak g_2$; Lie generation; quaternionic subalgebras; compact symmetric pairs; nongenerating loci; maximal subalgebras.

\section{Introduction}
The real octonion division algebra $\OO$ contains associative quaternionic subalgebras.  We call such a unital subalgebra $H$ a \emph{quaternionic context} and write $P_H=H\cap V$, where $V=\operatorname{Im}\OO$.
Let $G_2=\Aut(\OO)$ and $\mathfrak g_2=\Der(\OO)$ be the compact real group and its Lie algebra.
The stabilizer of a context is $K_H=\{g\in G_2:g(H)=H\}$, with Lie algebra $\mathfrak k_H\cong\mathfrak{so}(4)$.
A local generator relative to a fixed context $H_A$ is an element $Z\in\mathfrak k_A:=\mathfrak k_{H_A}$, and its pulse is $e^{tZ}$, $t\in\RR$.
Locality means preservation of $H_A$, not identity on its orthogonal complement and not locality in a tensor-product model.

The quaternionic reflection associated with $H$ is
\[
\sigma_H(h+y)=h-y,\qquad h\in H,\quad y\in H^\perp.
\]
It is an octonion automorphism by the stabilizer model in \Cref{prop:stabilizer}.
On $V$ it has matrix $\operatorname{diag}(I_3,-I_4)$ in an adapted orthonormal basis, where $I_n$ denotes the identity matrix.  Thus it is not a hyperplane reflection.
We write $I_E$ for the identity on a vector space $E$.
The reflecting context $H$ need not be $H_A$: the condition $T(H_A)\ne H_A$ concerns the local context, although $T=\sigma_H$ fixes its own reflecting context pointwise.

The available pulse and switch produce
\[
Te^{tZ}T^{-1}=e^{tTZT},\qquad T^{-1}=T.
\]
Here $TZT$ is composition of real-linear operators.  For two Lie-algebra elements, $\Lie\{U,V\}$ denotes the smallest real Lie subalgebra containing them; this is not a topological closure.
There are two distinct parameter spaces.  With $T$ fixed, put
\[
A_+=\frac{Z+TZT}{2},\qquad A_-=\frac{Z-TZT}{2}.
\]
If $Z$ ranges over all of $\mathfrak g_2$, then $Z\mapsto(A_+,A_-)$ is a linear bijection onto $\mathfrak k_T\oplus\mathfrak m_T$, the even and odd eigenspaces of $\Ad_T$.
The inverse is $(A_+,A_-)\mapsto A_++A_-$, and
$\Lie\{Z,TZT\}=\Lie\{A_+,A_-\}$.
Thus the unrestricted reflected-pair problem is precisely two-generation with one independently chosen element of each parity, not an additional coupling of those two elements.
The local constraint $Z\in\mathfrak k_A$ is different: it imposes
\[
(A_+,A_-)\in\mathcal W_A(T):=
\{(A_+,A_-)\in\mathfrak k_T\oplus\mathfrak m_T:A_++A_-\in\mathfrak k_A\},
\]
a six-dimensional linear subspace of the fourteen-dimensional parity parameter space.

Classical two-generation results allow independently chosen elements in a semisimple or simple Lie algebra \cite{Kuranishi1951,AlbuquerqueSilvaLeite1989,Bois2009}; explicit constructions include $G_2$ \cite{DetinkoDeGraaf2020}.
Reflected pairs have a direct antecedent in Bauer, Levaillant, and Freedman, who use a generator and its swap conjugate in $\mathfrak{su}(4)$ \cite{BauerLevaillantFreedman2014}.
General openness and genericity statements for generating tuples and conjugate subgroups are treated by Chirvasitu \cite{Chirvasitu2021}.
The context space $G_2/\SO(4)$ and its stabilizers are classical \cite{HarveyLawson1982,Cacciatori2005,KnarrStroppel2025}; invariant associative planes and principal subalgebras are described by Draper and Mart\'in-Gonz\'alez \cite{DraperMartin2025}.
These are structural inputs, not novelty claims of the present results.

The existence of some polynomial test for generation is also formal: finitely many brackets span the generated algebra, so the sum of squares of their maximal minors vanishes exactly on the nongenerating set.
The specific uniform construction in \Cref{thm:universal-locus} instead uses
\[
\mathcal U_{38}(T,Z)=\mathcal R_{30}(T,Z)\mathcal H_8(T,Z).
\]
The first factor is a $15\times15$ Gram determinant for symmetric operators commuting with $T$ and detects reducibility on $V$.  The second is a degree-eight expression in two double commutators and vanishes precisely for generated algebras of dimension at most three.
The compact maximal-subalgebra list makes these tests exhaustive.
The point is their explicit construction and their interpretation, not the bare existence of an algebraic equation or an assertion of optimal degree or faster evaluation.
\Cref{cor:universal-spectral} characterizes every nonzero pulse that can give the remaining principal closure: exactly those with positive frequencies in the ratio $1:2:3$.

For the actual local constraint, \Cref{thm:universal-reflections} gives a fixed list of six elements of $\mathfrak k_A$ that contains a generating choice for every $T$ with $T(H_A)\ne H_A$.
The reverse obstruction is immediate: if $T$ preserves $H_A$, both generators remain in $\mathfrak k_A$.
The substantial assertion is the uniform sufficient list, obtained from a forced context and an affine obstruction of dimension at most one.
The chosen member may depend on $T$; one member working for every reflection is not asserted.

For the explicit reflection $S$ in \Cref{sec:models}, set
\[
L_Z=\Lie\{Z,\Ad_S Z\},\qquad
\mathcal N=\{Z\in\mathfrak k_A:L_Z\ne\mathfrak g_2\},\qquad
\mathcal G=\mathfrak k_A\setminus\mathcal N.
\]
\Cref{thm:single} supplies a generating witness.
\Cref{thm:complete-failure-set} describes $\mathcal N$ as the union of the common-vector, projective, resonant, and rank-four invariant-context families.
\Cref{thm:low-complete} gives all low-dimensional normal forms, and \Cref{cor:dimension-connectivity} determines $\dim_\RR\mathcal N=4$.
\Cref{thm:scalar-classifier,prop:irredundant-tests} give six irredundant systems of scalar equations.  Their homogeneous encoding $\mathcal P_{34}$ is an elementary consequence, not a separate classification.
Although a degree-thirty Gram determinant also decides membership for this $S$, it does not replace these geometric descriptions or their scalar elimination.
The four families need not be disjoint or irreducible components.  No geometric classification of all fibers as $T$ varies is claimed.

The local extension formulas are preparatory: they distinguish a prescribed three-plane rotation from its octonion-product-preserving extensions.
The minimum-norm choice follows from orthogonal splitting of the stabilizer; the formulas fix its exact normalization.
The two coordinate maps in \Cref{thm:charts} have explicit nonsingular differentials for a fixed finite-word alphabet.  This finite-time choice requires its own certificate and is not inferred merely from infinitesimal generation.
No physical implementation, computational speedup, globally optimal word length, or minimal number of switches is asserted.

\Cref{sec:setting,sec:models,sec:g2closure,sec:single} fix the conventions and initial generation certificates.
The uniform symmetric-pair criterion is proved first in \Cref{sec:universal-locus}, independently of the fixed-fiber classification.
\Cref{sec:failure,sec:scalar} then determine the finer geometry and scalar tests for $S$; \Cref{sec:varying-switch} proves the six-pulse local selection theorem.
\Cref{sec:charts} treats the coordinate maps.  All finite calculations used in proofs are determined by the displayed conventions and the certificates in the appendices; the exact reconstruction guide records their locations.

\section{Octonions and quaternionic contexts}
\label{sec:setting}

An orthogonal transformation of $V=\operatorname{Im}\OO$ need not preserve octonion multiplication.  We therefore fix the multiplication convention and the standard $\SO(4)$ stabilizer model before comparing the two kinds of local rotation used later.

\subsection{Conventions}

All vector spaces, Lie algebras, ranks, and dimensions are real unless a complexification is explicitly indicated.  Matrices act on column vectors, with composition from right to left.  For matrices $A,B$ we use $[A,B]=AB-BA$ and $A^{\mathsf T}$ for transpose; for octonions, $[a,b]=ab-ba$ denotes the octonion commutator.  An octonion automorphism is an invertible real-linear map preserving the unit and multiplication.  A derivation is a real-linear map $D$ satisfying $D(xy)=D(x)y+xD(y)$, hence $D(1)=0$.

For a real Euclidean space $E$, $\SO(E)$ is its determinant-one orthogonal group and $\mathfrak{so}(E)$ is the Lie algebra of skew-symmetric endomorphisms.

We use $\SU(m)$ for the determinant-one unitary group and $\mathfrak{su}(m)$ for its Lie algebra of trace-zero skew-Hermitian matrices; $A^*=\overline A^{\mathsf T}$ denotes conjugate transpose.  We identify $G_2$ and $\mathfrak g_2$ with their faithful actions on $V$; on $\OO$ their elements respectively fix and annihilate $1$.  Thus $TZT$ means composition of linear maps, not octonion multiplication.  The symbol $I$ without a subscript denotes the identity on the space then in use, except for the explicitly named scalar in \Cref{sec:scalar}; $i$ denotes the complex imaginary unit.  For an invertible linear map $g$, $\operatorname{Ad}_g Z=gZg^{-1}$.  We use $\Lie(K)$ for the Lie algebra of a Lie group $K$, and $\langle K_1,K_2\rangle$ for the subgroup of finite products of elements of the two groups and their inverses.
Let
\[
\OO=\Span_{\RR}\{1,e_1,\ldots,e_7\}
\]
with $1$ a two-sided identity, $e_i^2=-1$, and multiplication determined by the oriented Fano triples
\begin{equation}
(123),\ (145),\ (176),\ (246),\ (257),\ (347),\ (365).
\label{eq:fano}
\end{equation}
For an oriented triple $(abc)$,
\[
e_ae_b=e_c,\qquad e_be_c=e_a,\qquad e_ce_a=e_b,
\]
and reversing the order changes the sign.  Conjugation and norm are
\[
\overline{x_0+\sum_{i=1}^7x_ie_i}=x_0-\sum_{i=1}^7x_ie_i,
\qquad
N(x)=x\bar x.
\]

Here $N(x)$ is identified with its real scalar value.  The Euclidean inner product and norm are
\[
\langle x,y\rangle=\operatorname{Re}(x\bar y)
=\sum_{j=0}^{7}x_jy_j,\qquad \|x\|=\sqrt{N(x)}.
\]
The real part $\operatorname{Re}(x)$ is the coefficient of $1$, and $\operatorname{Im}(x)=x-\operatorname{Re}(x)1$.  This inner product fixes all vector-space orthogonality conventions.
The associator $[x,y,z]=(xy)z-x(yz)$ is alternating.  Put $V=\imO\simeq\RR^7$.  For $u,v\in V$,
\[
uv=-\braket{u}{v}+u\times v,
\]
which defines the cross product on $V$ and the alternating three-form $\varphi(u,v,w)=\braket{u\times v}{w}$.  A three-plane $P\subset V$ is \emph{associative} when $\RR1\oplus P$ is a quaternionic subalgebra; equivalently, it is spanned by $u,v,u\times v$ for orthonormal $u,v$.  Its orientation is the one for which $(u,v,u\times v)$ is positive.  Orthogonal complements of $P\subset V$ are taken in $V$, whereas those of $H\subset\OO$ are taken in $\OO$; in particular $P_H^\perp=H^\perp$ as four-dimensional subspaces.  With these conventions,
\[
G_2=\Aut(\OO)=\{g\in\SO(V):g^*\varphi=\varphi\},
\qquad
\mathfrak g_2=\Der(\OO).
\]

Here $(g^*\varphi)(u,v,w)=\varphi(gu,gv,gw)$; preservation of $\varphi$ also makes the cross product $G_2$-invariant.  Ranks, kernels, and characteristic polynomials of derivations below refer to their seven-dimensional action on $V$, unless stated otherwise.
Standard references for these octonionic and $G_2$ conventions are \cite{Baez2002,SpringerVeldkamp2000,ChemtovKarigiannis2022}.

For $a\in\OO$, let $L_a(x)=ax$ and $R_a(x)=xa$.  The standard inner derivations are
\[
D_{a,b}=[L_a,L_b]+[L_a,R_b]+[R_a,R_b],
\]
and satisfy
\[
D_{a,b}(x)=[[a,b],x]-3[a,b,x].
\]
We abbreviate $D_{ij}:=D_{e_i,e_j}$.  These maps obey the Leibniz rule, and the $D_{e_i,e_j}$ span the fourteen-dimensional algebra $\mathfrak g_2$ \cite{SpringerVeldkamp2000,RauschSlupinski2022}.  We use the ordered basis
\begin{equation}
\mathcal B_{G_2}=(D_{12},D_{13},D_{14},D_{15},D_{16},D_{17},D_{23},D_{24},D_{25},D_{26},D_{27},D_{45},D_{46},D_{47}).
\label{eq:g2basis}
\end{equation}

For later parametrizations, $\RR P(E)$ denotes the set of one-dimensional real subspaces of $E$, and $\RR P^m=\RR P(\RR^{m+1})$.  Thus $[v]$ is the line $\RR v$ for $v\ne0$, and homogeneous coordinates such as $[c:s]$ are defined up to common nonzero real scaling.  The notation $\CC P(E)$ and $\CC P^m$ has the analogous meaning for complex lines.

\subsection{Quaternionic subalgebras and their stabilizers}
For a context $H$, the plane $P_H=H\cap V$ is associative and $H=\RR1\oplus P_H$.  Conversely, every oriented associative three-plane $P\subset V$ determines the quaternionic subalgebra $\RR1\oplus P$.  The seven coordinate contexts are
\[
H_{abc}=\Span\{1,e_a,e_b,e_c\}
\]
for the triples in \eqref{eq:fano}.  We write
\[
H_A=H_{123},\qquad H_B=H_{145},\qquad P_A=P_{H_A},\quad P_B=P_{H_B}.
\]
Figure~\ref{fig:fano} records the seven coordinate contexts in the Fano-plane convention used throughout.

\begin{figure}[htbp]
\centering
\begin{tikzpicture}[scale=1.25,
  pt/.style={circle,fill=black,inner sep=1.6pt},
  lbl/.style={font=\small},
  ln/.style={line width=0.7pt,gray!70},
  hl/.style={line width=1.6pt,black}]
% vertices of the triangle: 1, 2, 4 ; midpoints: 3 = 1+2, 5 = 1+4, 6 = 2+4 ;
% centre: 7 = 1+2+4
\coordinate (p1) at (90:2);
\coordinate (p2) at (210:2);
\coordinate (p4) at (330:2);
\coordinate (p3) at ($(p1)!0.5!(p2)$);
\coordinate (p5) at ($(p1)!0.5!(p4)$);
\coordinate (p6) at ($(p2)!0.5!(p4)$);
\coordinate (p7) at (0,0);
% the six straight lines
\draw[hl]  (p1) -- (p2);              % A = {1,2,3}
\draw[hl]  (p1) -- (p4);              % B = {1,4,5}
\draw[ln]  (p2) -- (p4);              % D = {2,4,6}
\draw[hl]  (p1) -- (p6);              % C = {1,7,6}
\draw[ln]  (p2) -- (p5);              % E = {2,5,7}
\draw[ln]  (p4) -- (p3);              % F = {3,4,7}
% the circle
\draw[ln]  (p7) circle (1);           % G = {3,6,5}
% points
\foreach \i in {1,2,3,4,5,6,7}{\node[pt] at (p\i) {};}
\node[lbl,above]       at (p1) {$e_1$};
\node[lbl,left]        at (p2) {$e_2$};
\node[lbl,left]        at (p3) {$e_3$};
\node[lbl,right]       at (p4) {$e_4$};
\node[lbl,right]       at (p5) {$e_5$};
\node[lbl,below right] at (p6) {$e_6$};
\node[lbl,above right] at (p7) {$e_7$};
% line names
\node[lbl] at ($(p1)!0.5!(p2)+(-0.55,0.30)$) {$A$};
\node[lbl] at ($(p1)!0.5!(p4)+( 0.55,0.30)$) {$B$};
\node[lbl] at ($(p1)!0.5!(p6)+( 0.30,0.45)$) {$C$};
\node[lbl] at ($(p2)!0.22!(p4)+( 0.00,-0.32)$) {$D$};
\node[lbl] at ($(p2)!0.32!(p5)+(0.30,-0.16)$) {$E$};
\node[lbl] at ($(p4)!0.32!(p3)+(-0.30,-0.16)$) {$F$};
\node[lbl] at (203:1.34) {$G$};
\end{tikzpicture}
\caption{The seven coordinate contexts as the lines of the Fano plane,
$A=H_{123}$, $B=H_{145}$, $C=H_{176}$, $D=H_{246}$, $E=H_{257}$,
$F=H_{347}$, $G=H_{365}$.  The three bold lines are the three coordinate contexts containing $e_1$, called the pencil through $e_1$.}
\label{fig:fano}
\end{figure}

The transitive action of $G_2$ on associative three-planes and the identification of the stabilizer with $\SO(4)$ are classical \cite{HarveyLawson1982,Cacciatori2005,KnarrStroppel2025,ChemtovKarigiannis2022}.  Thus the space of quaternionic contexts is the homogeneous space $G_2/\SO(4)$.
We record the standard coordinate form of the stabilizer because it is used repeatedly below.

For any context $H$, set
\[
K_H=\Stab_{G_2}(H)=\{g\in G_2:g(H)=H\},\qquad
\mathfrak k_H=\Lie(K_H)=\{Z\in\mathfrak g_2:Z(H)\subseteq H\}.
\]
This is a setwise stabilizer, not a pointwise stabilizer.

For a coordinate context we also write $\mathfrak k_{abc}=\mathfrak k_{H_{abc}}$.  We abbreviate $K_{H_A},\mathfrak k_{H_A}$ by $K_A,\mathfrak k_A$, and use the analogous convention for $H_B$ and for the pointwise stabilizers below.  Fix $H_0=H_A=H_{123}$ and identify it with $\HH$, so that $\OO=\HH\oplus\HH e_4$.  Write $\Sp(1)=\{a\in\HH:N(a)=1\}$ for the group of unit quaternions; its Lie algebra $\mathfrak{sp}(1)=\operatorname{Im}\HH$ uses the quaternionic commutator.  For $a,b\in\Sp(1)$ define
\begin{equation}
\gamma_{a,b}(x+ye_4)=axa^{-1}+(bya^{-1})e_4.
\label{eq:gamma}
\end{equation}

For an arbitrary context $H$, write
\[
K_H^{\mathrm{pt}}=\{g\in K_H:g|_H=I\},\qquad
\mathfrak k_H^{\mathrm{pt}}=\Lie(K_H^{\mathrm{pt}})
=\{Z\in\mathfrak k_H:Z|_{P_H}=0\}.
\]

\begin{lemma}[Standard stabilizer model]
\label{prop:stabilizer}
The maps \eqref{eq:gamma} are octonion automorphisms preserving $H_0$.  The homomorphism $(a,b)\mapsto\gamma_{a,b}$ has kernel $\{(1,1),(-1,-1)\}$, and
\[
K_{H_0}=\Stab_{G_2}(H_0)
\cong(\Sp(1)\times\Sp(1))/\{\pm(1,1)\}
\cong\SO(4).
\]
The pointwise stabilizer of $H_0$ is $K_{H_0}^{\mathrm{pt}}\cong\Sp(1)$.  Formula \eqref{eq:gamma} is the usual $\SO(4)\subset G_2$ model; see, for example, \cite{Cacciatori2005}.
\end{lemma}

\begin{proof}
Using the Cayley--Dickson product
\[
(x+ye_4)(u+ve_4)=(xu-\bar v y)+(vx+y\bar u)e_4,
\]
associativity in $\HH$ gives $\gamma_{a,b}(zw)=\gamma_{a,b}(z)\gamma_{a,b}(w)$.  Its restriction to $H_0$ is $x\mapsto axa^{-1}$.  The only pairs acting trivially on both summands are $(1,1)$ and $(-1,-1)$.  Conversely, if $g\in G_2$ preserves $H_0$, then $g|_{H_0}$ is an inner quaternion automorphism $x\mapsto axa^{-1}$.  Orthogonality gives $g(e_4)=ue_4$ with $u\in\Sp(1)$, and multiplicativity forces
\[
g(ye_4)=(ua\,y\,a^{-1})e_4.
\]
Taking $b=ua$ gives \eqref{eq:gamma}.  Setting $a=1$, modulo the diagonal kernel, gives the pointwise stabilizer.
\end{proof}

For $H_0$, differentiating \eqref{eq:gamma} at $(1,1)$ gives convenient coordinates on $\mathfrak k_{H_0}$.  If $q_a,q_b\in\operatorname{Im}\HH$, the corresponding derivation acts by
\[
Z(q_a,q_b)(x+ye_4)=[q_a,x]+(q_by-yq_a)e_4.
\]
Thus every element of $\mathfrak k_{H_0}$ is represented uniquely by a pair $(q_a,q_b)\in\operatorname{Im}\HH\oplus\operatorname{Im}\HH$.  The same notation is transported to any context by conjugation.  The pair $(q_a,q_b)$ consists of imaginary quaternions, not group elements.  When it is written as two triples, their coordinates are taken in the ordered basis $(e_1,e_2,e_3)$ of $\operatorname{Im}H_A$.

\subsection{Reflections, parity, and the local domain}\label{subsec:parity-domain}
Recall the reflection $\sigma_H(h+y)=h-y$ from the introduction.  In the adapted model \eqref{eq:gamma}, it is $\gamma_{1,-1}$, so \Cref{prop:stabilizer} proves multiplicativity.  On $V$ define
\[
\mathscr C=\{\sigma_H:H\text{ a quaternionic context}\}
 =\{T\in G_2:T^2=I,\ \operatorname{tr}_V T=-1\}.
\]
The equality follows because the three-dimensional positive eigenspace of such a $T$ is closed under the cross product and hence associative.  Transitivity identifies $\mathscr C$ with the connected compact real algebraic manifold $G_2/\SO(4)$.

The trace condition excludes the identity and fixes the positive and negative eigenspace dimensions as three and four on $V$.  In statements about almost every reflection, invariant measure is the normalized $G_2$-invariant measure on this compact homogeneous space.  


For $T\in\mathscr C$, set
\[
\mathfrak k_T=\{A\in\mathfrak g_2:TAT=A\},\qquad
\mathfrak m_T=\{B\in\mathfrak g_2:TBT=-B\}.
\]
The fixed algebra is the stabilizer of the reflecting context, of dimension six; its complementary odd eigenspace has dimension eight.
For every $Z$, the parity coordinates $A_+=(Z+TZT)/2$, $A_-=(Z-TZT)/2$ give a linear bijection
$\mathfrak g_2\to\mathfrak k_T\oplus\mathfrak m_T$ with inverse addition.
The two ordered pairs $(Z,TZT)$ and $(A_+,A_-)$ have the same span, so they generate the same algebra.
That algebra is $\Ad_T$-stable, since a homogeneous bracket has its parity sign.
No local coupling is imposed when $Z$ is unrestricted.
For $Z\in\mathfrak k_A$, the allowed pairs are instead the six-dimensional subspace $\mathcal W_A(T)$ defined in the introduction.

We write $\alpha=\|q_a\|$ and $\beta=\|q_b\|$ for the local quaternionic velocity norms.
The action $[q_a,x]$ on $\operatorname{Im}\HH$ has frequency $2\alpha$; on $\HH e_4$ the commuting left and right multiplications give frequencies $|\beta-\alpha|$ and $\beta+\alpha$.
These conventions allow zero frequencies.  A \emph{principal spectrum} means the nonzero positive frequencies are $\lambda,2\lambda,3\lambda$ for some $\lambda>0$.

\subsection{Compact maximal subalgebras}\label{subsec:compact-maximals}
For $0\ne v\in V$, write $\mathfrak g_2^v=\{Z\in\mathfrak g_2:Zv=0\}$.  We use the compact maximal-subalgebra classification in its real form; see Draper and Mart\'in-Gonz\'alez \cite{DraperMartin2025}: every proper subalgebra of the compact $\mathfrak g_2$ is contained in a vector stabilizer $\mathfrak{su}(3)$, a quaternionic stabilizer $\mathfrak{so}(4)$, or a principal $\mathfrak{su}(2)$.  A subalgebra of a compact Lie algebra is reductive, because the ambient positive invariant inner product restricts to it.  The maximal compact embeddings correspond to the maximal reductive complex types $A_2$, $A_1+A_1$, and the principal $A_1$ of Dynkin index $28$ \cite{Dynkin1952b}.  This is not a claim about all maximal complex subalgebras, which also include nonreductive parabolic subalgebras.  Compact realizations of each type are conjugate by $G_2$: in a faithful unitary realization, if $g$ conjugates a compact subgroup $K_1$ into the compact group, then $h^*g^*gh=g^*g$ for all $h\in K_1$.  Hence the positive factor $(g^*g)^{1/2}$ centralizes $K_1$, leaving the compact factor in the polar decomposition of $g$ to perform the conjugation.  This explains the passage from the complex list to the compact statement being used.

The symbols $A_2$, $A_1+A_1$, and $A_1$ denote complex root-system types; the corresponding compact algebras here are $\mathfrak{su}(3)$, $\mathfrak{su}(2)\oplus\mathfrak{su}(2)$, and $\mathfrak{su}(2)$.  For the cited Dynkin index, invariant forms are normalized so that long roots have squared length two, and the index is the factor by which the ambient normalized form restricts to the embedded simple algebra.  The numerical index is used only to identify the embedding in the classical list.

Here \emph{principal} denotes the class of $\mathfrak{su}(2)$ subalgebras acting irreducibly on $V$; a nonzero element has three positive frequencies in the ratio $1:2:3$.  The first two containing algebras act reducibly on $V$, with invariant subspaces of dimensions one and three, respectively.  These three containing types are the structural input to the uniform criterion and the later fixed-fiber reduction.


\section{Local extension conventions}
\label{sec:models}

This preliminary section fixes what a local generator means.  A rotation prescribed only on $P_H$ does not determine a transformation of all seven imaginary octonion directions.  We compare two natural extensions.  The first fixes $P_H^\perp$ pointwise and is merely orthogonal; the second preserves octonion multiplication and therefore lies in $G_2$.

A \emph{context switch} from $H$ to $H'$ is an orthogonal map $S\in\SO(V)$ satisfying $S(P_H)=P_{H'}$; it is \emph{product preserving} when $S\in G_2$.  For a set $\mathcal A$ of matrices, $\Lie(\mathcal A)$ denotes the smallest real Lie subalgebra containing $\mathcal A$.

For a fixed context $H$, the frozen-complement group and its Lie algebra are
\[
B_H=\{R\oplus I_{P_H^\perp}:R\in\SO(P_H)\}\subset\SO(7),
\qquad
\mathfrak b_H=\mathfrak{so}(P_H)\oplus0.
\]
The product-preserving local group is instead
\[
K_H=\Stab_{G_2}(H)\cong\SO(4),
\qquad
\mathfrak k_H=\operatorname{Lie}(K_H)\subset\mathfrak g_2.
\]
Both groups restrict onto $\SO(P_H)$, but the restriction is one-to-one for $B_H$ and has kernel $K_H^{\mathrm{pt}}$ for $K_H$.  Thus prescribing the same rotation on $P_H$ gives a unique frozen extension and a family of product-preserving extensions with different actions on $P_H^\perp$.

\begin{proposition}[Frozen-complement obstruction]
\label{prop:frozen}
For every context $H$,
\[
B_H\cap G_2=\{I\},\qquad
\mathfrak b_H\cap\mathfrak g_2=\{0\}.
\]
Thus a nontrivial rotation of $P_H$ cannot preserve the octonion product while fixing $P_H^\perp$ pointwise.
\end{proposition}

\begin{proof}
Conjugate to $H_0$.  If $g=\gamma_{a,b}\in K_{H_0}$ fixes $H_0^\perp$ pointwise, then $bya^{-1}=y$ for every $y\in\HH$.  Taking $y=1$ gives $b=a$, and then $aya^{-1}=y$ for every quaternion $y$, so $a=\pm1$.  The diagonal kernel identifies both choices with the identity.  Differentiation gives the Lie-algebra statement, and conjugation by $G_2$ handles an arbitrary context.
\end{proof}

For $w\in V$ let $C_w\in\mathfrak{so}(V)$ be the cross-product operator
$C_w(x)=w\times x$.  The map $w\mapsto C_w$ is injective and its image is a
$\Gtwo$-submodule of $\mathfrak{so}(V)$ isomorphic to $V$; since $V$ and
$\gtwo$ are inequivalent $\gtwo$-modules, the classical splitting
\begin{equation}
\mathfrak{so}(V)=\gtwo\oplus\{C_w:w\in V\}
\label{eq:so7split}
\end{equation}
is orthogonal for the trace form $\langle A,B\rangle=-\operatorname{tr}(AB)$ \cite{ChemtovKarigiannis2022}.

The projection formulas in \cite[(2.23) and (4.5)]{ChemtovKarigiannis2022} give the underlying decomposition.  The next proposition records their normalization in terms of $D_{u,v}$ and includes a direct proof.

Fix an orthonormal oriented associative triple $(u,v,w)$, so $w=u\times v$.
The frozen generator $J_{u,v}\in\mathfrak{so}(V)$ sends $u$ to $v$, $v$ to $-u$,
and annihilates $\{u,v\}^\perp$.  We use
$\|B\|_{\mathrm{tr},E}^2=-\operatorname{tr}_E(B^2)$ for a skew operator on a Euclidean space $E$;
the subscript $E$ is omitted when $E=V$.

We abbreviate $J_{ij}=J_{e_i,e_j}$.  For a list of skew operators $A_1,\ldots,A_m$ on $V$, its trace-Gram matrix is
\[
\operatorname{Gram}_{\mathrm{tr}}(A_1,\ldots,A_m)
=\bigl(-\operatorname{tr}(A_iA_j)\bigr)_{i,j=1}^{m}.
\]
More generally, a Gram matrix is formed from the stated inner product.  For the positive definite real inner products used in the Gram certificates, its determinant is positive exactly when the listed vectors are independent.

\begin{proposition}[Derivation-coordinate form of the standard splitting]
\label{prop:split}
With the preceding notation,
\begin{equation}
J_{u,v}=C_w-\tfrac12\,[u,v,\,\cdot\,]
=\tfrac16\,D_{u,v}+\tfrac13\,C_w ,
\label{eq:split}
\end{equation}
and \eqref{eq:split} is the decomposition of $J_{u,v}$ along
\eqref{eq:so7split}.  Consequently
\[
\frac{\|\tfrac16 D_{u,v}\|_{\mathrm{tr}}^2}{\|J_{u,v}\|_{\mathrm{tr}}^2}=\frac23,
\qquad
\frac{\|\tfrac13 C_w\|_{\mathrm{tr}}^2}{\|J_{u,v}\|_{\mathrm{tr}}^2}=\frac13 ,
\]
so the two orthogonal components carry respectively two thirds and one third of the squared trace norm.  Equivariance and transitivity on orthonormal two-frames force these ratios to be constant; the displayed calculation determines their values.
\end{proposition}

\begin{proof}
Both sides of the first identity in \eqref{eq:split} are linear, so it is
enough to test them on $u,v,w$ and on $\{u,v,w\}^\perp$.  On $u$: $C_w(u)=w\times u=v$
and $[u,v,u]=0$ because the associator is alternating; symmetrically on $v$.
On $w$: $C_w(w)=0$ and $[u,v,w]=(uv)w-u(vw)=w\cdot w-u\cdot u=-1+1=0$, since
$uv=w$ and $vw=u$ inside the associative context.  For $x\perp\RR1\oplus\Span\{u,v,w\}$
one has $u(vx)=-(uv)x$, hence
$[u,v,x]=(uv)x-u(vx)=2wx=2\,C_w(x)$, while $J_{u,v}(x)=0$; this gives the
first identity.  The second follows from
$D_{u,v}(x)=[[u,v],x]-3[u,v,x]$ together with $[u,v]=2w$ and
$[w,x]=2\,w\times x$ for $x\in V$, so that $D_{u,v}=4C_w-3[u,v,\cdot\,]$.
Eliminating $[u,v,\cdot\,]$ between the two expressions yields
$J_{u,v}=\tfrac16 D_{u,v}+\tfrac13 C_w$.  For $\|A\|_{\mathrm{tr}}^2=-\operatorname{tr}(A^2)$, the norms are computed from
$\|J_{u,v}\|_{\mathrm{tr}}^2=2$, $\|D_{u,v}\|_{\mathrm{tr}}^2=48$ and $\|C_w\|_{\mathrm{tr}}^2=6$, and the two summands
are orthogonal by \eqref{eq:so7split}.
\end{proof}

\begin{remark}
\Cref{prop:split} strengthens \Cref{prop:frozen} by identifying both components of every frozen generator.  The term $D_{u,v}/6$ is the orthogonal projection of $J_{u,v}$ onto $\mathfrak g_2$, whereas $C_w/3$ is its component perpendicular to $\mathfrak g_2$.  Since $w=u\times v$, this decomposition is intrinsic to the associative plane and does not depend on the chosen coordinates.
\end{remark}

Prescribing a local generator fixes its action on $P_H$ but does not fix the complementary action.  The orthogonal complement of the restriction kernel selects the minimum-norm extension, by the standard least-norm property of an orthogonal splitting.  The next calculation gives its normalization in the stabilizer model.

For $J\in\mathfrak{so}(P_H)$ there is a unique $w\in P_H$ with $Jx=w\times x$.
In any quaternionic coordinates adapted to $H$, define the candidate extension
\[
\mathcal L_H(J)=Z(w/2,0).
\]
The proposition proves that this expression is independent of the adapted coordinates by
characterizing it as the unique minimum of the invariant trace norm.

\begin{proposition}[Minimum-norm extension of an arbitrary local generator]
\label{cor:forced-motion}
For every quaternionic context $H$ and every $J\in\mathfrak{so}(P_H)$, the extensions of $J$ in $\mathfrak g_2$ form
\[
\{L\in\mathfrak g_2:L|_{P_H}=J\}
 =\mathcal L_H(J)+\mathfrak k_H^{\mathrm{pt}}.
\]
The distinguished extension $\mathcal L_H(J)$ is the unique extension of minimum trace norm.  More precisely, for every extension $L$,
\begin{align}
 \|L\|_{\mathrm{tr}}^2
 &=\frac32\|J\|_{\mathrm{tr},P_H}^2+
       \|L-\mathcal L_H(J)\|_{\mathrm{tr}}^2,\label{eq:minimum-lift-norm}\\
 \|L|_{H^\perp}\|_{\mathrm{tr},H^\perp}^2
 &\ge\frac12\|J\|_{\mathrm{tr},P_H}^2.\nonumber
\end{align}
Equality in the inequality holds precisely for $L=\mathcal L_H(J)$.  The map $\mathcal L_H$ is linear, preserves Lie brackets, and is equivariant under simultaneous $G_2$-conjugation of $H$ and $J$.

For an orthonormal associative triple $(u,v,w=u\times v)$ spanning $P_H$ and $J=J_{u,v}|_{P_H}$, one has
\[
\mathcal L_H(J)=\frac14D_{u,v},\qquad
\mathcal L_H(J)|_{H^\perp}=-\frac12C_w.
\]
Thus the half-speed complementary motion belongs to the minimum-norm extension, not to every extension with the same local angular velocity.
\end{proposition}

\begin{proof}
Conjugate to the model \eqref{eq:gamma}.  On $P_H$, the derivation $Z(q_a,q_b)$ acts as $2q_a\times x$, so restriction to $J$ fixes $q_a=w/2$ and leaves $q_b$ arbitrary.  On $\HH e_4$ its action is $L_{q_b}-R_{q_a}$.  The traces of $L_{q_b}^2$ and $R_{q_a}^2$ on $\HH$ are $-4\|q_b\|^2$ and $-4\|q_a\|^2$, while $\operatorname{tr}(L_{q_b}R_{q_a})=0$ for imaginary $q_a,q_b$.  The three-dimensional block contributes $8\|q_a\|^2$.  Hence
\[
\|Z(q_a,q_b)\|_{\mathrm{tr}}^2=12\|q_a\|^2+4\|q_b\|^2.
\]
Since $\|J\|_{\mathrm{tr},P_H}^2=2\|w\|^2$, the asserted identity and inequality follow, with equality exactly at $q_b=0$.  The first $\mathfrak{sp}(1)$ factor is a Lie subalgebra, and its restriction to $\mathfrak{so}(P_H)$ is an isomorphism; its inverse is $\mathcal L_H$.  Uniqueness of the minimum and invariance of the trace form prove equivariance, independently of the chosen adapted coordinates.

On $P_H$ the associator vanishes, so $D_{u,v}=4J_{u,v}$.  On $H^\perp$, \Cref{prop:split} gives $D_{u,v}=-2C_w$.  These restrictions identify $D_{u,v}/4$ with $Z(w/2,0)$ and prove the final assertions.
\end{proof}

\begin{example}[Equal local velocities with different complementary actions]
For $H_{123}$, both
\[
L_0=\frac14D_{12},\qquad L_1=\frac13D_{12}+\frac16D_{47}
\]
restrict to $J_{12}$ on $P_A$.  The exponential $e^{tL_0}$ rotates the $(e_4,e_7)$ and $(e_5,e_6)$ planes at half the prescribed angular speed, with opposite orientations.  The second instead satisfies
\[
L_1e_4=L_1e_7=0,\qquad L_1e_5=e_6,\qquad L_1e_6=-e_5.
\]
Both are derivations, so their exponentials preserve multiplication; only $L_0$ minimizes the trace norm.  Neither derivation annihilates the whole complement, consistently with \Cref{prop:frozen}.
\end{example}

An unsigned coordinate switch provides another illustration.  Let it fix $e_1,e_6,e_7$ and exchange $e_2\leftrightarrow e_4$, $e_3\leftrightarrow e_5$.  This map is orthogonal and involutive, but
\[
e_2e_6=-e_4,
\]
so the images of the two sides differ by a sign.  The corrected switch is
\begin{equation}
\begin{aligned}
S(e_1)&=e_1,&S(e_2)&=e_4,&S(e_3)&=e_5,\\
S(e_4)&=e_2,&S(e_5)&=e_3,&S(e_6)&=-e_6,&S(e_7)&=-e_7.
\end{aligned}
\label{eq:switch}
\end{equation}
Checking the seven oriented Fano products shows that $S(xy)=S(x)S(y)$ for basis elements and hence for all $x,y\in\OO$.  Therefore $S\in G_2$ and $S(H_A)=H_B$.

\section{Product-preserving closure from quaternionic stabilizers}
\label{sec:g2closure}

In the product-preserving model each local algebra is six-dimensional, and the closure depends on the relative position of two copies of $\mathfrak{so}(4)$ inside $\mathfrak g_2$.  The isotropy representation of $G_2/\SO(4)$ is irreducible; hence a Lie algebra that properly contains one context stabilizer must be all of $\mathfrak g_2$.

Put $\mathfrak k=\mathfrak k_{H_0}$ and let $\mathfrak m=\mathfrak k^\perp\cap\mathfrak g_2$ for the trace form.  Write $V_m$ for the $(m+1)$-dimensional irreducible complex representation of $\Sp(1)\simeq\SU(2)$, and use $\boxtimes$ for the external tensor product of the two $\Sp(1)$ factors in the stabilizer.

For a real vector space or Lie algebra $E$, write $E_{\CC}=E\otimes_{\RR}\CC$.

Inner products extend complex-bilinearly in the orthogonal-algebra decompositions below.  The external tensor product means $(g_1,g_2)(v\otimes w)=g_1v\otimes g_2w$; $\Lambda^2E$ is the exterior square.  The isotropy representation is the action of $K_{H_0}$ on the tangent space of $G_2/K_{H_0}$ at the identity coset, identified here with $\mathfrak m$.

\begin{lemma}[Isotropy branching and maximality]
\label{lem:maximality}
For $K=K_{H_0}\cong(\Sp(1)_a\times\Sp(1)_b)/\{\pm(1,1)\}$,
\begin{align}
V_{\CC}&\cong(V_2\boxtimes V_0)\oplus(V_1\boxtimes V_1),\label{eq:Vbranch}\\
\mathfrak g_{2,\CC}&\cong(V_2\boxtimes V_0)\oplus(V_0\boxtimes V_2)\oplus(V_3\boxtimes V_1).
\label{eq:gbranch}
\end{align}
Thus
\[
\mathfrak k_{\CC}\cong(V_2\boxtimes V_0)\oplus(V_0\boxtimes V_2),
\qquad
\mathfrak m_{\CC}\cong V_3\boxtimes V_1
\]
and $\mathfrak m$ is irreducible.  Therefore $\mathfrak k\cong\mathfrak{so}(4)$ is a maximal proper subalgebra of $\mathfrak g_2$, in agreement with the classical maximal-subalgebra classification \cite{Dynkin1952b}.
\end{lemma}

\begin{proof}
The action \eqref{eq:gamma} splits
\[
V=\operatorname{Im}H_0\oplus H_0e_4.
\]
The first summand is the adjoint three-dimensional module of $\Sp(1)_a$ and is trivial under $\Sp(1)_b$, giving $V_2\boxtimes V_0$.  On the second, $(a,b)$ acts on the quaternion coefficient by $y\mapsto bya^{-1}$, whose complexification is $V_1\boxtimes V_1$.  This proves \eqref{eq:Vbranch}.

Use $\mathfrak{so}(V_{\CC})\cong\Lambda^2V_{\CC}$ and
\[
\Lambda^2(A\oplus B)=\Lambda^2A\oplus(A\otimes B)\oplus\Lambda^2B.
\]
The tensor-product decomposition gives
\[
\Lambda^2(V_2)=V_2,
\quad
V_2\otimes V_1=V_3\oplus V_1,
\quad
\Lambda^2(V_1\boxtimes V_1)=(V_2\boxtimes V_0)\oplus(V_0\boxtimes V_2).
\]
Consequently
\[
\mathfrak{so}(V_{\CC})\cong
2(V_2\boxtimes V_0)\oplus(V_0\boxtimes V_2)
\oplus(V_3\boxtimes V_1)\oplus(V_1\boxtimes V_1).
\]
Removing the direct summand $V_{\CC}$ in \eqref{eq:so7split} yields \eqref{eq:gbranch}.  If $\mathfrak k\subseteq\mathfrak h\subseteq\mathfrak g_2$, the orthogonal splitting $\mathfrak g_2=\mathfrak k\oplus\mathfrak m$ gives $\mathfrak h=\mathfrak k\oplus W$ with $W$ a $\mathfrak k$-submodule of the irreducible module $\mathfrak m$.  Hence $W=0$ or $W=\mathfrak m$, proving maximality.
\end{proof}

\begin{lemma}[The stabilizer determines the context]
\label{lem:context-determined}
If $\mathfrak k_H=\mathfrak k_{H'}$, then $H=H'$.
\end{lemma}

\begin{proof}
Conjugate so that $H=H_0$.  Under $K_{H_0}$, the decomposition
\[
V=\operatorname{Im}H_0\oplus H_0e_4
\]
consists of irreducible inequivalent real modules of dimensions three and four.  Therefore $\operatorname{Im}H_0$ is the unique invariant three-dimensional subspace.  If $\mathfrak k_{H'}=\mathfrak k_{H_0}$, the associative plane $P_{H'}$ is invariant and must equal $P_{H_0}$; hence $H'=H_0$.
\end{proof}

\begin{proposition}[Two-context product-preserving closure]
\label{thm:two-context}
For any two distinct contexts $H\ne H'$,
\[
\Lie(\mathfrak k_H\cup\mathfrak k_{H'})=\mathfrak g_2,
\qquad
\langle K_H,K_{H'}\rangle=G_2.
\]
\end{proposition}

\begin{proof}
By \Cref{lem:context-determined}, $\mathfrak k_{H'}\ne\mathfrak k_H$.  The generated algebra strictly contains $\mathfrak k_H$, so maximality from \Cref{lem:maximality} forces it to be $\mathfrak g_2$.  The passage to the group is the usual Lie-rank argument \cite{JurdjevicSussmann1972}.  To make openness explicit, let $\Gamma=\langle K_H,K_{H'}\rangle$ and let $E$ be the span of all $\Ad_g U$ with $g\in\Gamma$ and $U\in\mathfrak k_H\cup\mathfrak k_{H'}$.  This space is $\Ad_\Gamma$-invariant.  Differentiating conjugation by the one-parameter groups shows that it contains the Lie algebra generated by the two local algebras, hence $E=\mathfrak g_2$.  Select fourteen such conjugates forming a basis.  The product of their one-parameter groups lies in $\Gamma$ and has invertible differential at zero.  The inverse-function theorem therefore gives an identity neighborhood in $\Gamma$.  Thus $\Gamma$ is open and equals the connected group $G_2$.
\end{proof}

Conjugation by a product-preserving switch carries the stabilizer of one context to the stabilizer of its image, giving the following immediate consequence.

\begin{corollary}
\label{cor:stabilizer-switch}
Let $H$ be a quaternionic context and let $S\in G_2$ satisfy $S(H)\neq H$.  Then
\[
\Lie\!\left(\mathfrak k_H\cup\Ad_S\mathfrak k_H\right)=\mathfrak g_2,
\qquad
\langle K_H,S\rangle=G_2.
\]
Thus one complete local context-control group plus one product-preserving context switch is sufficient for global $G_2$ control.
\end{corollary}

\begin{proof}
Conjugation gives $SK_HS^{-1}=K_{S(H)}$ and
$\Ad_S\mathfrak k_H=\mathfrak k_{S(H)}$.  Since $S(H)\neq H$, \Cref{thm:two-context} applies.  The group generated by $K_H$ and $S$ contains both $K_H$ and $K_{S(H)}$, hence all of $G_2$.
\end{proof}

Thus \Cref{cor:stabilizer-switch} reduces the continuously available directions to one six-dimensional local stabilizer and one fixed switch.  The next section reduces the continuous control further to one direction.

For the motivating pair the enlargement beyond a single stabilizer can also be seen directly.  An explicit basis is
\begin{align}
\mathfrak k_A&=\Span\{D_{12},D_{13},D_{23},D_{45},D_{46},D_{47}\},\\
\mathfrak k_B&=\Span\{D_{14},D_{15},D_{23},D_{26},D_{27},D_{45}\}.
\end{align}
Their intersection is $\Span\{D_{23},D_{45}\}$ and their sum has dimension ten.  The four identities
\begin{align}
[D_{12},D_{14}]&=-2D_{17}+4D_{24},\\
[D_{12},D_{15}]&= 2D_{16}+4D_{25},\\
[D_{12},D_{26}]&=-4D_{16}-2D_{25},\\
[D_{12},D_{27}]&=-4D_{17}+2D_{24}
\end{align}
generate the four missing directions.  These commutators exhibit directly the four directions missing from $\mathfrak k_A+\mathfrak k_B$ and provide a coordinate check of \Cref{thm:two-context} for the fixed pair.

\section{One pulse and one fixed switch}
\label{sec:single}

The preceding result allows all six directions in a local stabilizer.  We now impose the stronger restriction that only one continuously tunable direction is available and that the second generator is its conjugate by the fixed switch.

Fix the signed switch $S$ in \eqref{eq:switch} and define
\begin{equation}
X=D_{46}+D_{47}\in\mathfrak k_A,
\qquad
Y=\Ad_SX=-(D_{26}+D_{27})\in\mathfrak k_B.
\label{eq:XY}
\end{equation}

For this pair let $V_0=\Span\{X,Y\}$ and $V_{r+1}=V_r+[V_r,V_r]$.
We also fix bracket words by the following recursion for $(X,Y)$.
For another ordered pair $(U,V)$, $W_i(U,V)$ denotes the same word with
$X,Y$ replaced by $U,V$; $W_i$ without arguments means $W_i(X,Y)$.
\begin{align*}
W_1&=X,&W_2&=Y,&W_3&=[W_1,W_2],\\
W_4&=[W_1,W_3],&W_5&=[W_2,W_3],&W_6&=[W_1,W_4],\\
W_7&=[W_1,W_5],&W_8&=[W_2,W_5],&W_9&=[W_3,W_4],\\
W_{10}&=[W_3,W_5],&W_{11}&=[W_4,W_5],&W_{12}&=[W_1,W_9],\\
W_{13}&=[W_3,W_9],&W_{14}&=[W_3,W_{10}].
\end{align*}

\begin{theorem}[Constrained single-switch generation]
\label{thm:single}
The pair in \eqref{eq:XY} satisfies
\[
\Lie\{X,Y\}=\mathfrak g_2.
\]
The successive dimensions of the closure filtration are
\[
2\longrightarrow3\longrightarrow5\longrightarrow11\longrightarrow14\longrightarrow14.
\]
Consequently finite products of the pulse $e^{tX}$, the switch $S$, and their inverses reach every element of the connected group $G_2$.
\end{theorem}

\begin{proof}
All fourteen words belong to $\mathfrak g_2$ because $X,Y\in\mathfrak g_2$.
Their coordinate columns in the ordered basis \eqref{eq:g2basis} are printed in
Table~\ref{tab:brackets}, in Appendix~\ref{app:single-brackets}.  The determinant of that matrix is
\begin{equation}
6091593448313534987709865525248=2^{88}3^9\ne0.
\label{eq:bracketdet}
\end{equation}
Thus the fourteen words form a basis of $\mathfrak g_2$.

The intermediate dimensions follow from the same certificate.  By the recursion and the Jacobi identity
$[W_2,W_4]=[W_1,W_5]=W_7$,
\[
V_1=\Span\{W_1,W_2,W_3\},\qquad
V_2=\Span\{W_1,\ldots,W_5\},\qquad
V_3=\Span\{W_1,\ldots,W_{11}\}.
\]
Every bracket between the five generators of $V_2$ is in the displayed span for $V_3$,
and every listed $W_i$ with $i\le11$ occurs in that step.  The remaining words
$W_{12},W_{13},W_{14}$ belong to $V_4$, so $V_4=\mathfrak g_2$ and $V_5=V_4$.
Independence of the prefixes of the fourteen-word basis gives the stated dimensions.

Finally, $Se^{tX}S^{-1}=e^{tY}$, so the pulse--switch group contains both
one-parameter subgroups.  Their Lie algebra is $\mathfrak g_2$.  Applying the
product-map argument in the proof of \Cref{thm:two-context} gives an open identity
neighborhood in that group: choose fourteen conjugated infinitesimal directions
spanning $\mathfrak g_2$ and apply the inverse-function theorem to their product of
exponentials.  An open subgroup of the connected group $G_2$ is $G_2$ itself.
This is exact generation by finite products, not only density
\cite{JurdjevicSussmann1972,SilvaLeiteCrouch1988,AgrachevSachkov2004}.
\end{proof}

The constraint in \Cref{thm:single} is the fixed relation $Y=\Ad_S X$, not merely the existence of two independent generators.

Set $n=(e_2+e_3)/\sqrt2$, and define
\[
a(t)=\cos(\sqrt2t)-n\sin(\sqrt2t),\qquad
b(t)=\cos(3\sqrt2t)+n\sin(3\sqrt2t).
\]

\begin{proposition}[Closed form and period of the generating pulse]
\label{prop:pulse}
The pulse has the stabilizer form
\begin{equation}
e^{tX}=\gamma_{a(t),b(t)}.
\label{eq:pulse}
\end{equation}
The infinitesimal generator $X$ has characteristic polynomial on $V$
\[
\chi_X(\lambda)=\lambda(\lambda^2+8)^2(\lambda^2+32),
\]
which realizes the canonical $\mathfrak g_2$ frequency pattern $\{0,\pm ia,\pm ib,\pm i(a+b)\}$ with $a=b=2\sqrt2$ \cite{ChemtovKarigiannis2022}.  The one-parameter group $e^{tX}$ has minimum positive period $T_X=\pi/\sqrt2$.
\end{proposition}

\begin{proof}
Differentiate \eqref{eq:gamma}.  Imaginary velocities $(q_a,q_b)$ act as
\[
x\mapsto q_ax-xq_a,
\qquad
y\mapsto q_by-yq_a.
\]
Choosing $q_a=-(e_2+e_3)$ and $q_b=3(e_2+e_3)$ gives $X$, proving \eqref{eq:pulse}.  At $T_X$, $(a,b)=(-1,-1)$, the diagonal kernel element of \eqref{eq:gamma}, so the pulse returns to the identity.  The frequencies $2\sqrt2,2\sqrt2,4\sqrt2$ show that no smaller positive time makes all rotation blocks trivial, proving minimality of the period.
\end{proof}

\section{An explicit factorization for symmetric-pair generation}\label{sec:universal-locus}
Fix $T\in\mathscr C$.  By \Cref{subsec:parity-domain}, the pair $(Z,TZT)$ is equivalent to independently chosen $A_+\in\mathfrak k_T$, $A_-\in\mathfrak m_T$, with $Z=A_++A_-$.
We construct a factored test for this parity-restricted two-generator problem on the full algebra.  The local domain $\mathfrak k_A$ is not used in its main proof.

Polynomial definability itself is standard.  Starting from $E_1=\Span\{Z,TZT\}$ and setting
$E_{r+1}=E_r+[Z,E_r]+[TZT,E_r]$, stabilization gives the generated algebra, and strict growth cannot continue past dimension fourteen.
Thus finitely many words of length at most fourteen suffice.  If $\Delta_j$ are all their maximal coordinate minors, then $\sum_j\Delta_j^2=0$ is an equation for failure over the reals.
The construction below replaces that unrestricted collection with a fifteen-dimensional symmetric Gram determinant and a degree-eight double-commutator expression.
Its value is the explicit factorization and the separation of the reducible and principal mechanisms; no optimal-degree or runtime claim is made.
Commutant methods themselves are classical \cite{ZeierSchulteHerbrueggen2011}.

Throughout this section, matrices carry the Frobenius inner product
\[
\langle U,V\rangle_F=\operatorname{tr}(U^{\mathsf T}V).
\]
Its restriction to skew-symmetric matrices is the trace metric already used, whereas on symmetric matrices it is $\operatorname{tr}(UV)$.  For $T\in\mathscr C$ put
\[
\mathscr S_0=\{Q\in\operatorname{End}(V):Q^{\mathsf T}=Q,\ \operatorname{tr}Q=0\},
\qquad
\mathscr E_T=\{Q\in\mathscr S_0:TQT=Q\}.
\]
Since the eigenspaces of $T$ have dimensions three and four,
\[
\dim\mathscr S_0=27,\qquad \dim\mathscr E_T=6+10-1=15.
\]
For an algebra of skew-symmetric matrices, the orthogonal complement of every invariant subspace is also invariant.  We call such a representation \emph{reducible} when it has a nonzero proper invariant subspace.

The \emph{commutant} of a family of operators is the set of operators commuting with every member of that family.  The space $\mathscr E_T$ consists of symmetric trace-zero operators fixed by conjugation with $T$; equivalently, its elements commute with $T$.  A polynomial on $\mathscr C\times\mathfrak g_2$ means the restriction of a real polynomial in matrix entries.  All degrees stated for $\mathcal R_{30}$, $\mathcal H_8$, and $\mathcal U_{38}$ are degrees in $Z$, with $T$ held fixed, not total degrees in both variables.

\subsection{A symmetric commutant detects reducibility}

A proper invariant subspace gives an orthogonal projector that commutes with
both generators.  The next lemma averages that projector under the involution
without losing its non-scalar part.  This replaces a search over invariant
subspaces by linear equations for a symmetric matrix.

\begin{lemma}[Odd-dimensional involution test]
\label{lem:odd-commutant}
Let $V$ be an odd-dimensional real Euclidean space, $T$ an orthogonal involution, and $Z\in\mathfrak{so}(V)$.  Then $\Lie\{Z,TZT\}$ acts reducibly on $V$ if and only if there is a nonzero traceless symmetric operator $Q$ such that
\[
[Q,T]=[Q,Z]=0.
\]
\end{lemma}

\begin{proof}
Write $\mathfrak l=\Lie\{Z,TZT\}$.  If such a $Q$ exists, it also commutes with $TZT$.  A nonzero traceless symmetric operator is not scalar, and any one of its proper eigenspaces is invariant under both generators.

Conversely, suppose $\mathfrak l$ is reducible.  Choose an invariant subspace $W$ with $1\le d:=\dim W\le(\dim V-1)/2$, replacing $W$ by its orthogonal complement if necessary.  Let $P_W$ be its orthogonal projector.  It commutes with both generators.  The algebra $\mathfrak l$ is stable under $\Ad_T$, so $TP_WT$ also commutes with both.  Set
\[
Q=P_W+TP_WT-\frac{2d}{\dim V}I.
\]
Then $Q$ is symmetric, traceless, and commutes with $T$ and $Z$.  It is nonzero: $P_W+TP_WT$ has positive trace $2d$ and rank at most $2d<\dim V$, and hence cannot be a scalar operator.  The rank inequality is the only point at which odd dimension is used.
\end{proof}

For $T\in\mathscr C$ and $Z\in\mathfrak g_2$, choose any Frobenius-orthonormal basis $Q_1,\ldots,Q_{15}$ of $\mathscr E_T$ and define
\begin{equation}
\mathcal R_{30}(T,Z)
=\det\bigl(\langle[Z,Q_i],[Z,Q_j]\rangle_F\bigr)_{i,j=1}^{15}.
\label{eq:R30-gram}
\end{equation}
This definition is independent of that choice of orthonormal basis.  With an arbitrary basis, divide the determinant on the right by $\det(\langle Q_i,Q_j\rangle_F)$.

For a formula not requiring an eigenspace basis, define operators on $\mathscr S_0$ by
\[
\Pi_T^{\pm}(Q)=\tfrac12(Q\pm TQT),\qquad
\mathsf A_Z(Q)=[Z,Q].
\]

\begin{proposition}[A degree-thirty reducibility polynomial]
\label{prop:reducibility-polynomial}
The function $\mathcal R_{30}(T,Z)$ is nonnegative, is homogeneous of degree thirty in $Z$, and satisfies
\[
\mathcal R_{30}(T,Z)=0
\quad\Longleftrightarrow\quad
\Lie\{Z,TZT\}\text{ acts reducibly on }V.
\]
It is polynomial in the entries of $(T,Z)$ on $\mathscr C\times\mathfrak g_2$, without choosing an eigenspace frame.  The same function has the frame-free expression
\begin{equation}
\mathcal R_{30}(T,Z)
=\det_{\mathscr S_0}
\bigl(\Pi_T^- -\Pi_T^+\mathsf A_Z^2\Pi_T^+\bigr).
\label{eq:R30-framefree}
\end{equation}
Both formulas are invariant under simultaneous orthogonal conjugation of the matrix data.
\end{proposition}

\begin{proof}
The matrix in \eqref{eq:R30-gram} is the Gram matrix of the linear map
\[
\mathscr E_T\longrightarrow\mathscr S_0,\qquad Q\longmapsto[Z,Q].
\]
Its determinant is nonnegative and vanishes exactly when the map has a nonzero kernel.  \Cref{lem:odd-commutant} identifies that kernel condition with reducibility.  Each Gram entry is quadratic in $Z$, giving degree thirty.

The operator $\mathsf A_Z$ preserves $\mathscr S_0$ and is skew-adjoint for the Frobenius metric.  The orthogonal splitting $\mathscr S_0=\mathscr E_T\oplus\operatorname{im}\Pi_T^-$ has dimensions $15+12$.  Relative to it, the operator in \eqref{eq:R30-framefree} is the Gram operator of $\mathsf A_Z|_{\mathscr E_T}$ on the first summand and the identity on the second.  Its determinant is therefore \eqref{eq:R30-gram}.  The entries of $\Pi_T^\pm$ and $\mathsf A_Z$ are polynomial in $T,Z$, establishing the frame-free assertion.  Homogeneity of degree thirty is asserted on the reflection variety, where these are complementary projections.  Conjugation is an isometry of all the spaces and maps involved.
\end{proof}

The fifteen-dimensional formula is convenient for evaluation; the twenty-seven-dimensional formula shows that the coefficients vary polynomially with the reflection.  An explicit basis recipe and its normalization are given in Appendix~\ref{app:universal-certificates}.  Neither formula involves computing a generated Lie algebra.

\subsection{A degree-eight test for the low-dimensional branch}

The determinant $\mathcal R_{30}$ excludes reducible closures, but an irreducible
proper closure can still be a principal $\mathfrak{su}(2)$.  The compact maximal
list therefore leaves a question of dimension at most three.  That question
has an exact answer using only two double commutators.

For $T\in\mathscr C$ and $Z\in\mathfrak g_2$, put
\[
A_+=\tfrac12(Z+TZT),\qquad A_-=\tfrac12(Z-TZT),\qquad C=[A_+,A_-],
\]
\[
C_+=[A_+,C],\qquad C_-=[A_-,C].
\]
Define
\begin{equation}
\mathcal H_8(T,Z)
=\|A_-\|_F^2\|C_+\|_F^2
 +\|A_+\|_F^2\|C_-\|_F^2-2\|C\|_F^4.
\label{eq:H8}
\end{equation}

An invariant inner product on a real Lie algebra satisfies $\langle[X,Y],Z\rangle=\langle X,[Y,Z]\rangle$, equivalently every adjoint map is skew-adjoint.  An orthogonal involution preserves that inner product and has square equal to the identity.

More generally, if $\sigma$ is an orthogonal involution of a real Lie algebra
with a positive definite invariant inner product, use
$A_\pm=(Z\pm\sigma Z)/2$, $C=[A_+,A_-]$, and $C_\pm=[A_\pm,C]$ in the same
formula, and denote the resulting scalar by $\mathcal H_8(\sigma,Z)$.

\begin{lemma}[An exact small-closure criterion]
\label{lem:H8}
Let a finite-dimensional real Lie algebra have a positive definite invariant inner product and an orthogonal involutive automorphism $\sigma$.  For every $Z$ in that algebra, the expression $\mathcal H_8(\sigma,Z)$ defined above is nonnegative and vanishes if and only if
\[
\dim\Lie\{Z,\sigma Z\}\le3.
\]
If $C=0$, the closure is the abelian span of $A_+,A_-$.  If $C\ne0$ and $\mathcal H_8=0$, the closure is exactly $\Span\{A_+,A_-,C\}\cong\mathfrak{su}(2)$.
\end{lemma}

\begin{proof}
The vectors $A_+,A_-$ are orthogonal by parity.  Invariance gives
\[
\langle C,A_+\rangle=\langle C,A_-\rangle=0,
\quad
\langle C_+,A_-\rangle=-\|C\|^2,
\quad
\langle C_-,A_+\rangle=\|C\|^2.
\]
Thus $\mathcal H_8$ is the sum of the two nonnegative Gram determinants
\[
\det\operatorname{Gram}(A_-,C_+)
+\det\operatorname{Gram}(A_+,C_-).
\]
If $C=0$, both determinants vanish and the closure is abelian.  If $C\ne0$, then $A_+,A_-,C$ are nonzero and pairwise orthogonal.  Equality in the two Cauchy--Schwarz inequalities is equivalent to
\[
C_+=-\frac{\|C\|^2}{\|A_-\|^2}A_-,\qquad
C_-=\frac{\|C\|^2}{\|A_+\|^2}A_+.
\]
These relations close their span under brackets and give the compact three-dimensional simple algebra.  Conversely, if the closure has dimension three and $C\ne0$, this orthogonal triple spans it.  Parity and invariance force the two displayed relations.  This proves both directions, including every case in which one of the vectors vanishes.
\end{proof}

\subsection{A uniform exact equation for all reflections}

The two tests now exhaust the compact maximal-subalgebra alternatives:
reducible closures are detected by $\mathcal R_{30}$, and every remaining
proper closure has dimension at most three.  Their product consequently gives
one equation for the nongenerating set, including degenerate pulses.

For a quaternionic reflection $T\in\mathscr C$ and every $Z\in\mathfrak g_2$, set
\begin{equation}
\mathcal U_{38}(T,Z)=\mathcal R_{30}(T,Z)\mathcal H_8(T,Z).
\label{eq:U38}
\end{equation}

For a fixed local context $H_A$, also define the failed-pulse fiber
\[
\mathcal N_T=\{Z\in\mathfrak k_A:\Lie\{Z,TZT\}\ne\mathfrak g_2\}.
\]

\begin{theorem}[Reducibility--small-algebra factorization]
\label{thm:universal-locus}
The explicitly defined product $\mathcal U_{38}=\mathcal R_{30}\mathcal H_8$ is a nonnegative polynomial on $\mathscr C\times\mathfrak g_2$, homogeneous of degree thirty-eight in the pulse.  It gives the equivalence
\begin{equation}
\Lie\{Z,TZT\}=\mathfrak g_2
\quad\Longleftrightarrow\quad
\mathcal U_{38}(T,Z)>0.
\label{eq:universal-exact}
\end{equation}
In particular, for any prescribed local context and any prescribed reflection, the complete failed-pulse locus is
\begin{equation}
\mathcal N_T=\{Z\in\mathfrak k_A:\mathcal U_{38}(T,Z)=0\}.
\label{eq:NT-exact}
\end{equation}
There are exactly three mutually exclusive outcomes:
\begin{enumerate}[label=\textnormal{(\roman*)},leftmargin=8mm]
\item $\mathcal R_{30}=0$: the generated algebra acts reducibly and is proper;
\item $\mathcal R_{30}>0$ and $\mathcal H_8=0$: the generated algebra is a principal $\mathfrak{su}(2)$;
\item $\mathcal R_{30}>0$ and $\mathcal H_8>0$: the generated algebra is $\mathfrak g_2$.
\end{enumerate}
The criterion is invariant under simultaneous $G_2$ conjugation.
\end{theorem}

\begin{proof}
Fix $T\in\mathscr C$ and $Z\in\mathfrak g_2$, and write
$L=\Lie\{Z,TZT\}$.  By \Cref{prop:reducibility-polynomial,lem:H8},
\[
\mathcal R_{30}=0\ \Longleftrightarrow\ L\text{ is reducible on }V,
\qquad
\mathcal H_8=0\ \Longleftrightarrow\ \dim L\le3.
\]
Both factors are nonnegative.  If $L=\mathfrak g_2$, its seven-dimensional
representation is irreducible and its dimension is fourteen; hence both
factors are positive.

Conversely, suppose $L$ is proper and choose a maximal proper compact
subalgebra containing it.  The list in \Cref{subsec:compact-maximals} gives
three possibilities.  A vector stabilizer preserves a line in $V$, and a
quaternionic stabilizer preserves the splitting $3\oplus4$; containment in
either makes $L$ reducible and forces $\mathcal R_{30}=0$.  A principal
$\mathfrak{su}(2)$ has dimension three, so containment in it forces
$\mathcal H_8=0$.  Thus every proper closure makes the product vanish,
proving \eqref{eq:universal-exact} and its restriction
\eqref{eq:NT-exact}.

This reasoning also proves the three disjoint outcomes.  When
$\mathcal R_{30}>0$ and $\mathcal H_8=0$, the closure is irreducible.  It cannot
be contained in either reducible maximal type, so it lies in a principal
$\mathfrak{su}(2)$.  A proper subalgebra of $\mathfrak{su}(2)$ is abelian;
an abelian algebra of skew-symmetric operators in odd dimension has a common
fixed vector.  Thus $L$ equals that principal algebra.

By \Cref{prop:reducibility-polynomial} and the explicit bracket formula
\eqref{eq:H8}, the factors are polynomials on
$\mathscr C\times\mathfrak g_2$ of respective degrees thirty and eight in
$Z$.  Their product is homogeneous of degree thirty-eight.  Simultaneous
conjugation preserves the Frobenius metric, the brackets, and the spaces
$\mathscr E_T$, proving invariance.

\end{proof}

The formulas specify $\mathcal N_T$ for \emph{each} reflection, including exceptional relative positions and pulses with zero or repeated frequencies.  They are not a generic-only test and do not require choosing an invariant vector, a quaternionic context, or an iterated Lie-closure basis.  This is a uniform determinantal description; it does not assert that the irreducible components or orbit-type strata of $\mathcal N_T$ are independent of $T$.
The degree thirty-eight describes this explicit factorization only. There is no canonical bracket-minor polynomial against which it is compared, and no optimal-degree or complexity bound is asserted.

\subsection{The principal obstruction actually occurs}

The remaining irreducible proper outcome must be realized explicitly.  The following principal triple puts its first generator in the prescribed local algebra.

Define
\begin{align*}
K_{\rm p}&=\tfrac56D_{23}+\tfrac23D_{45}=Z(e_1,2e_1),\\
P_{\rm p}&=\tfrac{\sqrt6}{4}D_{14}
 +\tfrac{\sqrt{10}}6\bigl(\tfrac12D_{17}+D_{24}\bigr),\\
Q_{\rm p}&=\tfrac{\sqrt6}{4}D_{15}
 +\tfrac{\sqrt{10}}6\bigl(\tfrac12D_{16}-D_{25}\bigr).
\end{align*}
Put $N_{\rm p}=(K_{\rm p}+P_{\rm p})/\sqrt2$ and
\begin{equation}
T_*=I+\frac{136}{45}N_{\rm p}^2+\frac{10}{9}N_{\rm p}^4+\frac4{45}N_{\rm p}^6.
\label{eq:principal-reflection}
\end{equation}

For a triple with the cyclic commutator relations below, the quadratic Casimir in its action on $V$ is the operator $-(K_{\rm p}^2+P_{\rm p}^2+Q_{\rm p}^2)$.  On complexification, a spin-$j$ irreducible representation has dimension $2j+1$ and Casimir eigenvalue $j(j+1)$ in this normalization.  These conventions specify the irreducibility test in the proof.

\begin{proposition}[A moving reflection with principal closure]
\label{prop:principal-moving}
These derivations satisfy
\begin{equation}
[K_{\rm p},P_{\rm p}]=Q_{\rm p},\quad [P_{\rm p},Q_{\rm p}]=K_{\rm p},\quad [Q_{\rm p},K_{\rm p}]=P_{\rm p},
\qquad -(K_{\rm p}^2+P_{\rm p}^2+Q_{\rm p}^2)=12I_V.
\label{eq:principal-triple}
\end{equation}
Then $T_*\in\mathscr C$, $T_*(H_A)\ne H_A$, and
\[
T_*K_{\rm p}T_*=P_{\rm p},\qquad
\Lie\{K_{\rm p},T_*K_{\rm p}T_*\}=\Span\{K_{\rm p},P_{\rm p},Q_{\rm p}\}
\]
is a principal $\mathfrak{su}(2)$.  In particular,
\[
\mathcal R_{30}(T_*,K_{\rm p})>0,\qquad \mathcal H_8(T_*,K_{\rm p})=0.
\]
\end{proposition}

\begin{proof}
The identities in \eqref{eq:principal-triple} follow by substituting the fixed Fano derivations.  They give a compact $\mathfrak{su}(2)$ representation whose Casimir is $12I$.  In the normalization of the displayed brackets, the Casimir on a spin-$j$ irreducible module is $j(j+1)$.  Thus only spin three occurs; since $\dim V=7$, the representation is irreducible, and the subalgebra is principal.  Explicit principal triples in compact $\mathfrak g_2$ are classical; see, for example, Draper and Mart\'in-Gonz\'alez \cite{DraperMartin2025}. The displayed normalization is chosen to put $K_{\rm p}$ in the specified local algebra.

The element $N_{\rm p}$ is conjugate to $K_{\rm p}$ inside this compact subgroup, so its spectrum is $0,\pm i,\pm2i,\pm3i$.  The polynomial on the right of \eqref{eq:principal-reflection} takes values $1,-1,1,-1$ on the eigenspaces of $N_{\rm p}^2$ with eigenvalues $0,-1,-4,-9$.  Consequently it equals $\exp(\pi N_{\rm p})$, is an involution in $G_2$, and has trace $-1$ on $V$.  In the adjoint three-dimensional algebra it is the half-turn about the axis $K_{\rm p}+P_{\rm p}$, hence sends $K_{\rm p}$ to $P_{\rm p}$.  The latter does not belong to $\mathfrak k_A$, so $T_*$ cannot preserve $H_A$.  The two generators span the indicated principal algebra together with their bracket.  \Cref{prop:reducibility-polynomial} and \Cref{lem:H8} give the final two assertions.  A sparse matrix for $T_*$ is printed in Appendix~\ref{app:universal-certificates}.
\end{proof}

The example belongs to a continuous family.  For
\[
T_\theta=\exp\bigl(\pi(\cos\theta K_{\rm p}+\sin\theta P_{\rm p})\bigr),
\]
all $T_\theta$ are quaternionic reflections and
\[
T_\theta K_{\rm p}T_\theta=\cos(2\theta)K_{\rm p}+\sin(2\theta)P_{\rm p}.
\]
Whenever $\sin(2\theta)\ne0$, the reflection moves $H_A$ and the constrained closure is principal.  Thus the degree-eight factor in the universal criterion cannot be discarded.  This family also explains why three distinct frequencies, even with the principal ratio, do not determine the generated algebra: the relative reflection still matters.

\subsection{Exactly which pulses admit a principal reflected closure}
The example gives the sharp spectral boundary, using the classical vector stabilizer $\SU(3)$ and its standard action on the orthogonal complex three-space \cite{DraperMartin2025,Baez2002}.

\begin{corollary}[Spectral realizability of the principal obstruction]\label{cor:universal-spectral}
For a nonzero $Z\in\mathfrak g_2$, the following conditions are equivalent:
\begin{enumerate}[label=\textnormal{(\roman*)},leftmargin=9mm]
\item the positive rotation frequencies of $Z$ have ratio $1:2:3$;
\item $Z$ belongs to a principal $\mathfrak{su}(2)$ subalgebra of $\mathfrak g_2$;
\item there is $T\in\mathscr C$ such that $\Lie\{Z,TZT\}$ is a principal $\mathfrak{su}(2)$.
\end{enumerate}
For $Z=Z(q_a,q_b)\in\mathfrak k_A$ with $\alpha>0$, these conditions are equivalent to
\[
\frac{\beta}{\alpha}\in\left\{2,5,\frac13\right\}.
\]
Every reflection realizing (iii) for such a local $Z$ moves $H_A$.
Outside the principal spectrum, including $Z=0$, generation for any $T\in\mathscr C$ is equivalent to $\mathcal R_{30}(T,Z)>0$.
\end{corollary}
\begin{proof}
A nonzero element of the irreducible spin-three representation has positive frequencies in the ratio $1:2:3$, so (iii) implies (ii) and (ii) implies (i).
To prove the converse, choose a unit vector $v\in\ker Z$.  Under (i) the kernel is a line.  The vector stabilizer $\mathfrak g_2^v\cong\mathfrak{su}(3)$ acts on $v^\perp$ as a complex three-space with complex structure $x\mapsto v\times x$.
The skew-Hermitian operator $Z|_{v^\perp}$ has eigenvalues $it_1,it_2,it_3$, where
\[
t_1+t_2+t_3=0,\qquad \{|t_1|,|t_2|,|t_3|\}=\{\lambda,2\lambda,3\lambda\}.
\]
The only signed triples with these absolute values and zero sum are permutations of $\lambda(1,2,-3)$ and its negative.
Conjugate $v$ to $e_1$ using $G_2$ and diagonalize in its $\SU(3)$ stabilizer.  Unitary diagonalization may be taken in $\SU(3)$ by changing the scalar phase of the diagonalizing matrix.
The element $K_{\rm p}=Z(e_1,2e_1)$ has the signed triple $(2,1,-3)$ in this complex structure.
It follows that $Z=c\Ad_g K_{\rm p}$ for some $g\in G_2$ and $c\ne0$.
Taking $T=gT_*g^{-1}$ and using \Cref{prop:principal-moving} gives
\[
TZT=c\Ad_g P_{\rm p},\qquad
\Lie\{Z,TZT\}=\Ad_g\Span\{K_{\rm p},P_{\rm p},Q_{\rm p}\},
\]
which proves (iii).  This argument needs conjugacy only on the principal spectral cone, not a general inverse-spectral claim.

For local $Z$ with $\alpha>0$, divide the frequencies $2\alpha,|\beta-\alpha|,\beta+\alpha$ by $\alpha$ and put $r=\beta/\alpha$.
For $r>1$ their signed complex triple is $(2,r-1,-r-1)$; the ratio $1:2:3$ gives $r=2$ or $r=5$.
For $0\le r<1$, the largest frequency is $2$, and $1-r$, $1+r$ must be $2/3$, $4/3$, giving $r=1/3$.
The cases $r=1$ or $\alpha=0$ have a zero frequency and are not principal.
If a realizing $T$ preserved $H_A$, its generated principal algebra would lie in $\mathfrak k_A$, which preserves the proper subspace $P_A$.  This contradicts irreducibility.
Finally, without a principal spectrum the irreducible proper case of \Cref{thm:universal-locus} is impossible, leaving exactly the asserted reducibility test.
\end{proof}

The existential statement describes the projection of the principal-pair locus onto the pulse variable.  It does not say that every reflection works for such a pulse.  The following subsection describes the realizing-reflection fiber itself.
For a local pulse the projection is the nonzero part of
\[
(\beta^2-4\alpha^2)(\beta^2-25\alpha^2)(9\beta^2-\alpha^2)=0.
\]


\subsection{All reflections realizing a fixed principal pulse}
The preceding corollary describes the projection onto the pulse variable.  We now fix a nonzero pulse with principal spectrum and determine every realizing reflection.
Put
\[
\mathfrak s_{\rm p}=\Span\{K_{\rm p},P_{\rm p},Q_{\rm p}\},\qquad
C_{\rm p}=\{u\in G_2:\Ad_uK_{\rm p}=K_{\rm p}\},
\]
and let $S_{\rm p}\subset G_2$ be the connected subgroup with Lie algebra $\mathfrak s_{\rm p}$.
Its spin-three action identifies it with $\SO(3)$.
For a Lie subalgebra $\mathfrak s$, its normalizer is $N_{G_2}(\mathfrak s)=\{u\in G_2:\Ad_u\mathfrak s=\mathfrak s\}$.
For any fixed $Z$, define
\[
\mathscr P_Z=\{T\in\mathscr C:\Lie\{Z,TZT\}\text{ is principal}\}.
\]
For $Z$ with principal spectrum choose, once and for all, $c\ne0$ and $g\in G_2$ with $Z=c\Ad_gK_{\rm p}$, as in \Cref{cor:universal-spectral}.
The family $T_\theta$ is the one defined after \Cref{prop:principal-moving}.

\begin{proposition}[The realizing-reflection fiber]\label{prop:principal-reflection-fiber}
The centralizer $C_{\rm p}$ is a two-dimensional torus.  For the fixed choice $Z=c\Ad_gK_{\rm p}$,
\[
\mathscr P_Z
 =\{g uT_\theta u^{-1}g^{-1}:u\in C_{\rm p},\quad 0<\theta<\pi/2\}.
\]
Every reflection on the right has a unique pair $(u,\theta)$.  This parametrization is a real-analytic diffeomorphism onto $\mathscr P_Z$ with the subspace topology inherited from $\mathscr C$.  The fiber is therefore an embedded real-analytic submanifold of dimension three and codimension five in the eight-dimensional reflection manifold $\mathscr C$.
If $Z\in\mathfrak k_A$, all these reflections move $H_A$.
\end{proposition}
\begin{proof}
Scaling and conjugation reduce the proof to $Z=K_{\rm p}$.
First compute its group centralizer.  The kernel line is $\RR e_1$.  An element centralizing $K_{\rm p}$ must preserve this line.
If it sent $e_1$ to $-e_1$, its restriction to $e_1^\perp$ would be conjugate-linear for $C_{e_1}$, since it preserves the cross product.
Commutation with $K_{\rm p}$ would then identify the signed complex spectrum $(2,1,-3)$ with its negative, which is impossible.
It therefore fixes $e_1$ and lies in its $\SU(3)$ stabilizer.
There $K_{\rm p}$ has three distinct complex eigenvalues, so its centralizer is the diagonal determinant-one unitary torus.  Hence $C_{\rm p}\cong(S^1)^2$.

The group centralizer of $\mathfrak s_{\rm p}$ in $G_2$ is trivial.  Indeed, its real seven-dimensional action is irreducible.
A commuting symmetric operator is scalar, and a commuting skew-symmetric operator in odd dimension has a nonzero kernel and must vanish by irreducibility.
Thus the real commutant consists of scalars; its orthogonal elements are $\pm I_7$, of which only $I_7$ lies in $\SO(7)$.
Every automorphism of $\mathfrak{su}(2)$ is inner and is realized by $S_{\rm p}$, so
\[
N_{G_2}(\mathfrak s_{\rm p})=S_{\rm p}.
\]
Every principal algebra containing $K_{\rm p}$ is a $C_{\rm p}$-conjugate of $\mathfrak s_{\rm p}$.
To see this, choose $h$ carrying $\mathfrak s_{\rm p}$ to such an algebra, using the compact conjugacy in \Cref{subsec:compact-maximals}.
Then $\Ad_{h^{-1}}K_{\rm p}\in\mathfrak s_{\rm p}$ has the same trace norm as $K_{\rm p}$.
Rotations in $S_{\rm p}$ are transitive on that sphere, so $\Ad_vK_{\rm p}=\Ad_{h^{-1}}K_{\rm p}$ for some $v\in S_{\rm p}$.
The element $hv$ centralizes $K_{\rm p}$ and carries $\mathfrak s_{\rm p}$ to the prescribed algebra.

If $T\in\mathscr P_{K_{\rm p}}$, it normalizes its generated principal algebra: conjugation by $T$ interchanges the two generators.
Choose $u\in C_{\rm p}$ with $\Lie\{K_{\rm p},TK_{\rm p}T\}=\Ad_u\mathfrak s_{\rm p}$, as in the preceding paragraph.
Then $u^{-1}Tu\in N_{G_2}(\mathfrak s_{\rm p})=S_{\rm p}$ is a nonidentity involution.
It is therefore a half-turn about an unoriented axis in $S_{\rm p}\cong\SO(3)$.  Write its unit axis in the normalization $(K_{\rm p},P_{\rm p},Q_{\rm p})$ as
$\xi_1K_{\rm p}+\xi_2P_{\rm p}+\xi_3Q_{\rm p}$, where $\sum\xi_j^2=1$.
Its conjugate of $K_{\rm p}$ is independent of $K_{\rm p}$ exactly when $0<|\xi_1|<1$.
Choose the sign of the axis so that $\xi_1>0$ and rotate its transverse direction by $\exp(tK_{\rm p})\in C_{\rm p}$.
The axis becomes $\cos\theta K_{\rm p}+\sin\theta P_{\rm p}$ with $0<\theta<\pi/2$.
Conversely, every such half-turn generates $\mathfrak s_{\rm p}$ with $K_{\rm p}$ and is a quaternionic reflection, by the spin-three weights.
This proves the asserted equality of sets.

The angle is recovered from
\[
\frac{\langle K_{\rm p},TK_{\rm p}T\rangle_{\rm tr}}{\|K_{\rm p}\|_{\rm tr}^2}=\cos(2\theta).
\]
If two centralizer elements give the same reflection at this angle, their quotient commutes with $K_{\rm p}$ and $T_\theta$.
It therefore commutes with $K_{\rm p}$ and $T_\theta K_{\rm p}T_\theta$, hence with all of $\mathfrak s_{\rm p}$, and is the identity.
The parametrization is thus injective.  Its torus-orbit differential is injective by the same stabilizer argument; the angle has nonzero derivative in the displayed scalar on the open interval.
It is an analytic immersion.  On every compact subinterval the domain is compact, and the displayed scalar recovers $\theta$ continuously in the subspace topology of $\mathscr C$.  It follows that the map is an embedding with analytic inverse onto its image.  Its codimension is $\dim\mathscr C-3=8-3=5$.
The last assertion follows from \Cref{cor:universal-spectral}.
\end{proof}

The set $\mathscr P_Z$ is defined intrinsically by $Z$ and does not depend on the choice of $(c,g)$; only the coordinates on it do.
For fixed $c$, replacing $g$ by $gv$ with $v\in C_{\rm p}$ changes the coordinate of the same reflection from $u$ to $v^{-1}u$, with $\theta$ unchanged.
Arbitrary admissible choices of $(c,g)$ still give the same set, including choices with the opposite sign of $c$.

The single curve $T_\theta$ does not by itself exhaust the fiber: it fixes one plane of axes in one principal subgroup.  Centralizer conjugation supplies the other axes and the other principal subgroups containing the pulse.
For example, conjugating $T_*$ by $u_0=\exp(\pi K_{\rm p}/2)$ gives the reflection $T_{\rm az}$ carrying $K_{\rm p}$ to $Q_{\rm p}$, not to a vector in $\Span\{K_{\rm p},P_{\rm p}\}$.  Its explicit matrix is printed in Appendix~\ref{app:universal-certificates}.
The result concerns only the principal branch; it does not classify all reducible failure fibers for arbitrary relative positions.

\paragraph{Dimension in both variables.}
These fibers fit into the principal-pair locus
\[
\mathscr M_{\rm p}
 =\{(T,Z)\in\mathscr C\times\mathfrak g_2:
       \Lie\{Z,TZT\}\text{ is principal}\}.
\]
The zero pulse is excluded by this definition.  In fact the map
\[
\begin{split}
\Phi:G_2\times(0,\infty)\times(0,\pi/2)&\longrightarrow\mathscr M_{\rm p},\\
(g,r,\theta)&\longmapsto
       (gT_\theta g^{-1},\ r\Ad_g K_{\rm p})
\end{split}
\]
is a real-analytic diffeomorphism onto the locus with its subspace topology.
To obtain a positive coefficient $r$ in \Cref{cor:universal-spectral}, absorb a negative sign by the half-turn $\exp(\pi P_{\rm p})$, whose adjoint sends $K_{\rm p}$ to $-K_{\rm p}$.
Then \Cref{prop:principal-reflection-fiber} gives surjectivity of $\Phi$.
For a pair in its image, the two parameters are recovered by
\[
r=\frac{\|Z\|_{\rm tr}}{\|K_{\rm p}\|_{\rm tr}},\qquad
\chi(T,Z):=\frac{\langle Z,TZT\rangle_{\rm tr}}{\|Z\|_{\rm tr}^2}
                 =\cos(2\theta).
\]
Once $r,\theta$ are fixed, two choices of $g$ differ by an element commuting with $K_{\rm p}$ and $T_\theta$, hence by the identity as proved above.
The differential is injective for the same reason: the differentials of $r$ and $\chi$ first eliminate the two parameter directions, after which a zero orbit tangent centralizes $K_{\rm p}$ and $T_\theta$.
The derivative of $\cos(2\theta)$ is nonzero on the open interval.
Compactness of $G_2$, together with the continuous recovery of $r,\theta$, makes this injective analytic immersion an embedding with analytic inverse.
Consequently
\[
\dim_{\RR}\mathscr M_{\rm p}=14+1+1=16,
\qquad
\operatorname{codim}_{\mathscr C\times\mathfrak g_2}\mathscr M_{\rm p}
       =22-16=6.
\]
For comparison, its projection onto the pulse variable is
$(0,\infty)\times G_2/C_{\rm p}$, of dimension $1+14-2=13$;
the fibers have dimension three, as in \Cref{prop:principal-reflection-fiber}.
The full pair locus is not a single orbit of simultaneous $G_2$ conjugation: $r$ and $\chi$ are two independent invariants.
For each fixed $(r,\theta)$ it is instead one free $G_2$ orbit.
These statements concern the full pulse domain $\mathfrak g_2$, not an additional restriction to $\mathfrak k_A$.


\section{The fixed involution and the nongenerating locus}
\label{sec:failure}

Because the second generator is $\Ad_SZ$, every generated algebra is invariant under $\Ad_S$.  This symmetry strongly restricts its possible proper over-algebras.  We first record the elementary involution argument and then identify the corresponding $6+8$ symmetric decomposition of $\mathfrak g_2$.

In this section a context is $S$-invariant when $S(H)=H$, not necessarily when $S|_H=I_H$.  A Lie subalgebra is $\operatorname{Ad}_S$-stable when conjugation by $S$ maps it to itself.  The words invariant and stable below refer to these setwise conditions.

\subsection{Symmetric-pair reduction and genericity}
We record the elementary parity observation from \Cref{subsec:parity-domain} in the general notation used for the fixed-fiber calculations; compactness is not needed.

For an involutive automorphism $\sigma$ of a real Lie algebra and an element $Z$, write
\[
Z^+=\tfrac12(Z+\sigma Z),\qquad Z^-=\tfrac12(Z-\sigma Z).
\]

\begin{lemma}[Generation under an involution]
\label{lem:general-symmetric-pair}
Let $\mathfrak g$ be a finite-dimensional real Lie algebra and $\sigma$ an involutive automorphism.  For every $Z\in\mathfrak g$,
\[
\Lie\{Z,\sigma Z\}=\Lie\{Z^+,Z^-\},
\]
and this Lie algebra is $\sigma$-stable.
\end{lemma}

\begin{proof}
The pairs $(Z,\sigma Z)$ and $(Z^+,Z^-)$ have the same linear span, so they generate the same Lie algebra.  Since $\sigma Z^+=Z^+$ and $\sigma Z^-=-Z^-$, every homogeneous bracket is carried to itself up to its parity sign, proving stability.
\end{proof}

The later maximal-subalgebra argument chooses an ordinary maximal proper subalgebra and intersects it with its image under the involution.

Let
\[
H_S=\RR1\oplus\Span\left\{e_1,\frac{e_2+e_4}{\sqrt2},\frac{e_3+e_5}{\sqrt2}\right\}.
\]
The displayed imaginary frame is associative, and the fixed switch satisfies
\[
S|_{H_S}=I,\qquad S|_{H_S^\perp}=-I.
\]
Thus $S$ is the orthogonal reflection that is $+I$ on the quaternionic subalgebra $H_S$ and $-I$ on $H_S^\perp$.

For the fixed involution, use $\sigma=\Ad_S$ in the preceding notation and set
$\mathfrak m_{H_S}=\{B\in\mathfrak g_2:\Ad_SB=-B\}$.
In the bracket relations of the following proposition, $\mathfrak k$ and $\mathfrak m$ abbreviate $\mathfrak k_{H_S}$ and $\mathfrak m_{H_S}$, respectively, rather than the subspaces associated with $H_0$ in \Cref{sec:g2closure}.  The projections $Z^+$ and $Z^-$ are onto these two eigenspaces of $\operatorname{Ad}_S$ on $\mathfrak g_2$, not onto the eigenspaces of $S$ on $V$.

\begin{proposition}[Symmetric-pair form of the constrained problem]
\label{prop:symmetric-switch}
The involution $\sigma=\Ad_S$ gives the symmetric decomposition
\[
\gtwo=\mathfrak k_{H_S}\oplus\mathfrak m_{H_S},
\qquad \dim\mathfrak k_{H_S}=6,
\quad \dim\mathfrak m_{H_S}=8,
\]
where $\mathfrak k_{H_S}$ and $\mathfrak m_{H_S}$ are the $+1$ and $-1$ eigenspaces of $\Ad_S$, respectively.  They satisfy
\[
[\mathfrak k,\mathfrak k]\subseteq\mathfrak k,
\qquad [\mathfrak k,\mathfrak m]\subseteq\mathfrak m,
\qquad [\mathfrak m,\mathfrak m]\subseteq\mathfrak k.
\]
For every $Z\in\mathfrak g_2$,
\[
\Lie\{Z,\Ad_SZ\}=\Lie\{Z^+,Z^-\},
\]
and this Lie algebra is $\Ad_S$-stable.
\end{proposition}

\begin{proof}
The three displayed imaginary generators of $H_S$ form an oriented quaternionic triple.  Applying the signed permutation $S$ to this frame and to its orthogonal complement gives the stated $\pm1$ eigenspaces.  Hence the fixed algebra of $\Ad_S$ is $\mathfrak k_{H_S}\cong\mathfrak{so}(4)$.  The symmetric-pair bracket relations follow from the eigenspace decomposition of an involutive Lie-algebra automorphism.  The final closure identity and its $\Ad_S$-stability are the specialization of \Cref{lem:general-symmetric-pair}.
\end{proof}

A subset of a real coordinate space is real algebraic when it is a common zero set of real polynomials; a real Zariski-open set is the complement of such a set.  The same terminology on a real algebraic manifold means the relative topology.  Full Lebesgue measure means that the complement has measure zero.

For $Z\in\mathfrak k_A$, let $M_{58}(Z)$ be the coordinate matrix whose columns are
$W_i(Z,\Ad_SZ)$ in \eqref{eq:g2basis}, and define $\Delta_{58}(Z)=\det M_{58}(Z)$.

\begin{remark}[Zariski-open success]
\label{thm:generic}
The polynomial $\Delta_{58}$ is homogeneous of degree $58$ and nonzero at $X$.
Its nonvanishing locus is a nonempty real Zariski-open subset of the generating set
$\mathcal G$; in particular, $\mathcal G$ has full Lebesgue measure in $\mathfrak k_A$.
\end{remark}

\begin{proof}
The recursion for the fourteen words gives their degrees in the coordinates of $Z$:
\[
(1,1,2,3,3,4,4,4,5,5,6,6,7,7).
\]
Their sum is $58$, so every term in $\det M_{58}(Z)$ has degree $58$.
By \eqref{eq:bracketdet}, $\Delta_{58}(X)=2^{88}3^9\ne0$.
Whenever $\Delta_{58}(Z)\ne0$, the fourteen columns span $\mathfrak g_2$ and
therefore $Z,\Ad_SZ$ generate it.  The complement of this sufficient open set
is the zero set of a nonzero real polynomial and has measure zero.
This proves the full-measure claim without asserting that every zero of
$\Delta_{58}$ is nongenerating.
\end{proof}

\subsection{Vector-stabilizer failures}
A \emph{common fixed vector} for the generated infinitesimal action means a nonzero $w\in V$ annihilated by both generators: $Zw=(\operatorname{Ad}_S Z)w=0$.  Equivalently, all their exponential flows fix $w$.  This differs from a subspace merely being preserved setwise.

For $0\ne w\in V$, write
\[
\mathfrak g_2^w=\{A\in\mathfrak g_2:Aw=0\}.
\]
After rescaling $w$ to unit length this is the standard maximal stabilizer $\mathfrak{su}(3)$.  The $\Ad_S$-stability from \Cref{prop:symmetric-switch} turns containment in such a stabilizer into a common-kernel condition.

For $Z=Z(q_a,q_b)\in\mathfrak k_A$, write $\alpha=\|q_a\|$ and $\beta=\|q_b\|$.

The rotation frequencies of a skew operator on $V$ are the three nonnegative numbers $\omega_j$, counted with multiplicity, for which its complex spectrum is $0,\pm i\omega_1,\pm i\omega_2,\pm i\omega_3$.  Below, \emph{spectrally regular} means that these three numbers are positive and pairwise distinct; this is the frequency condition used here, not an unstated regularity assumption.  A rank stratum is the subset on which the rank of this seven-dimensional operator is fixed.
For each $0\ne w\in V$, define the local annihilator
\[
\mathcal N_w=\{Z\in\mathfrak k_A:Zw=0\}.
\]
We shall use the fixed witness
\begin{equation}
Z_*=D_{23}+2D_{13}-4D_{45}-4D_{46}.
\label{eq:su3-witness}
\end{equation}

\begin{proposition}[Complete common-vector criterion and the rank-six slice]
\label{prop:failure}
Let $Z\in\mathfrak k_A$.  The following conditions are equivalent:
\begin{enumerate}[label=\textnormal{(\roman*)},leftmargin=9mm]
\item $\Lie\{Z,\Ad_SZ\}\subseteq\gtwo^w$ for some $0\ne w\in V$;
\item $\ker Z\cap S(\ker Z)\ne\{0\}$;
\item $\ker Z$ contains a nonzero vector in $P_{H_S}$ or a nonzero vector in $H_S^\perp$.
\end{enumerate}
Its three rotation frequencies, allowing zeros, are
\[
2\alpha,\qquad |\beta-\alpha|,\qquad \beta+\alpha.
\]
On the rank-six stratum $\alpha\ne0$, $\beta\ne\alpha$, one has $\ker Z=\RR q_a\subset P_A$, and
\[
\left\{Z\in\mathfrak k_A:\operatorname{rank}Z=6,
\ \Lie\{Z,\Ad_SZ\}\subseteq\gtwo^w\text{ for some }w\ne0\right\}
=\mathcal N_{e_1}\cap\{\operatorname{rank}Z=6\},
\]
where
\[
\mathcal N_{e_1}
=\Span\left\{D_{23},D_{45},D_{46}-\tfrac12D_{13},D_{47}+\tfrac12D_{12}\right\}.
\]
A nonempty Zariski-open subset of this four-plane generates the full maximal stabilizer $\gtwo^{e_1}\cong\mathfrak{su}(3)$.  The element $Z_*$ satisfies
$\Lie\{Z_*,\Ad_SZ_*\}=\mathfrak g_2^{e_1}$.
\end{proposition}

\begin{proof}
Since $S=S^{-1}$,
\[
(\Ad_SZ)w=0\iff Z(Sw)=0\iff w\in S(\ker Z),
\]
which proves the equivalence of the first two conditions.  The common-kernel space is $S$-invariant and therefore contains an eigenvector of the involution, proving the third condition.  The differential of \eqref{eq:gamma} yields
\[
\chi_Z(\lambda)=\lambda(\lambda^2+4\alpha^2)
(\lambda^2+(\beta-\alpha)^2)(\lambda^2+(\beta+\alpha)^2).
\]
On rank six the unique kernel is $\RR q_a$.  Because
\[
P_A\cap P_{H_S}=\RR e_1,
\qquad P_A\cap H_S^\perp=0,
\]
the criterion reduces to $Ze_1=0$.  If
\[
Z=x_0D_{12}+x_1D_{13}+x_2D_{23}+x_3D_{45}+x_4D_{46}+x_5D_{47},
\]
then $Ze_1=0$ is exactly $x_5=2x_0$ and $x_4=-2x_1$, giving the displayed basis.  The full-$\mathfrak{su}(3)$ witness and determinant are recorded in Appendix~\ref{app:failure-certificate}.
\end{proof}

For the rank-four common-vector calculation, fix
$Z'=D_{12}-3D_{13}+2D_{47}$.

\begin{proposition}[A ruled rank-four vector-stabilizer component]
\label{prop:rankfour-failure}
For $0\ne w=ae_6+be_7$, the linear space $\mathcal N_w$ has dimension three and is spanned by
\begin{align*}
&D_{45}-D_{23},\\
&3abD_{12}-(2a^2-b^2)D_{13}+(a^2+b^2)D_{46},\\
&-(a^2-2b^2)D_{12}-3abD_{13}+(a^2+b^2)D_{47}.
\end{align*}
The union over $[a:b]\in\RR P^1$ is a ruled four-dimensional cone of common-vector failures.  At $a=b=1$, the witness $Z'$ has
\[
\ker Z'=\Span\{e_2,e_4+e_5,e_6+e_7\},
\qquad Z'e_1=-12e_3,
\]
and
\[
\Lie\{Z',\Ad_SZ'\}=\gtwo^{e_6+e_7}\cong\mathfrak{su}(3).
\]
\end{proposition}

\begin{proof}
In the ordered basis $(D_{12},D_{13},D_{23},D_{45},D_{46},D_{47})$, the map $Z\mapsto Zw$ has exact rank three for $(a,b)\ne(0,0)$.  Solving it gives the three displayed generators.  The incidence space over $\mathbb{RP}^1$ has dimension four, giving an upper bound for the dimension of its image.  At the displayed element $Z'$, its kernel intersects $\Span\{e_6,e_7\}$ in exactly $\RR(e_6+e_7)$.  This uniqueness persists on a relatively open part of the incidence space because the restriction to that two-plane has rank one there.  The projection to the derivation therefore has zero-dimensional fibers on that open part, proving that its image has dimension four.  The stated kernel and $Z'e_1$ follow by substitution, while the eight-dimensional closure is certified by the determinant in Appendix~\ref{app:projective-failure-certificates}.
\end{proof}

The positive and negative eigenvectors used in the incidence description are parametrized by
\[
w_+(a,b,c)=ae_1+b(e_2+e_4)+c(e_3+e_5),
\]
\[
w_-(a,b,r,t)=a(e_2-e_4)+b(e_3-e_5)+re_6+te_7.
\]
Let $Z_+(a,b,c)$ and $Z_-(a,b,r,t)$ denote the homogeneous cubic expressions
specified in Appendix~\ref{app:rankfour-parametrization}.
Their annihilator properties are proved below; the formulas fix the chosen representatives.

\begin{proposition}[Complete incidence parametrization of rank-four vector failures]
\label{prop:rankfour-complete}
Let $Z\in\mathfrak k_A$ be nonzero of rank four.  Then
\[
\Lie\{Z,\Ad_SZ\}\subseteq\gtwo^w
\quad\text{for some }w\ne0
\]
if and only if $Z$ belongs to at least one of the following four explicitly parametrized families:
\begin{enumerate}[label=\textnormal{(\alph*)},leftmargin=9mm]
\item $\mathcal N_{e_1}\cap\{\operatorname{rank}Z=4\}$;
\item the cubic line $\RR Z_+(a,b,c)$ over an eigenline
\[
[w_+(a,b,c)]\in\RR P(P_{H_S}),
\qquad (b,c)\ne(0,0);
\]
\item the cubic line $\RR Z_-(a,b,r,t)$ over an eigenline
\[
[w_-(a,b,r,t)]\in\RR P(H_S^\perp),
\qquad (a,b)\ne(0,0);
\]
\item the three-plane $\mathcal N_{re_6+te_7}$ from \Cref{prop:rankfour-failure}, for $[r:t]\in\RR P^1$.
\end{enumerate}
The explicit homogeneous cubic generators $Z_+$ and $Z_-$ are printed in Appendix~\ref{app:rankfour-parametrization}.  Away from the exceptional eigenlines $[e_1]$ and $\RR P\Span\{e_6,e_7\}$, the annihilator fiber in $\mathfrak k_A$ is one-dimensional.  Consequently the entire rank-four common-vector mechanism is closed by this incidence parametrization.
\end{proposition}

\begin{proof}
Write an $S$-eigenvector as $w=p+ye_4$, with $p\in P_A$ and quaternion coefficient $y\in H_A$.  Since every $Z\in\mathfrak k_A$ is block diagonal for $P_A\oplus H_A^\perp$, the equation $Zw=0$ is equivalent, in the quaternionic coordinates of \eqref{eq:gamma}, to
\[
[q_a,p]=0,
\qquad q_by=yq_a.
\]
If $p\ne0$ and $y\ne0$, the first equation forces $q_a$ to be a scalar multiple of $p$, and the second then determines $q_b=yq_ay^{-1}$; the annihilator is therefore a line.  For a positive eigenvector, $y=0$ occurs only on $\RR e_1$, giving the four-plane $\mathcal N_{e_1}$.  For a negative eigenvector, $p=0$ occurs exactly on $\Span\{e_6,e_7\}$, where $q_a$ is free and the annihilator has dimension three.  Substitution gives the cubic generators in Appendix~\ref{app:rankfour-parametrization}.

Every nonzero solution outside the $e_1$ fiber has a kernel vector with nonzero components in both invariant blocks, or a nonzero kernel in the four-dimensional normal block.  Skew symmetry then forces a kernel of dimension at least three.  Since a nonzero element of $\mathfrak g_2$ cannot have rank two \cite{ChemtovKarigiannis2022}, it has rank four.  Conversely, \Cref{prop:failure} says that every rank-four vector failure contains an $S$-eigenvector, which lies in one of the four cases above.
\end{proof}

\begin{remark}[Associative kernel reduction]
The kernel of every rank-four element of $\mathfrak g_2$ is an associative three-plane \cite{ChemtovKarigiannis2022}.  Writing $P_Z=\ker Z$, the proposition can equivalently be viewed as a complete parametrization of those associative kernel planes arising from $\mathfrak k_A$ that satisfy $P_Z\cap S(P_Z)\ne0$.  The invariant-context mechanism outside this locus is classified in \Cref{thm:rankfour-context-complete}, and \Cref{thm:low-complete} treats all low-dimensional closures.
\end{remark}

\subsection{Quaternionic-stabilizer failures}
We first classify all $S$-invariant quaternionic subalgebras.  Only a lower-dimensional part of this four-dimensional family meets the fixed local algebra $\mathfrak k_A$ in generators that produce proper closures of dimension greater than three.

The following description is the complex-line form of the invariant-plane family described by Draper and Mart\'in-Gonz\'alez \cite{DraperMartin2025}.  Its short proof fixes the parametrization used in the pulse classification.

Let $P_{H_S}$ and $Q_{H_S}=H_S^\perp$ be the $+1$ and $-1$ eigenspaces of $S$ in $V$.  

\begin{lemma}[Invariant associative planes]
\label{thm:invariant-contexts}
An associative three-plane $R\subset V$ is $S$-invariant if and only if either
\[
R=P_{H_S},
\]
or
\[
R=\RR m\oplus\Pi,
\]
where $0\ne m\in P_{H_S}$ and $\Pi\subset Q_{H_S}$ is a complex line for the complex structure $C_m/\lVert m\rVert$.  Hence the nontrivial family is a $\CC P^1$-bundle over $\RR P^2$ and has real dimension four.
\end{lemma}

\begin{proof}
Write $P=P_{H_S}$ and $Q=H_S^\perp$.  If an associative three-plane $R$ is $S$-invariant, then
\[
R=(R\cap P)\oplus(R\cap Q).
\]
Let the two dimensions be $(r_+,r_-)$, so $r_++r_-=3$.  The case $(3,0)$ gives $R=P$.  The case $(0,3)$ is impossible: if $x,y\in R\subset Q$ are independent, multiplicativity of $S$ gives
\[
S(x\times y)=Sx\times Sy=(-x)\times(-y)=x\times y.
\]
Thus the nonzero vector $x\times y$ lies in $P$, whereas associativity of $R$ requires it to lie in $R\subset Q$.  The case $(2,1)$ is also impossible: two independent vectors of the associative plane $P$ have cross product in $P$ and generate its third quaternionic direction, so an associative plane containing them cannot have only a two-dimensional intersection with $P$.

It remains to consider $(1,2)$.  Write $R\cap P=\RR m$ and $\Pi=R\cap Q$.  For $x\in\Pi$, associativity of $R$ implies $m\times x\in\Pi$.  Since $C_m^2=-\lVert m\rVert^2 I$ on $Q$, the operator $C_m/\lVert m\rVert$ is a complex structure on $Q$, and $\Pi$ is a complex line.  Conversely, if $\Pi$ is such a complex line, choose a unit $x\in\Pi$ and set $y=(m/\lVert m\rVert)\times x$.  Then $(m/\lVert m\rVert,x,y)$ is an oriented associative triple, so $\RR m\oplus\Pi$ is associative.  The positive line is parametrized by $\RR P^2$ and, for fixed $m$, the complex lines in the complex two-space $(Q,C_m/\lVert m\rVert)$ form $\CC P^1$.
\end{proof}

\begin{lemma}[Intersections of distinct context stabilizers]
\label{lem:nonfixed-context-bound}
For any two distinct quaternionic contexts $H,H'$, one has
\[
\dim\bigl(\mathfrak k_H\cap\mathfrak k_{H'}\bigr)\le3.
\]
The bound is sharp.
\end{lemma}

\begin{proof}
Set $\mathfrak h=\mathfrak k_H\cap\mathfrak k_{H'}$ and identify $\mathfrak k_H=\mathfrak s_1\oplus\mathfrak s_2$ with $\mathfrak s_i\cong\mathfrak{su}(2)$.  Let $p_i$ be the two coordinate projections and put $\mathfrak n_1=\mathfrak h\cap(\mathfrak s_1\oplus0)$ and $\mathfrak n_2=\mathfrak h\cap(0\oplus\mathfrak s_2)$.  The spaces $p_i(\mathfrak h)$ are subalgebras of $\mathfrak{su}(2)$, so their dimensions are $0$, $1$, or $3$, and $\mathfrak n_i$ is an ideal in $p_i(\mathfrak h)$.  Moreover, pairing the two coordinates of elements of $\mathfrak h$ induces an isomorphism
\[
p_1(\mathfrak h)/\mathfrak n_1\simeq p_2(\mathfrak h)/\mathfrak n_2.
\]
If $\dim\mathfrak h\ge4$, this relation and the simplicity of $\mathfrak{su}(2)$ force one of the $\mathfrak n_i$ to be a full three-dimensional factor.  Indeed, if both projections are onto and both kernels vanish, $\mathfrak h$ is the graph of an isomorphism and has dimension three; all remaining possibilities of dimension at least four contain a full factor.  Thus $\mathfrak h$ contains one of the two simple factors of $\mathfrak k_H$.

The branching in \Cref{lem:maximality} recovers $P_H$ from either factor: the second factor fixes $P_H$ pointwise, while the first has $P_H$ as its unique three-dimensional invariant summand.  In either case $P_H$ is the only invariant three-dimensional subspace for that fixed factor: for the pointwise factor the other real summand is irreducible of dimension four, and for the first factor the summands of dimensions three and four are irreducible and inequivalent.  Since $P_{H'}$ is also invariant under the factor contained in $\mathfrak h$, it must equal $P_H$, forcing $H'=H$.  The contrapositive gives $\dim\mathfrak h\le3$.  For sharpness, take
\[
H=\exp(\pi D_{25}/12)H_A,
\]
whose associative frame may be taken as
\[
\frac{\sqrt3e_1+e_6}{2},\qquad
\frac{e_2+\sqrt3e_5}{2},\qquad
\frac{\sqrt3e_3-e_4}{2}.
\]
Direct substitution gives $S(H)\ne H$ and
\[
\mathfrak k_H\cap\mathfrak k_{S(H)}
=\Span\left\{
\tfrac12D_{17}+D_{24},\,-D_{15}+D_{26},\,D_{13}+D_{46}
\right\},
\]
which has dimension three.
\end{proof}

Write $q_a=ae_1+ce_2+se_3$.  For parameters with $n=c^2+s^2>0$, set
\[
\zeta=c+is,\qquad r=\operatorname{Re}(\zeta^3)=c^3-3cs^2,
\qquad t=\operatorname{Im}(\zeta^3)=3c^2s-s^3.
\]
The projective family and its two candidate intersection generators are
\begin{equation}
H_{[c:s]}=\RR1\oplus\Span\{ce_2+se_3,ce_4+se_5,(c^2-s^2)e_6+2cs\,e_7\},
\qquad[c:s]\in\RR P^1,
\label{eq:projective-context-family}
\end{equation}
\begin{align}
\mathsf P_{c,s}&=snD_{12}+rD_{46}+tD_{47},\notag\\
\mathsf Q_{c,s}&=cnD_{13}+rD_{46}+tD_{47}.
\label{eq:projective-intersection-basis}
\end{align}
For the resonant family, define
\begin{align*}
u_{a,c,s}&=2ae_1+c(e_2+e_4)+s(e_3+e_5),\\
v_{c,s}&=c(e_2-e_4)+s(e_3-e_5),\\
h_{a,c,s}&=-as(e_2-e_4)+ac(e_3-e_5)+(s^2-c^2)e_6-2cs\,e_7,
\end{align*}
\begin{equation}
H_{a;c,s}=\RR1\oplus\Span\{u_{a,c,s},v_{c,s},h_{a,c,s}\}.
\label{eq:resonant-context-family}
\end{equation}
The homogeneous pulse polynomial is
\[
\widehat Z_{a;c,s}
=2c^2sD_{12}-c(c^2-s^2)D_{13}
 +an(D_{23}+D_{45})+rD_{46}+tD_{47}.
\]
These subspaces are shown below to be quaternionic subalgebras.  To state their
full-generation certificates, use the following six-word pattern for an element $U$:
\[
F_1(U)=U,\quad F_2(U)=\Ad_SU,\quad F_3(U)=[F_1(U),F_2(U)],
\]
\[
F_4(U)=[F_1(U),F_3(U)],\quad
F_5(U)=[F_2(U),F_3(U)],\quad
F_6(U)=[F_1(U),F_4(U)].
\]
Arguments will be suppressed when the pulse is specified.

\begin{theorem}[Rank-six quaternionic-stabilizer failures]
\label{thm:visible-contexts}
Let $Z\in\mathfrak k_A$ have rank six and suppose
$L_Z\subseteq\mathfrak k_H$ with $\dim L_Z>3$ for a quaternionic context $H$.
Then $S(H)=H$.  If $Z\notin\mathcal N_{e_1}$, so that $n>0$, precisely the following cases occur.
\begin{enumerate}[label=\textnormal{(\roman*)},leftmargin=9mm]
\item If $a=0$, then $H=H_{[c:s]}$ and
\[
\mathfrak k_A\cap\mathfrak k_{H_{[c:s]}}
 =\Span\{\mathsf P_{c,s},\mathsf Q_{c,s}\}.
\]
\item If $a\ne0$, then $H=H_{a;c,s}$ and the second quaternionic coordinate is forced to be
\begin{equation}
q_b=3\left(ae_1+\frac r n e_2+\frac t n e_3\right).
\label{eq:resonant-qb}
\end{equation}
Equivalently,
\begin{equation}
nZ=\widehat Z_{a;c,s}.
\label{eq:resonant-pulse}
\end{equation}
The pulse $\widehat Z_{a;c,s}$ satisfies $\beta=3\alpha$ and generates the full stabilizer:
\[
\Lie\{\widehat Z_{a;c,s},\Ad_S\widehat Z_{a;c,s}\}
 =\mathfrak k_{H_{a;c,s}}\cong\mathfrak{so}(4).
\]
Its six-word trace-Gram certificate is
\begin{equation}
\det\bigl(-\operatorname{tr}(F_i(\widehat Z)F_j(\widehat Z))\bigr)_{i,j=1}^{6}
 =2^{46}3^{11}n^{34}(a^2+n)^2(2a^2+n)^6\ne0.
\label{eq:resonant-so4-gram}
\end{equation}
\end{enumerate}
Conversely, every pulse in the projective intersection of (i) and every homogeneous pulse
$\widehat Z_{a;c,s}$ preserves the stated $S$-invariant context.
The intersection formula is valid at every projective point, and the homogeneous certificate
also holds at $a=0$.
\end{theorem}

\begin{proof}
Because $L_Z$ is $\Ad_S$-stable,
\[
L_Z\subseteq\mathfrak k_H\cap\mathfrak k_{S(H)}.
\]
The hypothesis $\dim L_Z>3$ and \Cref{lem:nonfixed-context-bound} force $S(H)=H$.  The possible invariant associative planes are therefore exactly those described in \Cref{thm:invariant-contexts}.  Since $Z$ has rank six, its kernel is the line $\RR q_a$.  The restriction of a skew operator to the invariant odd-dimensional plane $P_H$ has a zero eigenvalue, so $q_a\in P_H$.  Hence $S(q_a)\in P_H$ as well.  Outside $\mathcal N_{e_1}$ these two vectors are independent and, by alternativity, determine the unique quaternionic context generated by them.  This argument uses no spectral-separation hypothesis.

If $a=0$, the two vectors are $ce_2+se_3$ and $ce_4+se_5$; their cross product is $(c^2-s^2)e_6+2cs\,e_7$, giving \eqref{eq:projective-context-family}.  Imposing invariance on a general element of $\mathfrak k_A$ gives the $12\times6$ system printed in Appendix~\ref{app:projective-failure-certificates}; its row reduction there proves rank four for every $n>0$.  The displayed vectors solve it.  In the invariant trace metric,
\begin{align*}
\langle\mathsf P,\mathsf P\rangle_{\mathrm{tr}}&=48n(c^4+3s^4),\\
\langle\mathsf Q,\mathsf Q\rangle_{\mathrm{tr}}&=48n(3c^4+s^4),\\
\langle\mathsf P,\mathsf Q\rangle_{\mathrm{tr}}&=24n(3c^4-2c^2s^2+3s^4),
\end{align*}
and
\[
\det\operatorname{Gram}_{\mathrm{tr}}(\mathsf P,\mathsf Q)
=2^6 3^3 n^6>0.
\]
They therefore span the full kernel of the constraint system.  The polynomial $n^4(c^4+3c^2s^2+s^4)$ is only the coordinate Gram determinant in the nonorthonormal displayed derivation basis.

Now suppose $a\ne0$.  The sum and difference of $q_a$ and $S(q_a)$ are $u_{a,c,s}$ and $v_{c,s}$.  Direct multiplication gives
\[
u_{a,c,s}\times v_{c,s}=2h_{a,c,s},
\quad
v_{c,s}\times h_{a,c,s}=n\,u_{a,c,s},
\quad
h_{a,c,s}\times u_{a,c,s}=(2a^2+n)v_{c,s}.
\]
Thus \eqref{eq:resonant-context-family} is associative and is fixed by $S$, with $S(u)=u$, $S(v)=-v$, and $S(h)=-h$.

Write $q_b=b_1e_1+b_2e_2+b_3e_3$.  Substitution in the invariance equations for this forced context first yields
\[
b_1=3a,
\]
and then
\[
\begin{pmatrix}c^2-s^2&2cs\\2cs&-(c^2-s^2)\end{pmatrix}
\begin{pmatrix}b_2\\b_3\end{pmatrix}
=
\begin{pmatrix}3cn\\-3sn\end{pmatrix}.
\]
The determinant is $-n^2$, so the solution is unique and equals \eqref{eq:resonant-qb}.  Inverting the linear $(q_a,q_b)$ coordinate map gives \eqref{eq:resonant-pulse}.  In particular,
\[
\alpha^2=a^2+n,
\qquad
\beta^2=9(a^2+n),
\]
for $Z=\widehat Z/n$, and hence $\beta=3\alpha$.

For the converse, the polynomial representative satisfies
\[
\widehat Z u=-2nh,
\qquad
\widehat Z v=2nh,
\qquad
\widehat Z h=n^2u-n(2a^2+n)v.
\]
It therefore preserves $H_{a;c,s}$, and so does its $S$-conjugate.  The normalized Schur-pivot calculation and equivariance in Appendix~\ref{app:projective-failure-certificates} give \eqref{eq:resonant-so4-gram}.  The six words are independent and lie in the six-dimensional stabilizer $\mathfrak k_{H_{a;c,s}}$, proving equality.  At $a=0$, the context reduces to $H_{[c:s]}$ and
\[
\widehat Z_{0;c,s}
=\frac{2c^2}{n}\mathsf P_{c,s}
 -\frac{c^2-s^2}{n}\mathsf Q_{c,s},
\]
so the resonant cone attaches to the projective component along a distinguished pulse line.
\end{proof}

The projective endpoints $[1:0]$ and $[0:1]$ give $H_{246}$ and $H_{365}$;
write $H_\star=H_{[1:1]}$ for the noncoordinate member.  The remaining coordinate
fixed context satisfies
\[
\mathfrak k_A\cap\mathfrak k_{176}=\Span\{D_{23},D_{45}\}\subset\mathcal N_{e_1}.
\]
Individual intersection generators may have rank below six.  For example,
$\mathsf P_{c,s}$ has $\alpha=0$ at $[c:s]=[\sqrt3:1]$, and
$\mathsf Q_{c,s}$ has $\alpha=0$ at $[c:s]=[1:\sqrt3]$.
These are points of $\mathcal N_{e_1}$, not exceptions to the two-dimensional intersection formula.



\begin{corollary}[Resonant rank-six stabilizer cone]
\label{cor:resonant-cone}
The homogeneous family $\widehat Z_{a;c,s}$, with $n>0$, has real affine dimension
three and projective dimension two.  Its members with $a\ne0$ lie outside
$x_2=x_3=0$ and generate full quaternionic stabilizers on the rank-six resonance
$\beta=3\alpha$.
\end{corollary}
\begin{proof}
Every coefficient of $\widehat Z_{a;c,s}$ is homogeneous of degree three in $(a,c,s)$.
The last two coefficients obey $x_4+ix_5=(c+is)^3$, so a nonzero image has only
finitely many choices of $(c,s)$.  Once they are fixed, $x_2=x_3=an$ determines $a$.
Thus the three-dimensional parameter domain has finite fibers and its image has
real dimension three.  Nonzero real scaling acts by scaling all parameters, so
projectivization has dimension two.  If $a\ne0$ then $x_2=x_3\ne0$; the generation
and resonance claims follow from \Cref{thm:visible-contexts}.
\end{proof}


Write
\[
Z=x_0D_{12}+x_1D_{13}+x_2D_{23}+x_3D_{45}+x_4D_{46}+x_5D_{47}
\]
For complex scalars, a superscript $\ast$ denotes complex conjugation. Introduce
\[
w=x_0+ix_1,
\qquad z=x_4+ix_5,
\qquad z^\ast=x_4-ix_5,
\qquad \omega=z-2iw.
\]
Here $z^\ast$ denotes the complex conjugate of $z$.



\begin{corollary}[Projective quartic stabilizer cone]
\label{cor:quartic-cone}
$\mathcal N_{e_1}=\{\omega=0\}$.  In the four-plane $x_2=x_3=0$, the union of the projective context-intersection planes is exactly
\begin{equation}
\operatorname{Im}(z^\ast\omega^3)=0.
\label{eq:so4-quartic-cone}
\end{equation}
It is a real homogeneous quartic cone of dimension three.  Its singular locus is $\omega=0$, a real two-plane in the ambient four-plane.
\end{corollary}

\begin{proof}
For the basis in \eqref{eq:projective-intersection-basis}, one has
\[
z=\mu\zeta^3,
\qquad \omega=\lambda\zeta
\]
with real $\lambda,\mu$, so \eqref{eq:so4-quartic-cone} follows.  Conversely, if $\omega\ne0$, choose a projective representative $\zeta$ in its real direction.  The quartic equation forces $z$ to have the phase of $\zeta^3$ up to sign.  For a pulse $\xi\mathsf P_{c,s}+\eta\mathsf Q_{c,s}$, its two real coefficients satisfy
\[
\begin{pmatrix}\lambda\\\mu\end{pmatrix}
=\begin{pmatrix}c^2-3s^2&3c^2-s^2\\1&1\end{pmatrix}
\begin{pmatrix}\xi\\\eta\end{pmatrix}.
\]
The coefficient matrix has determinant $-2(c^2+s^2)\ne0$, so the required
combination exists.  When $\omega=0$ and $z\ne0$, choose $\zeta\ne0$ with the phase of a cube root of $z$, and take $\lambda=0$ and real $\mu$ satisfying $z=\mu\zeta^3$.  The same invertible coefficient map supplies the required linear combination.  If $z=\omega=0$, the element is zero.  This proves the converse also on the singular locus.

For the singular locus put $F=\operatorname{Im}(z^\ast\omega^3)$.  Direct differentiation in the original real coordinates gives
\[
\frac{\partial F}{\partial x_0}=-6\operatorname{Re}(z^\ast\omega^2),
\qquad
\frac{\partial F}{\partial x_1}=6\operatorname{Im}(z^\ast\omega^2).
\]
Thus a singular point satisfies $z^\ast\omega^2=0$.  If $\omega=0$, all first derivatives vanish.  If instead $z=0$ and $\omega\ne0$, then
\[
\left.\frac{\partial F}{\partial x_4}\right|_{z=0}=\operatorname{Im}(\omega^3),
\qquad
\left.\frac{\partial F}{\partial x_5}\right|_{z=0}=-\operatorname{Re}(\omega^3),
\]
which cannot both vanish.  Hence the singular locus is exactly $\omega=0$.
\end{proof}

Choose a nonzero representative $(c,s)$ of $[c:s]\in\RR P^1$, and set
\[
Z_{[c:s]}=sD_{12}-cD_{13}.
\]
The pulse expression depends on the representative, but its generated algebra does not.
We also define a spectrally regular choice
\begin{equation}
Z^{\mathrm{reg}}_{[c:s]}
 =(3c^2-5s^2)\mathsf P_{c,s}+(c^2+9s^2)\mathsf Q_{c,s}.
\label{eq:uniform-regular-pulse}
\end{equation}
Below $F_i=F_i(Z_{[c:s]})$ and $G_i=F_i(Z^{\mathrm{reg}}_{[c:s]})$
use the six-word pattern preceding \Cref{thm:visible-contexts}.

\begin{proposition}[Uniform generation of the projective stabilizers]
\label{prop:uniform-so4}
\[
Z_{[c:s]}\in\mathfrak k_A\cap\mathfrak k_{H_{[c:s]}},
\qquad
\Ad_SZ_{[c:s]}=sD_{14}-cD_{15},
\]
and
\[
\Lie\{Z_{[c:s]},\Ad_SZ_{[c:s]}\}=\mathfrak k_{H_{[c:s]}}\cong\mathfrak{so}(4).
\]
The closure ranks are uniformly
\[
2\longrightarrow3\longrightarrow5\longrightarrow6\longrightarrow6,
\]
and the invariant trace-Gram determinant is
\begin{equation}
\det\bigl(-\operatorname{tr}(F_iF_j)\bigr)_{i,j=1}^6
=2^{46}3^{11}(c^2+s^2)^{14}\ne0.
\label{eq:uniform-so4-gram}
\end{equation}
The pulse $Z^{\mathrm{reg}}_{[c:s]}$ has stabilizer parameters
\[
\alpha=6(c^2+s^2)^{5/2},
\qquad
\beta=12(c^2+s^2)^{5/2}=2\alpha,
\]
so the pulse is rank six with three pairwise distinct positive frequencies for every projective point.  For the same six bracket pattern its trace-Gram determinant is
\begin{equation}
\det\bigl(-\operatorname{tr}(G_iG_j)\bigr)_{i,j=1}^6
=2^{46}3^{39}(c^2+s^2)^{70}\ne0,
\label{eq:uniform-regular-gram}
\end{equation}
and hence it also generates the full $\mathfrak k_{H_{[c:s]}}$ uniformly.
\end{proposition}

\begin{proof}
The identity
\[
\mathsf P_{c,s}-\mathsf Q_{c,s}=(c^2+s^2)Z_{[c:s]}
\]
places the simple pulse in the stated intersection.  The switch-conjugate formula follows from the exact action of $S$ on the derivation basis.  The principal trace-Gram minors printed in Appendix~\ref{app:projective-failure-certificates}, of orders $2,3,5,6$, are nonzero for $n>0$ and give the rank ladder and \eqref{eq:uniform-so4-gram}.  The six words lie in the six-dimensional stabilizer, so this certificate proves equality of the generated algebra with that stabilizer.

More generally, for $U=u\mathsf P_{c,s}+v\mathsf Q_{c,s}$ one finds
\[
\alpha=|\lambda|\sqrt n,
\qquad
\beta=3|\mu|n^{3/2},
\]
where $n=c^2+s^2$, $\mu=u+v$, and
\[
\lambda=u(c^2-3s^2)+v(3c^2-s^2).
\]
For the coefficients in \eqref{eq:uniform-regular-pulse}, $\mu=4n$ and $\lambda=6n^2$, giving $\beta=2\alpha$.  The regular-pulse Schur pivots and homogeneity in Appendix~\ref{app:projective-failure-certificates} give \eqref{eq:uniform-regular-gram}.
\end{proof}

The regular pulse in \eqref{eq:uniform-regular-pulse} has $\beta=2\alpha$, so its three frequencies have ratio $1:2:3$.  That numerical spectrum alone does not identify a containing principal algebra: \Cref{prop:uniform-so4} proves that this pulse generates the full quaternionic stabilizer with the fixed switch.  The principal-exclusion statement in \Cref{cor:no-principal} is about the generated algebra, not about the spectrum of a single pulse.  For varying reflections, principal closures do occur, as shown in \Cref{prop:principal-moving}.

\subsection{Maximal-type reduction and the low-dimensional branch}
The compact classification recalled in \Cref{subsec:compact-maximals} gives the following reduction for the fixed reflection.

\begin{proposition}[Maximal-container reduction]
\label{prop:maximal-reduction}
Let
\[
L_Z=\Lie\{Z,\Ad_SZ\}\subsetneq\gtwo.
\]
Then one of the following holds:
\begin{enumerate}[label=\textnormal{(\roman*)},leftmargin=9mm]
\item $L_Z\subseteq\gtwo^w\cong\mathfrak{su}(3)$ for an $S$-eigenvector $w$;
\item $L_Z\subseteq\mathfrak k_H\cong\mathfrak{so}(4)$ for an $S$-invariant quaternionic context $H$;
\item $\dim L_Z\le3$.
\end{enumerate}
Consequently the maximal-container type classification of proper constrained closures with dimension greater than three is complete: they are of vector or invariant-context type.  The intersection bound is sharp, as shown in \Cref{lem:nonfixed-context-bound}.
\end{proposition}

\begin{proof}
By the compact maximal-subalgebra classification, a proper closure is contained in a maximal subalgebra of vector-stabilizer, context-stabilizer, or principal-$\mathfrak{su}(2)$ type \cite{Dynkin1952b,DraperMartin2025}.  Because $L_Z$ is $\Ad_S$-stable, vector containment in $\gtwo^w$ also gives containment in $\gtwo^{Sw}$.  If $Sw\not\parallel w$, their intersection is the stabilizer of two independent vectors and has dimension three; otherwise (i) holds.

Likewise, containment in $\mathfrak k_H$ gives containment in $\mathfrak k_H\cap\mathfrak k_{S(H)}$.  If $S(H)\ne H$, \Cref{lem:nonfixed-context-bound} gives dimension at most three; otherwise (ii) holds.  A principal $\mathfrak{su}(2)$ is itself three-dimensional, so it belongs to case (iii).
\end{proof}

\subsection{The complete low-dimensional locus}
The maximal-container reduction leaves a finite-dimensional question which can be solved without a large projective elimination.  A circle of symmetries preserves both the local control domain and the fixed switch.  It removes one angular variable and turns the three-dimensional closure condition into two elementary real normal forms.

Write the local coordinates as
\begin{equation}
 q_a=(a_1,a_2,a_3)=(2x_2-x_3,-2x_1-x_4,2x_0-x_5),\qquad
 q_b=(b_1,b_2,b_3)=3(x_3,x_4,x_5),
 \label{eq:terminal-qa-qb}
\end{equation}
and set $q=a_2+ia_3$, $p=b_2+ib_3$.  In particular $q=-\omega$ and $p=3z$ in the notation of \Cref{cor:quartic-cone}.  To avoid confusion, $p$ and $q$ here are complex coordinates, not quaternionic contexts.  The inverse of \eqref{eq:terminal-qa-qb} is
\[
(x_0,x_1,x_2,x_3,x_4,x_5)
=\left(\frac{a_3}{2}+\frac{b_3}{6},
-\frac{a_2}{2}-\frac{b_2}{6},\frac{a_1}{2}+\frac{b_1}{6},
\frac{b_1}{3},\frac{b_2}{3},\frac{b_3}{3}\right).
\]
  Let
\[
 \Xi=\frac12(D_{23}+D_{45}),\qquad g_\theta=\exp(\theta\Xi).
\]
The transformation $\Ad_{g_\theta}$ preserves $\mathfrak k_A$ and commutes with $\Ad_S$.  Direct differentiation of \eqref{eq:terminal-qa-qb} gives
\begin{equation}
 (a_1,b_1,q,p)\longmapsto(a_1,b_1,e^{i\theta}q,e^{3i\theta}p).
 \label{eq:residual-circle}
\end{equation}
This is a symmetry of the constrained closure problem, not an additional available control.

The weights one and three mean exactly the factors $e^{i\theta}$ and $e^{3i\theta}$ in \eqref{eq:residual-circle}.  A normal form below is taken up to this circle action and nonzero real rescaling of $Z$; these operations conjugate the generated algebra or leave it unchanged.  Scalar coordinates $a,b,c,d$ in a normalized pair such as $q_a=(a,1,0)$, $q_b=(d,b,c)$ are local to that normal-form calculation.

\begin{theorem}[Complete classification below dimension four]
\label{thm:low-complete}
Let $Z\in\mathfrak k_A$ and $L_Z=\Lie\{Z,\Ad_SZ\}$.
\begin{enumerate}[label=\textnormal{(\roman*)},leftmargin=9mm]
\item The closure is abelian if and only if
\[
 x_0=x_1=x_4=x_5=0.
\]
On this plane $Z=x_2D_{23}+x_3D_{45}$.  Its closure has dimension zero at the origin, dimension one on the two nonzero lines $x_2=x_3$ and $x_2=-x_3$, and dimension two elsewhere.
\item Outside this plane, $\dim L_Z\le3$ if and only if $a_1=b_1=0$ and either
\[
 q=0,\quad p\ne0,
\]
or
\begin{equation}
 q\ne0,\qquad |q|^2p=\eta q^3,
 \qquad \eta\in\{1,-1,-3\}.
 \label{eq:three-small-families}
\end{equation}
Every closure in (ii) is isomorphic to $\mathfrak{su}(2)$ and has dimension exactly three.
\end{enumerate}
The dimensions of the common fixed spaces in $V$ are, respectively, one for $q=0$, and three, one, zero for $\eta=1,-1,-3$.
\end{theorem}
\begin{proof}
Put $A=(Z+\Ad_SZ)/2$, $B=(Z-\Ad_SZ)/2$, and $W=[A,B]$.  Substitution of the matrices of $A$ and $B$ gives
\[
 W_{24}+W_{35}=12(x_0^2-x_0x_5+x_5^2+x_1^2+x_1x_4+x_4^2).
\]
The quadratic form is positive definite in $(x_0,x_1,x_4,x_5)$.  Thus $W=0$ forces those four coordinates to vanish; conversely $D_{23}$ and $D_{45}$ commute and are exchanged by $\Ad_S$.  This proves (i).

If $W\ne0$, the three vectors $A,B,W$ are nonzero and pairwise trace-orthogonal.  A three-dimensional closure must therefore equal their span.  Parity under $\Ad_S$ and invariance of the trace metric show that this is equivalent to
\[
[A,W]=-kB,\qquad [B,W]=lA
\]
for real $k,l>0$.  Conversely these relations close the span and give a compact simple three-dimensional algebra.

When $q\ne0$, use \eqref{eq:residual-circle} and a nonzero scalar rescaling to normalize
\[
 q_a=(a,1,0),\qquad q_b=(d,b,c).
\]
Using the normalized matrix of $Z$ printed in Appendix~\ref{app:terminal-certificates}, the matrices $U=[A,W]+kB$ and $V_0=[B,W]-lA$ have entries
\[
 U_{12}=3c,\quad (V_0)_{25}=-3\{a(b^2+c^2+2)-d\},\quad
 U_{16}=-2\{3ab+(b+2)d\}.
\]
Their vanishing implies $c=0$, $d=a(b^2+2)$, and
\[
 a(b+1)(b^2+b+4)=0.
\]
The final factor is strictly positive for real $b$.  If $b=-1$, then $d=3a$ and $U_{37}=16a^2$, so $a=d=0$ in all cases.  The remaining equations give
\[
 k=l=\frac{3b^2+6b+7}{2},\qquad (b-1)(b+1)(b+3)=0.
\]
Substitution verifies sufficiency for all three roots.  Undoing the normalization gives \eqref{eq:three-small-families}.

When $q=0$ but $p\ne0$, normalize instead to $q_a=(a,0,0)$, $q_b=(d,1,0)$.  Now $(V_0)_{25}=-3a$ forces $a=0$, and the remaining equations imply $d(d^2+1)=0$, hence $d=0$.  Substitution again gives a nonabelian three-dimensional closure.  The stated fixed-space dimensions follow from the kernels of the two normalized generators; their trace norms and Casimir spectra are recorded in Appendix~\ref{app:terminal-certificates}.  All these reductions use real equations; no assertion about a complex projective solution set is needed.
\end{proof}

\begin{corollary}[Exclusion of the principal component]
\label{cor:no-principal}
No nonzero $Z\in\mathfrak k_A$ has $L_Z$ contained in a principal $\mathfrak{su}(2)$ of $\mathfrak g_2$.
\end{corollary}
\begin{proof}
A principal $\mathfrak{su}(2)$ acts irreducibly on $V$ and every nonzero element has three distinct positive frequencies in the ratio $1:2:3$.  The nonabelian cases of \Cref{thm:low-complete} instead have rank four, or rank six with frequencies $1:1:2$.  The nonzero one-dimensional cases in (i) have the same alternatives.  An abelian two-dimensional algebra cannot lie in $\mathfrak{su}(2)$, and a closure of dimension greater than three cannot lie there either.
\end{proof}

\subsection{Rank-four quaternionic-stabilizer failures}
The rank-four common-vector locus was already parametrized in \Cref{prop:rankfour-complete}.  Outside that locus, the invariant-context mechanism is nonempty but considerably more rigid than a general four-parameter family of invariant contexts.  A normalized pulse determines its candidate context from the positive part of the symmetric decomposition, even though its own kernel is three-dimensional.

For $0<\tau<1$, put
\begin{equation}
 D_\tau=\tau^2+3,\qquad
 \kappa^2=\frac{(1-\tau)D_\tau}{\tau},\qquad
 h=-\frac{\tau(\tau+1)}{D_\tau}\kappa,
 \label{eq:rankfour-parameter}
\end{equation}
where either real sign of $\kappa$ is allowed.  The sign is not an additional discrete modulus: modulo the residual circle and nonzero rescaling,
\[
E_{\tau,-\kappa}=-\Ad_{g_\pi}E_{\tau,\kappa}.
\]
Define
\begin{align}
 E_{\tau,\kappa}={}&(\tau-3)D_\tau D_{13}
 +(\tau^2-\tau+6)\kappa D_{23}\notag\\
 &+(2\tau^2+\tau+3)\kappa D_{45}
 -2\tau D_\tau D_{46},
 \label{eq:rankfour-new-pulse}\\
 m={}&\frac\kappa2e_1+e_2+e_4,\qquad
 v=h(e_3-e_5)-(\tau+1)e_6,\qquad r=m\times v.
 \label{eq:rankfour-new-frame}
\end{align}
Set $H_{\tau,\kappa}=\RR1\oplus\Span\{m,v,r\}$.
The notation $E_{\tau,\kappa}$ denotes a derivation, not a matrix unit.

The symbol $D_\tau$ is a positive scalar polynomial, unlike the derivations $D_{ij}$.  The parameters $\tau,\kappa$ and the frame $m,v,r$ in this subsection do not redefine the scalar coordinates or the frame parameters used in the rank-six family.

\begin{theorem}[Complete rank-four invariant-context mechanism]
\label{thm:rankfour-context-complete}
Every $E_{\tau,\kappa}$ has rank four, no common kernel with its switch conjugate, and
\[
 \Lie\{E_{\tau,\kappa},\Ad_SE_{\tau,\kappa}\}
 =\mathfrak k_{H_{\tau,\kappa}}\cong\mathfrak{so}(4).
\]
Conversely, if $Z\in\mathfrak k_A$ has rank four, $\ker Z\cap S\ker Z=0$, and $L_Z$ is contained in an $S$-invariant context stabilizer, then
\[
 Z=\lambda\Ad_{g_\theta}E_{\tau,\kappa}
\]
for some $\lambda\ne0$, $\theta\in\RR$, and parameters satisfying \eqref{eq:rankfour-parameter}.  Thus the entire mechanism is a three-real-parameter family after allowing rescaling and the residual circle.
\end{theorem}
\begin{proof}
Outside the common-vector locus, $q\ne0$: otherwise $Ze_1=0$.  Normalize to $q_a=(a,1,0)$ and $q_b=(d,b,c)$ using \eqref{eq:residual-circle}.  Rank four is equivalent to
\begin{equation}
 d^2+b^2+c^2=a^2+1.
 \label{eq:rankfour-sphere}
\end{equation}
Write $A=(Z+\Ad_SZ)/2$ and $B=(Z-\Ad_SZ)/2$.  On the positive eigenspace of $S$, in the ordered frame $(e_1,e_2+e_4,e_3+e_5)$, the even part $A$ has matrix
\[
 \begin{pmatrix}0&0&2\\0&0&-(a+d)/2\\-1&(a+d)/2&0\end{pmatrix}.
\]
The frame is not orthonormal: its Gram matrix is $\operatorname{diag}(1,2,2)$.  The unique kernel line is generated by $m=(a+d)e_1/2+e_2+e_4$.  Every invariant context must have this positive line; the exceptional all-positive context would give a common fixed vector and is excluded.  Its negative plane contains
\[
 v=Bm=(a-d)(e_3-e_5)+(b-1)e_6+ce_7.
\]
This vector is nonzero, since otherwise $m$ would be a common kernel vector.  Associativity forces the entire context to be $\RR1+\Span\{m,v,m\times v\}$.

Preservation by $B$ requires $Bv\parallel m$.  Since $m$ has zero $e_3$ component and $(Bv)_3=c$, its coordinate equations give $c=0$.  With $h=a-d$, $k=a+d$, the remaining equation and \eqref{eq:rankfour-sphere} reduce to
\[
hk=b^2-1,\qquad (b^2+3)h+b(b-1)k=0.
\]
If $b=1$, these equations force $h=0$, so $v=0$.  The value $b=0$ is impossible.  For $b\ne0,1$ they imply
\[
 k^2=-\frac{(b+1)(b^2+3)}b.
\]
Real solutions have $-1\le b<0$.  At $b=-1$, $k=h=0$, and \Cref{thm:low-complete} gives a common-vector closure.  The remaining interval is exactly $b=-\tau$, $0<\tau<1$.  Solving for $a,d$ yields
\[
 a=\frac{3-\tau}{2D_\tau}\kappa,\qquad
 d=\frac{2\tau^2+\tau+3}{2D_\tau}\kappa.
\]
Multiplying the normalized pulse by $6D_\tau$ gives \eqref{eq:rankfour-new-pulse}.  This proves necessity without eliminating any complex variables.

For sufficiency, use the even and odd parts $A_E,B_E$ of $E_{\tau,\kappa}$.  Exact substitution into \eqref{eq:rankfour-new-frame} gives
\[
 A_Em=0,\quad B_Em=6D_\tau v,\quad
 A_Ev=-\frac{12D_\tau}{\tau+1}r,\quad
 B_Ev=-24\tau(\tau+1)m.
\]
The derivation rule then supplies the action on $r=m\times v$.  Hence both generators preserve the same associative plane.  They have rank four by \eqref{eq:rankfour-sphere}.  To prove full six-dimensional closure, take
\[
 F=(A_E,B_E,[A_E,B_E],[A_E,[A_E,B_E]],
 [B_E,[A_E,B_E]],[A_E,[A_E,[A_E,B_E]]]).
\]
Their trace-Gram determinant is
\begin{equation}
 \det\operatorname{Gram}_{\rm tr}(F)=
 \frac{2^{32}3^{31}(3-\tau)^{10}(1-\tau)^6(1+\tau)^{10}
 D_\tau^{24}(\tau^2-2\tau+3)^4}{\tau^6}>0.
 \label{eq:rankfour-new-gram}
\end{equation}
A short block-factorized certificate is given in Appendix~\ref{app:terminal-certificates}.  Thus the closure is the full stabilizer.  Its $3\oplus4$ representation on $V$ has no fixed vector, proving the common-kernel assertion as well.
\end{proof}

For example, at $\tau=1/2$ and $\kappa=\sqrt{13}/2$, an integral-quadratic rescaling gives
\[
Z_{\rm new}=-65D_{13}+23\sqrt{13}D_{23}
             +16\sqrt{13}D_{45}-26D_{46}.
\]
It satisfies $\alpha^2=\beta^2=36036$, has rank four and no common vector with $\Ad_SZ_{\rm new}$, and generates a full $\mathfrak{so}(4)$.  This is an explicit rank-four quaternionic-stabilizer failure without a common fixed vector.
The open interval $0<\tau<1$ is forced by the real non-common-vector branch rather than chosen by convention.  As $\tau\to1^-$ one has $\kappa\to0$ and
\[
E_{\tau,\kappa}\longrightarrow-8(D_{13}+D_{46}),
\]
which is the $\eta=-1$ common-vector representative from the low-dimensional classification.  As $\tau\to0^+$, $|\kappa|\to\infty$ in the normalization \eqref{eq:rankfour-parameter}.

We now assemble the four failure families.  Define
\begin{align*}
\mathcal F_{\rm vec}&=\{Z\in\mathfrak k_A:\ker Z\cap S(\ker Z)\ne\{0\}\},\\
\mathcal F_{\rm proj}&=\bigcup_{[c:s]\in\RR P^1}
 (\mathfrak k_A\cap\mathfrak k_{H_{[c:s]}}),\\
\mathcal F_{\rm res}&=\{\widehat Z_{a;c,s}:a,c,s\in\RR,\ c^2+s^2>0\}\cup\{0\},\\
\mathcal F_{\rm rf}&=\{\lambda\Ad_{g_\theta}E_{\tau,\kappa}:\lambda\in\RR\setminus\{0\},\
\theta\in\RR,\ 0<\tau<1,\
\kappa^2=(1-\tau)(\tau^2+3)/\tau\}.
\end{align*}
The common-vector set is described in \Cref{prop:failure,prop:rankfour-complete},
the projective set in \Cref{cor:quartic-cone}, the resonant set in
\Cref{cor:resonant-cone}, and the last set in \Cref{thm:rankfour-context-complete}.
These families can overlap and are not asserted to be irreducible algebraic components.

\begin{theorem}[Exact nongenerating locus]
\label{thm:complete-failure-set}
For the fixed pair $(\mathfrak k_A,S)$, the complete nongenerating locus is
\begin{equation}
\{Z\in\mathfrak k_A:L_Z\ne\mathfrak g_2\}
 =\mathcal F_{\rm vec}\cup\mathcal F_{\rm proj}\cup\mathcal F_{\rm res}\cup\mathcal F_{\rm rf}.
\label{eq:complete-failure-set}
\end{equation}
\end{theorem}
\begin{proof}
First consider the inclusion from right to left.  The common-vector criterion
\Cref{prop:failure} places $\mathcal F_{\rm vec}$ in a vector stabilizer.
\Cref{thm:visible-contexts} places $\mathcal F_{\rm proj}$ and the nonzero part of
$\mathcal F_{\rm res}$ in quaternionic stabilizers.  Finally,
\Cref{thm:rankfour-context-complete} does the same for $\mathcal F_{\rm rf}$.
All of these containing algebras are proper, and $Z=0$ plainly fails to generate.

For the reverse inclusion, let $L_Z\ne\mathfrak g_2$.  If
$\dim L_Z>3$, \Cref{prop:maximal-reduction} places it in a vector stabilizer or
an $S$-invariant quaternionic stabilizer.  The vector case gives
$Z\in\mathcal F_{\rm vec}$.  In the other case we may assume there is no common
fixed vector; in particular $Z\notin\mathcal N_{e_1}$.  The nonzero element $Z$
has rank four or six.  At rank six, \Cref{thm:visible-contexts} gives
$Z\in\mathcal F_{\rm proj}\cup\mathcal F_{\rm res}$; the rescaling in
\eqref{eq:resonant-pulse} does not change this conical set.
At rank four, \Cref{thm:rankfour-context-complete} gives
$Z\in\mathcal F_{\rm rf}$.

If $\dim L_Z\le3$, apply \Cref{thm:low-complete}.
Its abelian plane and its $q=0$ family fix $e_1$, so they lie in
$\mathcal F_{\rm vec}$.  In the three remaining normal-form families,
$a_1=b_1=0$ and $|q|^2p=\eta q^3$ with $\eta\in\{1,-1,-3\}$.
Thus $x_2=x_3=0$ and $\operatorname{Im}(p(q^\ast)^3)=0$.
Since $q=-\omega$ and $p=3z$,
\[
\operatorname{Im}\bigl(p(q^\ast)^3\bigr)
 =3\operatorname{Im}(z^\ast\omega^3)=0.
\]
\Cref{cor:quartic-cone} therefore places these cases in $\mathcal F_{\rm proj}$.
The two inclusions cover all possibilities.
\end{proof}




For the four-family description, four witnesses already used above separate the mechanisms:
\[
\begin{split}
Z_{\rm v}&=D_{23},\quad Z_{\rm p}=D_{13},\quad
Z_{\rm r}=-D_{13}+D_{23}+D_{45}+D_{46},\\
Z_{\rm f}&=-65D_{13}+23\sqrt{13}\,D_{23}+16\sqrt{13}\,D_{45}-26D_{46}.
\end{split}
\]
\begin{corollary}[Irredundance of the geometric families]\label{cor:geometric-irredundance}
None of the four families in \Cref{thm:complete-failure-set} is contained in the union of the other three.
\end{corollary}
\begin{proof}
The element $Z_{\rm v}$ fixes $e_1$, has rank six, and has zero transverse coordinate $q$ and nonzero axial coordinate $a_1$; it belongs only to $\mathcal F_{\rm vec}$.
For $Z_{\rm p}$, the kernel is $\RR e_2$ and its reflected kernel is $\RR e_4$, so there is no common fixed vector.  It lies in $\mathcal F_{\rm proj}$ by \Cref{prop:uniform-so4}, has rank six, and has $\beta=0$, excluding the other mechanisms.
The element $Z_{\rm r}$ is $\widehat Z_{1;1,0}$: \Cref{thm:visible-contexts} gives full quaternionic-stabilizer closure, $\beta=3\alpha$, and $a_1=1$.  Its rank is six, it has no common fixed vector, and the nonzero axial coordinate excludes the projective family.
Finally, $Z_{\rm f}$ is the rank-four witness following \Cref{thm:rankfour-context-complete}.  It has no common fixed vector and $a_1=30\sqrt{13}\ne0$, so it is neither vector nor projective.  Its $\alpha=\beta>0$ excludes the resonant family.  Thus each displayed witness belongs to its named family alone.
\end{proof}
The irredundance is a statement about these four sets, not a claim that they are disjoint or are irreducible algebraic components.

\section{An explicit polynomial criterion}
\label{sec:scalar}

The geometric families in \Cref{thm:complete-failure-set} can be eliminated without a search over vectors, contexts, or circle parameters.  The resulting criterion involves only the six real coordinates of the local derivation.

Keep the coordinates of \eqref{eq:terminal-qa-qb}, and abbreviate
\[
a=a_1,\qquad d=b_1,\qquad q=a_2+ia_3,\qquad p=b_2+ib_3.
\]
Thus $(a,d,q,p)\in\RR^2\times\CC^2$ are linear coordinates on $\mathfrak k_A$.  Define real polynomials
\begin{align}
u&=|q|^2, & v&=|p|^2, & R+iI&=p(q^\ast)^3,\label{eq:scalar-invariants}\\
\rho&=a^2+d^2+u+v, & \delta&=a^2+u-d^2-v,\nonumber\\
\Theta&=u\bigl((2v+3u)a-3ud\bigr)-(a+d)R.&&&&\label{eq:scalar-theta}
\end{align}
Here $q^\ast$ denotes the complex conjugate of $q$, so the convention in \eqref{eq:scalar-invariants}
conjugates $q$ and not $p$. Equivalently, $R-iI=p^\ast q^3$.
For an explicit expression in real coordinates, write
$q=q_R+iq_I$ and $p=p_R+ip_I$, with all four components real. Then
\begin{align*}
R&=p_R(q_R^3-3q_Rq_I^2)+p_I(3q_R^2q_I-q_I^3),\\
I&=p_I(q_R^3-3q_Rq_I^2)-p_R(3q_R^2q_I-q_I^3).
\end{align*}
Here $I$ denotes this real polynomial, not an identity operator.  Under the residual circle \eqref{eq:residual-circle}, $q$ and $p$ have weights one and three.  Consequently $u,v,R,I,\rho,\delta,\Theta$ are circle-invariant.  Their respective homogeneous degrees are $2,2,4,4,2,2,5$.

The remaining degree-five invariant can likewise be evaluated entirely in real coordinates:
\[
\begin{aligned}
\Theta={}&(q_R^2+q_I^2)\bigl[(2p_R^2+2p_I^2+3q_R^2+3q_I^2)a
                       -3(q_R^2+q_I^2)d\bigr]\\
 &-(a+d)\bigl[p_R(q_R^3-3q_Rq_I^2)+p_I(3q_R^2q_I-q_I^3)\bigr].
\end{aligned}
\]
\begin{lemma}[Common fixed vectors in scalar coordinates]
\label{lem:scalar-vectors}
The pair $(Z,\Ad_S Z)$ has a nonzero common fixed vector if and only if at least one of the following holds:
\begin{alignat}{2}
&u=0; &&\label{eq:vector-scalar-one}\\
&a=0,\qquad d^2+v=u; &&\label{eq:vector-scalar-two}\\
&d=-a,\qquad v=u; &&\label{eq:vector-scalar-three}\\
&d=a,\qquad up=q^3. &&\label{eq:vector-scalar-four}
\end{alignat}
An equality between complex expressions represents two real polynomial equations.
\end{lemma}

\begin{proof}
When $q=0$, $q_a=ae_1$ and $Ze_1=0$, so $e_1$ is a common fixed vector.  Suppose $q\ne0$.  The residual circle and positive rescaling reduce the coordinates to
\[
q_a=(a,1,0),\qquad q_b=(d,b,c).
\]
By \Cref{prop:failure}, it suffices to look for a kernel vector in one of the two eigenspaces of $S$.

A vector in the positive eigenspace is
\[
s e_1+t(e_2+e_4)+r(e_3+e_5).
\]
Its component in $P_A$ must be parallel to $(a,1,0)$.  A nonzero solution therefore has $(s,t,r)=\lambda(a,1,0)$ with $\lambda\ne0$, and its quaternionic normal coefficient is the real scalar $\lambda$.  The equation $q_by=yq_a$ then gives $d=a$, $b=1$, $c=0$.

A negative eigenvector has the form
\[
s(e_2-e_4)+t(e_3-e_5)+r e_6+z e_7.
\]
If its component in $P_A$ is zero, its quaternionic normal coefficient is $y=r e_2+z e_3\ne0$.  Conjugation $q_b=yq_ay^{-1}$ reverses the $e_1$ coordinate and reflects the transverse unit direction.  As $[r:z]$ varies, all transverse unit directions occur.  This gives exactly $d=-a$ and $b^2+c^2=1$.

If the component in $P_A$ is nonzero, parallelism with $(a,1,0)$ forces $a=0$, $t=0$, and $s\ne0$.  Rescale to $s=1$.  The normal coefficient is $y=-1+r e_2+z e_3$, and $q_b=y e_2y^{-1}$ implies $d^2+b^2+c^2=1$.  Conversely, if this equation holds and $d\ne0$, the explicit vector
\[
e_2-e_4+\frac{c}{d}e_6+\frac{1-b}{d}e_7
\]
is annihilated by $Z$.  If $d=0$, the purely normal case just treated supplies a kernel vector.  Thus the negative-eigenvector possibilities are exactly the two stated families.

Undoing the circle action and rescaling sends the positive condition $b=1,c=0$ to $up=q^3$, and the two negative conditions to \eqref{eq:vector-scalar-two} and \eqref{eq:vector-scalar-three}.  This proves both directions.
\end{proof}

\begin{theorem}[Six finite polynomial tests]
\label{thm:scalar-classifier}
For the fixed pair $(\mathfrak k_A,S)$, one has $L_Z\ne\mathfrak g_2$ if and only if at least one of the following six conditions holds:
\begin{alignat}{2}
\mathrm{(i)}\quad &u=0; &&\nonumber\\
\mathrm{(ii)}\quad &a=0,\qquad d^2+v=u; &&\nonumber\\
\mathrm{(iii)}\quad &d=-a,\qquad v=u; &&\nonumber\\
\mathrm{(iv)}\quad &a=d=0,\qquad I=0; &&\label{eq:six-tests}\\
\mathrm{(v)}\quad &d=3a,\qquad up=3q^3; &&\nonumber\\
\mathrm{(vi)}\quad &\delta=0,\qquad I=0,\qquad\Theta=0. &&\nonumber
\end{alignat}
All these equations have degree at most five in the six real coordinates of $Z$.  No circle parameter or quaternionic context has to be chosen to apply the criterion.
\end{theorem}

\begin{proof}
Conditions (i)--(iii) are common-vector failures by \Cref{lem:scalar-vectors}.  Condition (iv) is exactly the projective quartic family: by \eqref{eq:terminal-qa-qb}, $a=d=0$ is equivalent to $x_2=x_3=0$, while $q=-\omega$ and $p=3z$ give
\[
I=\operatorname{Im}\bigl(p(q^\ast)^3\bigr)
 =-3\operatorname{Im}\bigl(z(\omega^\ast)^3\bigr)
 =3\operatorname{Im}(z^\ast\omega^3).
\]
For $q\ne0$, condition (v) is exactly the resonant family \eqref{eq:resonant-qb}: writing $q=c+is$ gives $d=3a$ and $p=3(c+is)^3/(c^2+s^2)$.  Its points with $q=0$ already satisfy (i).

Consider (vi) with $q\ne0$.  Normalize again to $q_a=(a,1,0)$ and $q_b=(d,b,c)$.  Then $I=c$, so $c=0$.  Put $h=a-d$ and $k=a+d$.  The remaining two equations become
\[
hk=b^2-1,\qquad (b^2+3)h+b(b-1)k=0.
\]
If $b=1$, then $h=0$, yielding the positive common-vector family.  The value $b=0$ is impossible.  Otherwise,
\[
k^2=-\frac{(b+1)(b^2+3)}{b}.
\]
Over the reals, this forces $-1\le b<0$.  At $b=-1$, both $h$ and $k$ vanish, so the closure has a common fixed vector.  On $-1<b<0$, set $\tau=-b$.  These are precisely the real normal forms in \Cref{thm:rankfour-context-complete}, which generate quaternionic stabilizers.  Thus (vi) always implies failure; its $q=0$ points are already in (i).

Conversely, the common-vector locus is exhausted by \Cref{lem:scalar-vectors}.  Its last case, $d=a$ and $up=q^3$, belongs to (vi) when $u>0$: it gives $v=u$, $I=0$, $R=u^2$, and hence $\delta=\Theta=0$.  It belongs to (i) when $u=0$.  The projective and resonant families have already been identified with (iv) and (v), up to common-vector points.  Every rank-four invariant-context family satisfies (vi) by its normal form.  The geometric classification \Cref{thm:complete-failure-set} therefore proves exhaustiveness.
\end{proof}

\begin{proposition}[Irredundance of the six scalar conditions]\label{prop:irredundant-tests}
None of the six conditions in \Cref{thm:scalar-classifier} is contained in the union of the other five.
\end{proposition}
\begin{proof}
The following six local coordinate vectors are witnesses.  Direct substitution in \eqref{eq:scalar-invariants}--\eqref{eq:scalar-theta} shows that row $j$ satisfies condition $j$ and no other condition in \eqref{eq:six-tests}.

\begin{center}
\begin{tabular}{cc}\toprule
Condition & $(x_0,x_1,x_2,x_3,x_4,x_5)$\\\midrule
1 & $(0,0,1,0,0,0)$ \\
2 & $(3,0,1,2,0,0)$ \\
3 & $(1,0,1,-1,0,-1)$ \\
4 & $(1,0,0,0,0,0)$ \\
5 & $(0,0,1,1,0,1)$ \\
6 & $(1,0,2,1,0,-1)$ \\
\bottomrule\end{tabular}
\end{center}
Thus omitting any condition would omit a genuine part of the nongenerating set.
\end{proof}

Define six nonnegative real homogeneous polynomials by
\begin{align}
F_1&=u,\nonumber\\
F_2&=a^2\rho+(d^2+v-u)^2,\nonumber\\
F_3&=(a+d)^2\rho+(v-u)^2,\nonumber\\
F_4&=(a^2+d^2)\rho^3+I^2,\label{eq:classifier-factors}\\
F_5&=(d-3a)^2\rho^2+|up-3q^3|^2,\nonumber\\
F_6&=\delta^2\rho^3+I^2\rho+\Theta^2,\nonumber
\end{align}
and put
\begin{equation}
\mathcal P_{34}(Z)=\prod_{j=1}^{6}F_j(Z).
\label{eq:exact-polynomial}
\end{equation}


\begin{corollary}[Homogeneous encoding of the six scalar tests]
\label{thm:polynomial-classifier}
$\mathcal P_{34}$ is circle-invariant and homogeneous of degree $34$, and
\[
\mathcal N=\{\mathcal P_{34}=0\},\qquad
\mathcal G=\{\mathcal P_{34}>0\}.
\]
In particular, the complete nongenerating set is real algebraic, and its complement is a nonempty real Zariski-open set.  The degree is neither claimed to be minimal nor intrinsic to the locus.
\end{corollary}

\begin{proof}
For $Z\ne0$ one has $\rho>0$.  Consequently a factor in
\eqref{eq:classifier-factors} vanishes if and only if all its nonnegative summands vanish.
In order, this gives exactly conditions (i)--(vi) of \Cref{thm:scalar-classifier}.
Thus $\mathcal P_{34}(Z)=0$ if and only if at least one of the six failure tests holds.
At the origin every factor vanishes and the generated algebra is zero, so the same
criterion remains valid there.

The factor degrees are $(2,4,4,8,6,10)$, whose sum is $34$.
The polynomials $u,v,R,I,\rho,\delta,\Theta$ are invariant under the residual circle;
$up-3q^3$ has weight three, so its squared modulus is invariant as well.
This proves homogeneity, nonnegativity, and invariance of the product.
By \Cref{thm:single} and the exact criterion just established,
$\mathcal P_{34}(X)>0$.  Hence its zero set is a proper real algebraic subset
and its complement is nonempty and real Zariski open.
\end{proof}

The factors in \eqref{eq:classifier-factors} express simultaneous real equations by sums of squares with positive powers of $\rho$.  This is the standard real sum-of-squares encoding.  The powers of $\rho$, positive away from the origin, merely equalize homogeneous degrees; other positive homogeneous multipliers give different degrees without changing the zero set.  Such a real equation may have real codimension greater than one.  This construction does not identify the irreducible components of a complex algebraic variety, and it does not extend the criterion from the compact real form to the complex Lie algebra by changing the ground field.

\begin{corollary}[Comparison with the uniform determinant]\label{cor:fixed-R30}
For the fixed local configuration $(\mathfrak k_A,S)$,
\begin{equation}
\mathcal N=\{Z\in\mathfrak k_A:\mathcal R_{30}(S,Z)=0\}
=\{Z\in\mathfrak k_A:\mathcal P_{34}(Z)=0\}.\label{eq:R30-P34}
\end{equation}
\end{corollary}
\begin{proof}
\Cref{cor:no-principal} excludes the irreducible proper outcome for every local pulse with this $S$.
The factorization theorem \Cref{thm:universal-locus} therefore reduces generation to $\mathcal R_{30}(S,Z)>0$.
The scalar encoding in \Cref{thm:polynomial-classifier} has the same zero set.
\end{proof}
The determinant has degree thirty rather than thirty-four.  The six scalar systems are therefore not needed to prove the existence of a decision criterion for this fiber.  Their distinct role is the explicit elimination of auxiliary variables and the finer geometry established in \Cref{sec:failure}; equality of these zero sets is not equality of the polynomials or a comparison of evaluation times.

For the comparison of the sufficient and exact criteria, fix
\[
Z_0=-2D_{12}+2D_{13}+D_{23}+D_{45}+2D_{46}.
\]

\begin{proposition}[The fixed bracket certificate is strictly sufficient]
\label{prop:strict-certificate}
The degree-$58$ determinant in \Cref{thm:generic} satisfies
\[
\{\Delta_{58}\ne0\}\subsetneq\mathcal G,
\qquad \mathcal N\subsetneq\{\Delta_{58}=0\}.
\]
The element $Z_0$ belongs to $\mathcal G\cap\{\Delta_{58}=0\}$.
\end{proposition}

\begin{proof}
Evaluate the words $W_1,\ldots,W_{14}$ from \Cref{thm:single} at $(Z_0,\Ad_S Z_0)$.  Reduction in the basis \eqref{eq:g2basis} gives the exact dependence
\[
-152064W_3+204W_6-24W_7+204W_8+W_{11}=0.
\]
Nevertheless, the ordered columns
\[
W_1,\ldots,W_{10},W_{12},W_{13},W_{14},[W_1,W_{10}]
\]
have determinant
\[
-2^{107}3^{23}7^4\cdot23\cdot37\cdot43^3\cdot79\ne0.
\]
Both equalities are obtained from the same Fano multiplication rule and fixed bracket words as the original certificate.  Hence $\Delta_{58}(Z_0)=0$ but $L_{Z_0}=\mathfrak g_2$.  Independently, substitution into \eqref{eq:classifier-factors} gives
\[
\mathcal P_{34}(Z_0)=180993418927422213570661711872>0.
\]
The nonvanishing bracket determinant proves strictness without using the classification theorem, while the polynomial value confirms consistency with that classification.
\end{proof}

A semialgebraic set is a finite union of sets specified by real polynomial equalities and inequalities.  Its dimension is the maximum dimension of its smooth strata.  Polygonal path connectedness means that any two points can be joined by finitely many line segments contained in the set; the next result gives the stronger bound of two segments.  Its codimension statement is relative to the six-dimensional space $\mathfrak k_A$.

\begin{corollary}[Codimension and connectivity]
\label{cor:dimension-connectivity}
The real algebraic set $\mathcal N\subset\mathfrak k_A\cong\RR^6$ has dimension four.  Its complement $\mathcal G$ has full Lebesgue measure, and any two elements of $\mathcal G$ can be joined by two line segments lying entirely in $\mathcal G$.  In particular, $\mathcal G$ and its intersection with the unit sphere are path connected.
\end{corollary}

\begin{proof}
The common-vector set in \Cref{lem:scalar-vectors} has dimension at most four.  Indeed, its first family is a four-plane.  Its second and third families are respectively defined, after one linear restriction, by a nonzero quadratic equation in five real variables.  In the fourth, $q\ne0$ determines $p=q^3/|q|^2$, leaving the three parameters $(a,q)$; its $q=0$ part is also at most three-dimensional.  The four-plane $u=0$ gives the matching lower bound.

The projective family is a quartic hypersurface in a four-plane and has dimension three by \Cref{cor:quartic-cone}.  The resonant and rank-four families have dimension at most three by their explicit parametrizations.  \Cref{thm:complete-failure-set} therefore gives $\dim_\RR\mathcal N=4$.  These dimension statements use the usual dimension of semialgebraic sets; the parameter descriptions also provide finite smooth stratifications, so a set of dimension at most five in $\RR^6$ has Lebesgue measure zero.

The remaining path-connectivity assertion is a general consequence for complements of closed semialgebraic sets of codimension at least two, applied here to the dimension just computed.  Fix $z\in\mathcal G$.  The intermediate endpoints $w$ for which the segment $[z,w]$ meets $\mathcal N$ are contained in the semialgebraic image
\[
B_z=\{z+t(n-z):n\in\mathcal N,\ t\ge1\}.
\]
This set has dimension at most five: the dimension of a semialgebraic image does not exceed that of its parameter space \cite[\S2.8]{BochnakCosteRoy1998}, here $\dim\mathcal N+1=5$.  Given $z_1,z_2\in\mathcal G$, choose $w\notin B_{z_1}\cup B_{z_2}$.  The two segments $[z_1,w]$ and $[w,z_2]$ avoid $\mathcal N$, as required.  The unit-sphere assertion follows by radial projection of this path: the origin lies in $\mathcal N$, and $\mathcal G$ is invariant under nonzero rescaling by homogeneity of $\mathcal P_{34}$.
\end{proof}

\section{Every context-moving reflection admits a generating pulse}
\label{sec:varying-switch}

The fixed involution admits a complete scalar classification.  A different question has an answer uniform over all quaternionic reflections: exactly which reflections admit at least one generating pulse from the local stabilizer?  The answer is simply that the reflection must move the local context.  Moreover, a fixed list of six pulses suffices to test every reflection.

In the notation $Z(\mathbf a,\mathbf b)$ used next, the bold vectors are precisely the quaternionic velocity coordinates $q_a,q_b$ from \eqref{eq:gamma}.
For $T\in\mathscr C$ and $\mathbf a\in P_A$ with $T\mathbf a\notin P_A$, define
\[
H(\mathbf a,T)=\RR1\oplus\Span\{\mathbf a,T\mathbf a,\mathbf a\times T\mathbf a\}.
\]

\begin{lemma}[A forced context and an affine obstruction]
\label{lem:universal-affine}
Let $T\in\mathscr C$ and let $0\ne \mathbf a\in P_A$ satisfy $T\mathbf a\notin P_A$.  The subspace $H(\mathbf a,T)$ is a $T$-invariant quaternionic context distinct from $H_A$.  Suppose that $Z=Z(\mathbf a,\mathbf b)\in\mathfrak k_A$ has three pairwise distinct positive frequencies which are not in the ratio $1:2:3$.  Then
\[
\Lie\{Z,TZT\}\ne\mathfrak g_2
\quad\Longleftrightarrow\quad Z\in\mathfrak k_{H(\mathbf a,T)}.
\]
For fixed $\mathbf a,T$, the set of $\mathbf b\in\operatorname{Im}\HH$ on the right is either empty or an affine subspace of dimension at most one.
\end{lemma}

\begin{proof}
The independent imaginary vectors $\mathbf a$ and $T\mathbf a$ generate the indicated quaternionic subalgebra by alternativity.  The involution exchanges them and sends their cross product to its negative, so the context is $T$-invariant.  Since $T\mathbf a\notin P_A$, it is different from $H_A$.

The frequency hypothesis gives $\operatorname{rank}Z=6$ and $\ker Z=\RR \mathbf a$.  The conjugate has kernel $\RR T\mathbf a$, so the pair has no common fixed vector.  If its Lie algebra $L$ is proper, the compact maximal-subalgebra list recalled in \Cref{subsec:compact-maximals} leaves a vector stabilizer, a quaternionic stabilizer, or a principal $\mathfrak{su}(2)$.  The first is excluded by the common kernels; the last is excluded by the frequency ratios.  Thus $L\subset\mathfrak k_H$ for some context $H$.

In this situation $\dim L>3$.  Indeed, an abelian $L$ would preserve the one-dimensional kernel of $Z$, and skew symmetry would force its elements to kill that line, contrary to the absence of a common fixed vector.  A nonabelian compact Lie algebra of dimension at most three is $\mathfrak{su}(2)$.  In $\mathfrak k_H\cong\mathfrak{su}(2)\oplus\mathfrak{su}(2)$, its projections onto the two factors are either zero or isomorphisms.  A factor embedding has either rank-four elements or positive frequencies in the ratio $1:1:2$; a graph embedding has equal quaternionic norms in the two factors and hence rank-four elements.  Indeed, in the coordinates $(q_a,q_b)$ both factors use the same quaternionic commutator, and a graph is the graph of an automorphism of $\mathfrak{su}(2)$.  Since $\Aut(\mathfrak{su}(2))\cong\SO(3)$, that automorphism preserves the quaternionic norm.  This concerns the coordinate norms, not the differently weighted restrictions of the ambient trace metric.  These assertions follow directly from the stabilizer frequencies $2\alpha,|\beta-\alpha|,\beta+\alpha$ in \Cref{prop:failure}.  None is compatible with the assumed spectrum.

The generated algebra is $\Ad_T$-stable, so it lies in $\mathfrak k_H\cap\mathfrak k_{T(H)}$.  By \Cref{lem:nonfixed-context-bound}, dimension greater than three forces $T(H)=H$.  The skew restriction of $Z$ to the odd-dimensional plane $P_H$ has a kernel, so its unique kernel line $\RR \mathbf a$ lies in $P_H$.  Invariance gives $T\mathbf a\in P_H$ and therefore $H=H(\mathbf a,T)$.  This proves necessity.  Conversely, $Z\in\mathfrak k_{H(\mathbf a,T)}$ and $T(H(\mathbf a,T))=H(\mathbf a,T)$ place both generators in the same proper stabilizer.

If the affine set of $\mathbf b$ is nonempty, its direction space identifies with
\[
\mathfrak k_A^{\mathrm{pt}}\cap\mathfrak k_{H(\mathbf a,T)}.
\]
This is a subalgebra of $\mathfrak k_A^{\mathrm{pt}}\cong\mathfrak{su}(2)$, so it has dimension zero, one, or three.  Dimension three would mean that the entire pointwise stabilizer preserves $P_{H(\mathbf a,T)}$.  Its representation on $V$ is the direct sum of the trivial three-space $P_A$ and an irreducible real four-space.  Thus $P_A$ is its only invariant three-plane, forcing $H(\mathbf a,T)=H_A$, a contradiction.  The dimension is therefore at most one.
\end{proof}

Choose the following fixed local coordinate vectors:
\[
\mathbf a^{(1)}=e_1,\qquad \mathbf a^{(2)}=e_2,\qquad
\mathbf b^{(1)}=e_1+e_2,\quad \mathbf b^{(2)}=e_1+e_3,\quad \mathbf b^{(3)}=e_2+e_3,
\]
and define six fixed derivations with integer coefficients in the local derivation basis below
\[
U_{ij}=6Z(\mathbf a^{(i)},\mathbf b^{(j)}),\qquad i\in\{1,2\},\quad j\in\{1,2,3\}.
\]
Their coefficients in the ordered local basis $(D_{12},D_{13},D_{23},D_{45},D_{46},D_{47})$ are the rows
\begin{equation}
\begin{pmatrix}
0&-1&4&2&2&0\\
1&0&4&2&0&2\\
1&-1&3&0&2&2\\
0&-4&1&2&2&0\\
1&-3&1&2&0&2\\
1&-4&0&0&2&2
\end{pmatrix},
\label{eq:six-universal-pulses}
\end{equation}
in the order $(11,12,13,21,22,23)$.  

\Needspace{10\baselineskip}
\begin{theorem}[A six-pulse test for all quaternionic reflections]
\label{thm:universal-reflections}
For every $T\in\mathscr C$, the following are equivalent:
\begin{enumerate}[label=\textnormal{(\roman*)},leftmargin=9mm]
\item $T(H_A)\ne H_A$;
\item some $Z\in\mathfrak k_A$ satisfies $\Lie\{Z,TZT\}=\mathfrak g_2$;
\item at least one of the six $U_{ij}$ satisfies $\Lie\{U_{ij},TU_{ij}T\}=\mathfrak g_2$.
\end{enumerate}
If these conditions hold, the generating elements form a nonempty real Zariski-open subset of $\mathfrak k_A$ and have full Lebesgue measure.  For any other local context, conjugating these six pulses gives the same assertion.
\end{theorem}

\begin{proof}
If $T$ preserves $H_A$, then $T\mathfrak k_AT=\mathfrak k_A$, so no such pair can generate $\mathfrak g_2$.  This proves (ii)$\Rightarrow$(i), while (iii)$\Rightarrow$(ii) is immediate.

Suppose $T(H_A)\ne H_A$.  At least one of $Te_1,Te_2$ is outside $P_A$: otherwise their cross product $Te_3=T(e_1\times e_2)$ also lies in $P_A$, forcing $T(P_A)=P_A$.  Fix such an $\mathbf a^{(i)}$.  Each $Z(\mathbf a^{(i)},\mathbf b^{(j)})$ has $\alpha=1$ and $\beta=\sqrt2$, hence positive frequencies
\[
2,\quad \sqrt2-1,\quad \sqrt2+1.
\]
They are pairwise distinct and not in the ratio $1:2:3$.  By \Cref{lem:universal-affine}, the failing $\mathbf b^{(j)}$ must lie in an affine subspace of dimension at most one.  The three distinct $\mathbf b^{(j)}$ lie on the sphere of radius $\sqrt2$; a line meets that sphere in at most two points.  Therefore at least one of them generates.  Scaling by six does not change Lie generation.  Inverting \eqref{eq:terminal-qa-qb} gives the integer coefficients in \eqref{eq:six-universal-pulses}.

To prove the assertion about the set of pulses, fix $T$ and set
\[
E_1=\Span\{Z,TZT\},\qquad E_{r+1}=E_r+[Z,E_r]+[TZT,E_r].
\]
Inductively, $E_r$ is spanned by iterated brackets of length at most $r$, with
successive brackets taken with $Z$ or $TZT$.  These words span the generated
Lie algebra: the Jacobi identity rewrites an arbitrary bracket in this form.
If $E_{r+1}=E_r$, the space is invariant under both adjoint generators, so all
later words lie in $E_r$ and $E_r=\Lie\{Z,TZT\}$.  Otherwise its dimension
strictly increases.  Since $\dim\mathfrak g_2=14$, words of length at most
fourteen suffice for every pulse.
Their coordinate columns are polynomial in $Z$.  Full generation is thus the
nonvanishing of at least one of finitely many $14\times14$ minors, a real
Zariski-open condition.  The generating pulse proved above makes one such
minor a nonzero polynomial.  The nongenerating set is contained in its zero
set and therefore has Lebesgue measure zero.  Finally, conjugation transports
the local stabilizer, the reflection, and the six candidates to any other
quaternionic context.
\end{proof}

The six pulses are a fixed menu from which one may select a successful pulse after $T$ is specified.  The theorem does not require six simultaneous continuous controls, nor does it assert that six is the smallest possible menu.  The reflections for which every local pulse fails are exactly those commuting with $\sigma_{H_A}$: preservation of $H_A$ is equivalent to preservation of its orthogonal splitting.  \Cref{thm:invariant-contexts}, applied to $H_A$, describes this exceptional reflection family as the single reflection $\sigma_{H_A}$ together with the four-dimensional family of other invariant contexts.

\begin{remark}[The original pulse works for generic reflections]
\label{thm:generic-reflections}
Keep the original pulse $X=D_{46}+D_{47}$ fixed.  There is a nonempty real Zariski-open subset $\mathscr U\subset\mathscr C$, dense and of full invariant measure, such that
\[
\Lie\{X,TXT\}=\mathfrak g_2\qquad(T\in\mathscr U).
\]
\end{remark}

\begin{proof}
Let $D(T,Z)$ be the determinant of the fourteen words in \Cref{thm:single}, evaluated at $(Z,TZT)$.  It is polynomial in the entries of $T$ and the coordinates of $Z$.  Since $D(S,X)=2^{88}3^9\ne0$, the function $T\mapsto D(T,X)$ is nonzero real analytic on the connected manifold $\mathscr C$.  Its zero set has empty interior and measure zero.  Its nonvanishing locus supplies $\mathscr U$.
\end{proof}

Combining this local existence theorem with \Cref{thm:universal-locus}, the restriction $\mathcal U_{38}(T,\cdot)|_{\mathfrak k_A}$ is nonzero exactly when $T$ moves $H_A$.
If $T$ preserves $H_A$, its orthogonal projector onto $P_A$, minus $3I/7$, is a nonzero symmetric traceless operator commuting with $T$ and every local $Z$; hence $\mathcal R_{30}(T,\cdot)|_{\mathfrak k_A}=0$.
Otherwise the menu supplies a point with $\mathcal U_{38}>0$.
This is the local nonvacuity statement, separate from the full-domain factorization.  The next section concerns a different use of an explicit generating pair: local coordinates from a finite-word alphabet.


\section{Restricted-alphabet local coordinates}
\label{sec:charts}
The failure analysis above is independent of the local chart construction.  We now return to the successful pulse $X$ and use its fixed orbit under the same switch to obtain explicit coordinates.

Set
\[
A=e^{\pi X/(8\sqrt2)},\qquad a=A^{-1}.
\]

The lower-case letter $a$ in this section is the inverse group element, not a scalar quaternionic coordinate.  A word in $\{A,a,S\}$ is the ordered matrix product of its symbols, acting on vectors from right to left; $I$ is the empty word.  The product defining $\Psi_{\mathcal W}$ is ordered by increasing $j$ from left to right.  A local chart here means an analytic diffeomorphism between a neighborhood of $0$ in $\RR^{14}$ and a neighborhood of the identity in $G_2$.
By \Cref{prop:pulse}, $A^8=I$ and $A^4\ne I$; hence both $A$ and $a=A^7$ are obtained by positive pulse time.  For a list $\mathcal W=(g_1,\ldots,g_{14})$ of fixed words in $\{A,a,S\}$, define
\[
\Psi_{\mathcal W}(t_1,\ldots,t_{14})
=\prod_{j=1}^{14}g_j e^{t_jX}g_j^{-1}.
\]
Its left-trivialized differential at the origin has columns $\Ad_{g_j}X$.

Use the following two fixed lists of conjugator words:
\begin{align*}
\mathcal W_{\mathrm{lean}}=(&I,S,AS,aS,AAS,aaS,SAS,SaS,AAAS,\\
&ASAS,aSAS,SAAS,ASAAS,SASAS)
\end{align*}
\begin{align*}
\mathcal W_{\mathrm{bal}}=(&I,AAS,aaS,AASAS,AASaS,aaSAS,aaSaS,\\
&ASAAAS,aSAAAS,SAAAAS,SAASAS,SAASaS,SaaSAS,SaaSaS)
\end{align*}
For either list, define its coordinate matrix $C_{\mathcal W}$ by
\[
\Ad_{g_j}X=\sum_{i=1}^{14}(C_{\mathcal W})_{ij}D_i,
\]
where $D_i$ is the $i$th element of the ordered basis \eqref{eq:g2basis}.
Write $C_{\rm lean}$ and $C_{\rm bal}$ for the two resulting matrices.

\begin{proposition}[Two explicit local charts]
\label{thm:charts}
The coordinate matrices of the two fixed word lists satisfy
\begin{equation}
\det C_{\mathrm{lean}}=\frac{10-7\sqrt2}{32}\ne0,
\label{eq:chartdet-lean}
\end{equation}
\begin{equation}
\det C_{\mathrm{bal}}=-\frac{140+91\sqrt2}{16}\ne0.
\label{eq:chartdet-balanced}
\end{equation}
Consequently $\Psi_{\mathcal W_{\mathrm{lean}}}$ and
$\Psi_{\mathcal W_{\mathrm{bal}}}$ are analytic local diffeomorphisms at the origin.
\end{proposition}

\begin{proof}
Every factor $g_j e^{t_jX}g_j^{-1}$ equals the identity at $t_j=0$.
Differentiating the product at the origin therefore gives
\[
(d\Psi_{\mathcal W})_0(t_1,\ldots,t_{14})
=\sum_{j=1}^{14}t_j\Ad_{g_j}X.
\]
In the basis \eqref{eq:g2basis}, its matrix is $C_{\mathcal W}$.
Appendix~\ref{app:chartcertificates} expresses $A$ by spectral projectors of
$X^2$, constructs all conjugator columns over $\mathbb Q(\sqrt2)$, and lists
the elimination pivots and row interchanges for both matrices.  Multiplying
those pivots with the corresponding permutation signs gives
\eqref{eq:chartdet-lean} and \eqref{eq:chartdet-balanced}.
Both differentials are therefore invertible.  The analytic inverse-function
theorem supplies neighborhoods of the origin and the identity on which each
map is an analytic diffeomorphism.
\end{proof}

For the conditioning comparison, set $G_{ij}=-\operatorname{tr}(D_iD_j)$.
The Gram matrix of the column operator is $C^{\mathsf T}GC$; its condition
number is the square root of the ratio of the largest to the smallest
eigenvalue.  If $\widehat G$ is the Gram matrix after normalizing every
column, the corresponding volume is $\sqrt{\det\widehat G}$.
With the invariant trace metric $\langle U,V\rangle_{\mathrm{tr}}=-\operatorname{tr}(UV)$, the corresponding column operators have condition numbers
\[
\operatorname{cond}_{\rm tr}(C_{\rm lean})\approx130.58,
\qquad
\operatorname{cond}_{\rm tr}(C_{\rm bal})\approx4.78.
\]
After normalizing each column, the parallelotope volumes are approximately $8.96\times10^{-6}$ and $4.79\times10^{-2}$, respectively.  The balanced list is therefore about $27$ times better conditioned and more than $5.3\times10^3$ times larger in normalized volume, at the cost of $28$ rather than $21$ raw occurrences of the switch symbol in the displayed conjugator words.  These raw symbol counts are a transparent word-level proxy, not an optimized physical switch count after cancellations between consecutive factors.


The number fourteen is the dimension of $G_2$, as for any local coordinate chart.  The specific assertion is the nonsingularity of the two prescribed finite-word differentials, not a separate parameter-optimality result.


Products of exponentials and their local differential equations are classical \cite{WeiNorman1964,Altafini2003}.  Here every variable direction is a conjugate of one local pulse by a word in the fixed switch.  The two charts show that the same restricted alphabet can have markedly different local conditioning without adding a continuous control direction.


\subsection{Density of the two finite-time letters}
The fixed finite-time pulse and the reflection have a stronger property than nondiscreteness.
Define $\Gamma=\langle A,S\rangle$, the subgroup of finite words in $A,A^{-1},S$.
The following proof uses an exact adjoint-commutant certificate in Appendix~\ref{app:density-certificate}.

\begin{proposition}[A dense finite alphabet]\label{prop:dense-alphabet}
The subgroup $\Gamma$ is dense in $G_2$:
\[
\overline{\langle A,S\rangle}=G_2.
\]
For every nonzero $U\in\mathfrak g_2$, the real span of $\{\Ad_wU:w\in\Gamma\}$ is $\mathfrak g_2$.
In particular, some list of fourteen finite words gives an invertible differential when used as conjugators of $X$.
\end{proposition}
\begin{proof}
Exact multiplication of the matrices in Appendix~\ref{app:chartcertificates} gives
\[
\operatorname{tr}(A^2SAS)=-\frac12.
\]
A finite-order complex matrix has roots of unity as eigenvalues, so its trace is an algebraic integer.
A rational algebraic integer is an integer; hence $A^2SAS$ has infinite order and $\Gamma$ is infinite.

Let $\mathsf A=\Ad_A$ and $\mathsf S=\Ad_S$ on $\mathfrak g_2$.
Appendix~\ref{app:density-certificate} gives an explicit nonzero $99\times99$ minor over $\mathbb Q(\sqrt2)$ which proves
\[
\{M\in\operatorname{End}_{\RR}(\mathfrak g_2):M\mathsf A=\mathsf AM,
 M\mathsf S=\mathsf SM\}=\RR I_{\mathfrak g_2}.
\]
Both adjoint operators preserve the positive trace metric.  If a proper nonzero invariant real subspace existed, its orthogonal projector would commute with both operators and would not be scalar.
Thus the adjoint action of $\Gamma$ is irreducible over $\RR$.

The topological closure $H=\overline\Gamma$ is a compact Lie subgroup by the closed-subgroup theorem; see \cite{Hall2000} for the matrix-group theory.
If $\Lie(H)$ were zero, $H$ would be a compact discrete Lie group and hence finite, contrary to the infinite order just proved.
The nonzero subspace $\Lie(H)$ is invariant under $\Ad_\Gamma$.
Irreducibility therefore gives $\Lie(H)=\mathfrak g_2$; consequently $H$ is an open subgroup of the connected group $G_2$, so $H=G_2$.

For nonzero $U$, the span of its $\Gamma$-orbit is a nonzero invariant subspace and is therefore all of $\mathfrak g_2$.
Choose fourteen independent orbit elements.  Their product of conjugated pulses has those elements as its differential at the origin, proving the final assertion.
\end{proof}

Density is an approximation statement for finite words in the two fixed letters, not equality of this countable subgroup with $G_2$.
It differs from the exact finite-product generation in \Cref{thm:single}, where a continuous pulse time is available.
The existence of some nonsingular word list follows from \Cref{prop:dense-alphabet}; the specific content of \Cref{thm:charts} is the two printed short lists, their length bounds, exact determinants, and conditioning data.  No optimality or convergence rate is asserted.

\section{Relation to classical generation problems}
Classical two-generation and generic generation provide the background \cite{Kuranishi1951,AlbuquerqueSilvaLeite1989,Bois2009,DetinkoDeGraaf2020,Chirvasitu2021}.
The reflected-pair formulation itself has an antecedent in the swap construction of Bauer, Levaillant, and Freedman \cite{BauerLevaillantFreedman2014}.
On the full algebra, the parity bijection makes the present problem exactly generation by an independent even and odd element of the compact symmetric pair $(\mathfrak g_2,\mathfrak{so}(4))$.
The additional local restriction is the six-dimensional subspace $\mathcal W_A(T)$, not the mere writing of a pair as $(Z,TZT)$.

The compact maximal-subalgebra classification and associative-plane geometry are imported structure \cite{DraperMartin2025,Dynkin1952b,Cacciatori2005,KnarrStroppel2025}; commutant tests for irreducibility are classical \cite{ZeierSchulteHerbrueggen2011}.
The explicit construction here combines odd-dimensional averaging with a Gram determinant on fifteen symmetric unknowns, and a degree-eight identity for the small-algebra branch.
The maximal list proves that no further proper type remains.
This is more informative than the automatic sum of squared bracket minors, but does not establish minimal polynomial degree or an algorithmic advantage.
The principal spectral realization in \Cref{cor:universal-spectral} uses the classical $\SU(3)$ description of a vector stabilizer and the displayed principal triple.  The normalizer and centralizer argument in \Cref{prop:principal-reflection-fiber} then gives all realizing reflections for each such pulse, without classifying the reducible branches of arbitrary fibers.

The fixed-fiber geometry is logically different from uniform decision.
\Cref{thm:visible-contexts,thm:low-complete,thm:rankfour-context-complete} identify the invariant-context and low-dimensional mechanisms.
Their irredundant geometric union (\Cref{cor:geometric-irredundance}), its dimension, and the six irredundant scalar systems are not obtained merely by writing $\mathcal U_{38}=0$.
The homogeneous polynomial $\mathcal P_{34}$ is only an encoding of those systems.
The path-connectivity conclusion uses a general codimension-two argument \cite{BochnakCosteRoy1998}; the specific input is the dimension calculation for this fiber.
The genericity statements follow from the explicit witnesses and standard polynomial or analytic zero-set arguments, and are recorded as consequences rather than separate novelty claims.

The local extension formulas are orthogonal-splitting calculations with the octonionic normalization \cite{ChemtovKarigiannis2022}.
Products of exponentials are classical \cite{WeiNorman1964,Altafini2003}; the content of \Cref{thm:charts} is the specified finite-time word lists and their exact invertibility certificates.
The two finite-time letters generate a dense subgroup by \Cref{prop:dense-alphabet}, so the existence of some nonsingular finite-word list is automatic.  The prescribed short lists and their conditioning are the specific content.
The charts do not claim optimal word length, switching cost, or physical realization.

\section{Conclusion}
For a quaternionic reflection, generation by an even and an odd element of the associated compact symmetric pair is determined by an explicit product $\mathcal R_{30}\mathcal H_8$.
The factors detect reducibility and generated dimension at most three.  The principal spectral cone is exactly the set of nonzero pulses for which a reflected principal closure can be realized.
These are uniform statements on the full algebra, without the additional local constraint.

When the pulse is required to lie in one quaternionic stabilizer, the six-pulse theorem gives a fixed successful menu for every reflection moving that context.
For the distinguished reflection $S$, the four geometric families and six irredundant systems describe the failure set more finely than a uniform determinant alone.
The sum-of-squares polynomial $\mathcal P_{34}$, genericity observations, and codimension consequences organize that classification but do not constitute independent classification mechanisms.
The printed certificates also give a concrete generating pulse and two finite-word local coordinate maps.  The fixed finite-time letters generate a dense subgroup; the chart certificates provide explicit short lists rather than merely asserting their existence.

\section*{Exact reconstruction}
\addcontentsline{toc}{section}{Exact reconstruction}
\label{sec:reproducibility}

The Fano multiplication convention, the ordered derivation basis, and the matrices and elimination certificates in the appendices determine all finite calculations used in the proofs.  Matrices act on column vectors; the $j$th column of a matrix is the image of the $j$th basis vector.  Matrix commutators are $[A,B]=AB-BA$.  Starting from the oriented Fano triples in \eqref{eq:fano}, form the $8\times8$ left- and right-multiplication matrices $L_{e_i},R_{e_i}$ in the basis $(1,e_1,\ldots,e_7)$ and then
\[
D_{ij}=[L_{e_i},L_{e_j}]+[L_{e_i},R_{e_j}]+[R_{e_i},R_{e_j}].
\]
Delete the scalar row and column to obtain the action on $V=\operatorname{Im}\OO$, and express every subsequent derivation in the ordered basis \eqref{eq:g2basis}.  These instructions determine every integer matrix used below without an unprinted sign convention.

The octonion multiplication, derivation matrices, and coordinate bracket certificates have rational coefficients.  Parametrized identities and the scalar classifier are taken in the corresponding polynomial or rational-function rings over $\mathbb Q$.  The chart matrices lie in $\mathbb Q(\sqrt2)$, the sharp intersection witness in $\mathbb Q(\sqrt3)$, and the principal construction in $\mathbb Q(\sqrt2,\sqrt3,\sqrt5)$.  The rank-four family is evaluated in
\[
\mathbb Q(\tau)[\kappa]\Big/\left(\kappa^2-\frac{(1-\tau)(\tau^2+3)}{\tau}\right).
\]
All determinant and rank identities are taken over the indicated exact coefficient domains.  The decimal condition numbers and normalized volumes are approximations used only to compare the two already nonsingular chart matrices.

The finite claims can be checked in the following order.
\begin{center}\small
\begin{tabular}{@{}>{\raggedright\arraybackslash}p{3.1cm}>{\raggedright\arraybackslash}p{5.0cm}>{\raggedright\arraybackslash}p{7.2cm}@{}}\toprule
Claim & Input to reconstruct & Printed certificate \\\midrule
\Cref{thm:single} & $S,X,Y$ and the fourteen bracket words & Appendix~\ref{app:single-brackets}: all columns and determinant $2^{88}3^9$ \\
Vector failures & local basis, $Z_*$ and $Z'$ & Appendices~\ref{app:failure-certificate} and \ref{app:projective-failure-certificates}: exact eight-word matrices \\
Projective/resonant $\mathfrak{so}(4)$ families & equations \eqref{eq:projective-context-family}--\eqref{eq:resonant-pulse} & Appendix~\ref{app:projective-failure-certificates}: constraint matrix, Schur pivots and Gram determinants \\
Low-dimensional and rank-four terminal strata & coordinates \eqref{eq:terminal-qa-qb} and \eqref{eq:rankfour-parameter}--\eqref{eq:rankfour-new-frame} & Appendix~\ref{app:terminal-certificates}: normalized matrices, Casimir data and exact Gram blocks \\
$\mathcal P_{34}$ & equations \eqref{eq:scalar-invariants}--\eqref{eq:classifier-factors} & \Cref{thm:scalar-classifier,thm:polynomial-classifier}; direct sums-of-squares evaluation \\
Six-pulse theorem & the six rows in \eqref{eq:six-universal-pulses} & proof of \Cref{thm:universal-reflections}; the forced-context and affine-obstruction arguments \\
$\mathcal R_{30},\mathcal H_8,\mathcal U_{38}$ & equations \eqref{eq:R30-gram}, \eqref{eq:H8}, \eqref{eq:U38} & Appendix~\ref{app:universal-certificates}: canonical symmetric basis, normalization and principal example \\
Finite-alphabet density & $A,S$ and their adjoint matrices & Appendix~\ref{app:density-certificate}: exact commutant minor and nonintegral trace \\
Local charts & the two word lists of \Cref{thm:charts} & Appendix~\ref{app:chartcertificates}: spectral-projector formula and exact elimination pivots \\
\bottomrule
\end{tabular}
\end{center}

Exhaustiveness of the universal criterion follows from the compact maximal-subalgebra alternatives stated in \Cref{subsec:compact-maximals}, together with the proved reducibility and small-closure criteria.  Likewise, the dimension argument in \Cref{cor:dimension-connectivity} uses the cited standard theorem on dimensions of semialgebraic images.  These are the principal external structural inputs; the paper-specific coordinate calculations are explicitly reconstructible from the preceding conventions.

\appendix
\section{Exact single-switch iterated-bracket certificate}
\label{app:single-brackets}

The following coordinate data give a finite certificate for \Cref{thm:single}.  In the ordered basis $(e_1,\ldots,e_7)$,
\[
S=\begin{pmatrix}
1&0&0&0&0&0&0\\
0&0&0&1&0&0&0\\
0&0&0&0&1&0&0\\
0&1&0&0&0&0&0\\
0&0&1&0&0&0&0\\
0&0&0&0&0&-1&0\\
0&0&0&0&0&0&-1
\end{pmatrix},
\]
\[
X=\begin{pmatrix}
0&2&-2&0&0&0&0\\
-2&0&0&0&0&0&0\\
2&0&0&0&0&0&0\\
0&0&0&0&0&-4&-4\\
0&0&0&0&0&-2&2\\
0&0&0&4&2&0&0\\
0&0&0&4&-2&0&0
\end{pmatrix},\qquad
Y=\begin{pmatrix}
0&0&0&2&-2&0&0\\
0&0&0&0&0&4&4\\
0&0&0&0&0&2&-2\\
-2&0&0&0&0&0&0\\
2&0&0&0&0&0&0\\
0&-4&-2&0&0&0&0\\
0&-4&2&0&0&0&0
\end{pmatrix}.
\]
The identity $Y=SXS^{-1}$ is immediate from these matrices.

\footnotesize
\Needspace{18\baselineskip}
\begin{longtable}{@{}r p{4.0cm} p{6.8cm}@{}}
\caption{Exact coordinates of the fourteen iterated brackets in the basis \eqref{eq:g2basis}.}\label{tab:brackets}\\
\toprule $i$&Word&Coordinates\\\midrule\endfirsthead
\toprule $i$&Word&Coordinates\\\midrule\endhead
1 & $X$ & $D_{46}+D_{47}$ \\
2 & $Y$ & $-D_{26}-D_{27}$ \\
3 & $[X,Y]$ & $-2D_{16}-2D_{17}+8D_{24}$ \\
4 & $[X,[X,Y]]$ & $36D_{14}+4D_{15}+40D_{26}+32D_{27}$ \\
5 & $[Y,[X,Y]]$ & $-36D_{12}-4D_{13}+40D_{46}+32D_{47}$ \\
6 & $[X,[X,[X,Y]]]$ & $304D_{16}+208D_{17}-352D_{24}-96D_{25}$ \\
7 & $[X,[Y,[X,Y]]]$ & $48D_{23}-48D_{45}$ \\
8 & $[Y,[Y,[X,Y]]]$ & $208D_{16}+208D_{17}-352D_{24}+96D_{25}$ \\
9 & $[[X,Y],[X,[X,Y]]]$ & $-2496D_{12}+64D_{13}+1760D_{46}+1888D_{47}$ \\
10 & $[[X,Y],[Y,[X,Y]]]$ & $-2496D_{14}+64D_{15}-1760D_{26}-1888D_{27}$ \\
11 & $[[X,[X,Y]],[Y,[X,Y]]]$ & $-16064D_{16}-13760D_{17}+19712D_{24}$ \\
12 & $[X,[[X,Y],[X,[X,Y]]]]$ & $5376D_{23}+768D_{45}$ \\
13 & $[[X,Y],[[X,Y],[X,[X,Y]]]]$ & $-157056D_{14}+1664D_{15}-92416D_{26}-120320D_{27}$ \\
14 & $[[X,Y],[[X,Y],[Y,[X,Y]]]]$ & $157056D_{12}-1664D_{13}-92416D_{46}-120320D_{47}$ \\
\bottomrule
\end{longtable}
\normalsize

The determinant of the coordinate columns is $2^{88}3^9$, as stated in \eqref{eq:bracketdet}.

\section{\texorpdfstring{Exact certificate for the $\mathfrak{su}(3)$ failure family}{Exact certificate for the su(3) failure family}}
\label{app:failure-certificate}

For the fixed vector $e_1$, use the ordered basis
\begin{align*}
\mathcal B_{e_1}=(&D_{23},\tfrac12D_{17}+D_{24},-\tfrac12D_{16}+D_{25},
\tfrac12D_{15}+D_{26},\\
&-\tfrac12D_{14}+D_{27},D_{45},-\tfrac12D_{13}+D_{46},
\tfrac12D_{12}+D_{47})
\end{align*}
of $\mathfrak g_2^{e_1}$.  Take the explicit pulse $Z=Z_*$ from \eqref{eq:su3-witness}, put $Y=\Ad_SZ$, and define
\[
F_1=Z,\ F_2=Y,\ F_3=[Z,Y],\ F_4=[Z,F_3],\ F_5=[Y,F_3],
\]
\[
F_6=[Z,F_4],\qquad F_7=[Z,F_5],\qquad F_8=[F_3,F_4].
\]
Their coordinate columns in $\mathcal B_{e_1}$ are
\[
\begin{pmatrix}
1&-4&0&-288&576&0&0&311040\\
0&0&48&0&0&-60480&58752&0\\
0&0&0&-1728&1728&0&0&93312\\
0&4&0&-720&576&0&0&-248832\\
0&0&24&0&0&-25056&24192&0\\
-4&1&0&-576&288&0&0&-279936\\
-4&0&0&576&-720&0&0&-311040\\
0&0&24&0&0&-27648&24192&0
\end{pmatrix}.
\]
Its determinant is
\[
41182968668869361664=2^{30}3^{20}\cdot11\ne0,
\]
which proves that the closure has dimension eight.  For comparison, the pulse $-3D_{12}+4D_{47}$ has exact closure dimensions $2\to3\to5\to6\to6$, showing that the vector-stabilizer criterion captures a substantial but not exhaustive component of the nongenerating locus.

\section{Projective-context and common-vector certificates}
\label{app:projective-failure-certificates}

The constraint matrix and uniform Gram determinants below support
\Cref{thm:visible-contexts,cor:quartic-cone,prop:uniform-so4}.
The subsequent eight-word matrix certifies the common-vector witness in
\Cref{prop:rankfour-failure}, and Appendix~\ref{app:rankfour-parametrization}
gives the cubic generators used in \Cref{prop:rankfour-complete}.
For the family in \eqref{eq:projective-context-family}, put
\[
\xi=ce_2+se_3,
\qquad \eta=ce_4+se_5,
\qquad \nu=(c^2-s^2)e_6+2cs\,e_7,
\qquad n=c^2+s^2.
\]
Direct multiplication gives
\[
\xi\times\eta=\nu,
\qquad \eta\times\nu=n\xi,
\qquad \nu\times\xi=n\eta,
\]
and the switch satisfies $S(\xi)=\eta$, $S(\eta)=\xi$, $S(\nu)=-\nu$.

For
\[
Z=x_0D_{12}+x_1D_{13}+x_2D_{23}+x_3D_{45}+x_4D_{46}+x_5D_{47},
\]
the twelve normal-block constraints expressing $Z(P_{H_{[c:s]}})\subseteq P_{H_{[c:s]}}$ are represented by
\[
\renewcommand{\arraystretch}{1.05}
\begin{pmatrix}
-4c&-4s&0&0&-2s&2c\\
0&0&0&0&0&0\\
0&0&0&0&0&0\\
0&0&4n&-2n&0&0\\
0&0&0&0&0&0\\
0&0&0&0&0&0\\
0&0&0&0&0&0\\
0&0&-2n&4n&0&0\\
-2cn&-2sn&0&0&4s(2c^2-s^2)&2c(-c^2+5s^2)\\
0&0&0&0&0&0\\
-2cn&-2sn&0&0&2s(-5c^2+s^2)&4c(c^2-2s^2)\\
0&0&2n^2&2n^2&0&0
\end{pmatrix}.
\]
Write $M_i$ for row $i$ of this matrix.  Only rows $1,4,8,9,11,12$
can be nonzero, and they satisfy
\[
M_{11}=nM_1-M_9,\qquad M_{12}=n(M_4+M_8).
\]
Rows $4,8$ are supported on columns $x_2,x_3$, with determinant $12n^2$.
On the other four columns, $M_1$ is nonzero in $x_0,x_1$, whereas
\[
M_9-\frac n2M_1
=(0,0,0,0,\,3s(3c^2-s^2),\,-3c(c^2-3s^2)).
\]
The sum of squares of the last two entries is $9n^3>0$.
Thus rows $1,4,8,9$ are independent and all other rows lie in their span,
proving rank four for every $n>0$.
Direct substitution annihilates the two coefficient columns of
$\mathsf P_{c,s},\mathsf Q_{c,s}$, and their positive Gram determinant below
proves independence.  Since the kernel has dimension two, those columns
form its complete basis.  The coordinate Gram determinant in this nonorthonormal coefficient basis is
\[
n^4(c^4+3c^2s^2+s^4),
\]
whereas the invariant trace-Gram determinant is
\[
2^6 3^3n^6.
\]
At $[1:0]$, $[0:1]$, and $[1:1]$ these specialize to the $H_{246}$, $H_{365}$, and $H_\star$ intersections, respectively.

In the complex coordinates of \Cref{cor:quartic-cone}, the expanded quartic is
\begin{align*}
\operatorname{Im}(z^\ast\omega^3)={}&8x_0^3x_4+24x_0^2x_1x_5-24x_0x_1^2x_4
-24x_0x_1x_4^2-24x_0x_1x_5^2\\
&-6x_0x_4^3-6x_0x_4x_5^2-8x_1^3x_5
+6x_1x_4^2x_5+6x_1x_5^3\\
&+2x_4^3x_5+2x_4x_5^3.
\end{align*}
The singular-locus calculation must account for the dependence $\omega=z-2iw$.  If $F=\operatorname{Im}(z^\ast\omega^3)$, then
\[
F_{x_0}=-6\operatorname{Re}(z^\ast\omega^2),
\qquad
F_{x_1}=6\operatorname{Im}(z^\ast\omega^2).
\]
Thus a singular point has $z^\ast\omega^2=0$.  When $z=0$ but $\omega\ne0$, the remaining derivatives are $F_{x_4}=\operatorname{Im}(\omega^3)$ and $F_{x_5}=-\operatorname{Re}(\omega^3)$, so they cannot both vanish.  Hence the singular locus is exactly $\omega=0$.

\paragraph{Uniform trace-Gram calculations.}
The three determinant identities for the simple, regular, and resonant pulse
families can be checked in one normalized slice.  Let
$T_0=(D_{23}+D_{45})/2$ and $g_\theta=\exp(\theta T_0)$, as in
\eqref{eq:residual-circle}.  For each of the three pulse polynomials $U$,
substitution of its coefficients gives
\[
[T_0,U]=-s\frac{\partial U}{\partial c}
       +c\frac{\partial U}{\partial s}.
\]
Consequently, writing $c=\varrho\cos\theta$, $s=\varrho\sin\theta$ with
$\varrho>0$, the simple and regular pulses are respectively
$\varrho\Ad_{g_\theta}U(1,0)$ and
$\varrho^5\Ad_{g_\theta}U(1,0)$.  For the resonant pulse,
\[
\widehat Z_{a;c,s}
=\varrho^3\Ad_{g_\theta}\widehat Z_{a/\varrho;1,0}.
\]
The circle commutes with $S$, and conjugation preserves the trace metric.
It remains to compute the six-word Gram matrices at $c=1,s=0$.

For clarity, the scalar pivots of symmetric elimination without row
interchanges are listed below.  Starting from a symmetric matrix
$\left(\begin{smallmatrix}\pi&r^{\mathsf T}\\r&M\end{smallmatrix}\right)$,
each elimination step retains the Schur complement $M-rr^{\mathsf T}/\pi$.
Set $h=a^2+1$ and $k=2a^2+1$ for the normalized resonant family.
\begin{center}
\renewcommand{\arraystretch}{1.3}
\begin{tabular}{c r r c}
\toprule
Pivot & Simple pulse & Regular pulse & Resonant pulse\\
\midrule
$1$ & $48$ & $1008$ & $48h$\\
$2$ & $48$ & $1008$ & $48k/h$\\
$3$ & $1344$ & $316224$ & $1344k$\\
$4$ & $9216$ & $20155392/7$ & $9216hk$\\
$5$ & $9216$ & $20155392/7$ & $9216k^2/h$\\
$6$ & $331776/7$ & $6530347008/61$ & $331776h^2k/7$\\
\bottomrule
\end{tabular}
\end{center}
Their products are, respectively,
\[
2^{46}3^{11},\qquad 2^{46}3^{39},\qquad
2^{46}3^{11}(a^2+1)^2(2a^2+1)^6.
\]
The six words have degrees $1,1,2,3,3,4$ in the pulse, so their Gram
determinant scales with the twenty-eighth power of its scalar multiplier.
The preceding equivariance and homogeneity therefore give
\eqref{eq:uniform-so4-gram}, \eqref{eq:uniform-regular-gram}, and
\eqref{eq:resonant-so4-gram} for all their parameters.

For the uniform pulse of \Cref{prop:uniform-so4}, the principal trace-Gram determinants along the bracket ladder are
\[
2^8 3^2n^2,
\qquad 2^{14}3^3\cdot7\,n^4,
\qquad 2^{34}3^7\cdot7\,n^{10},
\qquad 2^{46}3^{11}n^{14}.
\]
They certify ranks $2,3,5,6$ uniformly.  For the regular pulse \eqref{eq:uniform-regular-pulse}, the final six-word trace-Gram determinant is
\[
2^{46}3^{39}n^{70},
\]
which is nonzero at every projective point.  For the rank-four vector witness
\[
Z'=D_{12}-3D_{13}+2D_{47},\qquad Y'=\Ad_S Z',
\]
use the primitive integral basis
\begin{align*}
\mathcal B_{e_6+e_7}=(&D_{16}-D_{17},\ D_{16}+2D_{24},\ D_{16}-2D_{25},\\
&3D_{14}-D_{15}-2D_{26},\ D_{14}-3D_{15}-2D_{27},\ D_{23}-D_{45},\\
&3D_{12}-D_{13}+2D_{46},\ D_{12}-3D_{13}+2D_{47})
\end{align*}
of $\gtwo^{e_6+e_7}$.  Define
\[
F_1=Z',\quad F_2=Y',\quad F_3=[F_1,F_2],\quad F_4=[F_1,F_3],\quad F_5=[F_2,F_3],
\]
\[
F_6=[F_1,F_4],\qquad F_7=[F_1,F_5],\qquad F_8=[F_2,F_5].
\]
Their coordinate columns in $\mathcal B_{e_6+e_7}$ are
\[
\begin{pmatrix}
0&0&-60&0&0&29376&0&29376\\
0&0&24&0&0&-8640&0&-8640\\
0&0&0&0&0&5184&0&-5184\\
0&0&0&-216&0&0&0&0\\
0&1&0&-360&0&0&0&0\\
0&0&0&0&0&0&5184&0\\
0&0&0&0&216&0&0&0\\
1&0&0&0&360&0&0&0
\end{pmatrix}.
\]
Its determinant is
\[
2^{27}3^{20}=467988280328060928\ne0,
\]
so the closure is the full vector stabilizer.  Thus the closure is certified directly by the displayed matrix.

\subsection{Complete cubic parametrization of rank-four common-vector fibers}\label{app:rankfour-parametrization}
For a positive $S$-eigenvector
\[
w_+(a,b,c)=ae_1+b(e_2+e_4)+c(e_3+e_5),
\qquad (b,c)\ne(0,0),
\]
the unique projective annihilator in $\mathfrak k_A$ is generated by
\begin{align*}
Z_+(a,b,c)={}&-c(3b^2+c^2)D_{12}+2b^3D_{13}
-2a(b^2+c^2)D_{23}-a(b^2+c^2)D_{45}\\
&+b(3c^2-b^2)D_{46}+c(c^2-3b^2)D_{47}.
\end{align*}
Its normal-form vector satisfies
\[
q_a=-3(b^2+c^2)(ae_1+be_2+ce_3),
\qquad \lVert q_a\rVert=\lVert q_b\rVert,
\]
and
\[
\sum_{j=0}^{5}x_j(Z_+)^2
=(b^2+c^2)^2(5a^2+5b^2+2c^2),
\]
so the generator is nonzero away from the exceptional line $[e_1]$.

For a negative $S$-eigenvector
\[
w_-(a,b,r,t)=a(e_2-e_4)+b(e_3-e_5)+re_6+te_7,
\qquad (a,b)\ne(0,0),
\]
put $\upsilon=a^2t-2abr-b^2t$.  The unique projective annihilator is generated by
\begin{align*}
Z_-(a,b,r,t)={}&(3a^2b+art+b^3+br^2+2bt^2)D_{12}\\
&-(2a^3+2ar^2+at^2+brt)D_{13}+\upsilon D_{23}+2\upsilon D_{45}\\
&+(a^3-3ab^2+ar^2-at^2+2brt)D_{46}\\
&+(3a^2b+2art-b^3-br^2+bt^2)D_{47}.
\end{align*}
Here
\[
q_a=3(a^2+b^2+r^2+t^2)(ae_2+be_3),
\qquad \lVert q_a\rVert=\lVert q_b\rVert,
\]
so the generator is nonzero whenever $(a,b)\ne(0,0)$.  When $a=b=0$, the annihilator jumps from a line to the three-plane $\mathcal N_{re_6+te_7}$ of \Cref{prop:rankfour-failure}.  These formulas, together with the $e_1$ fiber $\mathcal N_{e_1}$, give the exhaustive incidence description in \Cref{prop:rankfour-complete}.

\section{Exact local-coordinate certificates and conditioning}
\label{app:chartcertificates}
For either word list in \Cref{thm:charts}, the coordinate matrix introduced in \Cref{sec:charts} satisfies
\[
\Ad_{g_j}X=\sum_{i=1}^{14}C_{ij}D_i
\]
in the ordered basis \eqref{eq:g2basis}, with $D_i$ denoting its $i$th element.  The exact determinants and the trace-metric diagnostics are summarized in the following table.

\begin{table}[ht]
\centering\small
\begin{tabular}{@{}lccccc@{}}
\toprule
Chart & $\det C$ & $\operatorname{cond}_{\rm tr}$ & Norm. volume & raw $S$ & max. length\\
\midrule
Lean & $(10-7\sqrt2)/32$ & $130.5778$ & $8.9648\times10^{-6}$ & 21 & 5\\
Balanced & $-(140+91\sqrt2)/16$ & $4.7828$ & $4.7934\times10^{-2}$ & 28 & 6\\
\bottomrule
\end{tabular}
\caption{Exact determinant and invariant trace-metric diagnostics for the two charts.  The raw $S$ count is taken before cancellations between consecutive factors.}
\end{table}

For example, the column corresponding to the word $AS$ is
\[
\Ad_{AS}X=\frac{\sqrt2}{2}D_{14}+\frac{\sqrt2-1}{2}D_{16}+\frac{1+\sqrt2}{2}D_{24}+\frac{1-\sqrt2}{2}D_{25}.
\]
The condition numbers are obtained from $C^TGC$, where $G_{ij}=-\operatorname{tr}(D_iD_j)$.  The coordinate determinant values depend on the chosen non-orthonormal derivation basis; nonsingularity and the condition numbers computed from the invariant trace metric do not.


To reconstruct the matrices without evaluating a transcendental exponential, set $Q=X^2$ and let $I_7$ be the identity on $V$.  The spectral projectors are
\[
P_0=\frac{(Q+8I_7)(Q+32I_7)}{256},\quad
P_8=-\frac{Q(Q+32I_7)}{192},\quad
P_{32}=\frac{Q(Q+8I_7)}{768}.
\]
Then the fixed pulse in \Cref{sec:charts} is exactly
\[
A=P_0+\frac{\sqrt2}{2}P_8+\frac14XP_8+\frac{\sqrt2}{8}XP_{32}.
\]
Multiplying the displayed conjugator words and expressing $gXg^{-1}$ in \eqref{eq:g2basis} determines every column over $\mathbb Q(\sqrt2)$.

\begin{samepage}
Here is a compact elimination certificate for the determinants.  In each column, choose the first nonzero pivot at or below the diagonal, swap its row with the diagonal row, and eliminate below without rescaling the pivot row.  Put $r=\sqrt2$.  The fourteen pivots, in order and split after the seventh entry, are
\begingroup\interdisplaylinepenalty=10000
\begin{align*}
\text{lean: }&\left(1,-1,\frac r2,1-r,-\frac12,-r,-\frac{r+1}{2};\right.\\[-1mm]
&\hspace{10mm}\left.r,2(r-2),\frac{1-r}{4},\frac r2,-1,-\frac12,\frac{1-r}{2}\right),\\
\text{balanced: }&\left(1,-\frac12,-r,-1,-\frac r2,\frac{r-3}{2},2;\right.\\[-1mm]
&\hspace{10mm}\left.-\frac{r+1}{4},-\frac{7r}{8},2,-\frac r2,
\frac{32+27r}{14},1,-4\right).
\end{align*}
\endgroup
\end{samepage}
The row swaps for the lean matrix are
\[
(1,13),(2,10),(4,5),(7,8),(8,13),(11,13),(12,14),(13,14),
\]
and those for the balanced matrix are
\[
(1,13),(2,3),(3,5),(9,10),(10,13),(12,13).
\]
Both permutations are even.  The products of the corresponding pivots are therefore the determinants in \eqref{eq:chartdet-lean} and \eqref{eq:chartdet-balanced}.  No numerical rank decision enters the certificate.


\subsection{Exact adjoint-commutant certificate}\label{app:density-certificate}
This certificate concerns \Cref{prop:dense-alphabet} and is independent of the two chart determinants.
Use the following ordered parity basis of $\mathfrak g_2$, with a semicolon separating the six positive and eight negative vectors for $\Ad_S$:
\[
\begin{split}
\mathcal E=(&D_{12}+D_{14},\ D_{13}+D_{15},\ -\tfrac12D_{16}+D_{25},\ D_{23}+D_{45},\\
 &-D_{26}+D_{46},\ -D_{27}+D_{47};\\
 &-D_{12}+D_{14},\ -D_{13}+D_{15},\ D_{16},\ D_{17},\ D_{24},\\
 &-D_{23}+D_{45},\ D_{26}+D_{46},\ D_{27}+D_{47}).
\end{split}
\]
If $P$ is the matrix of these columns in the ordered basis \eqref{eq:g2basis}, set
$\mathsf B=P^{-1}\mathsf A P$, where column $j$ of $\mathsf A$ is the coordinate vector of $AD_jA^{-1}$.
The exact spectral formula for $A$ above determines $\mathsf B$ over $\mathbb Q(\sqrt2)$; its entries lie in $\mathbb Z[1/2,\sqrt2]$.
In this basis $\mathsf S=\operatorname{diag}(I_6,-I_8)$, so every commuting endomorphism is block diagonal with $6\times6$ and $8\times8$ blocks.
Subtracting a scalar multiple of the identity, we may impose $M_{11}=0$, leaving $99$ unknowns.
Order their positions $(i,j)$ lexicographically within
\[
\mathcal J=\{(i,j):1\le i,j\le6\ \text{or}\ 7\le i,j\le14\}\setminus\{(1,1)\}.
\]
The $196\times99$ coefficient matrix $N$ of $[\mathsf B,M]=0$, with output entries listed row by row, is explicitly
\[
N_{14(r-1)+s,(i,j)}=\mathsf B_{ri}\delta_{js}-\delta_{ri}\mathsf B_{js},
\quad 1\le r,s\le14,\quad (i,j)\in\mathcal J.
\]
Select the following $99$ rows, in increasing order:
\[
\mathcal I=\{1,\ldots,70\}\cup\{77,\ldots,96\}\cup\{99\}
 \cup\{113,\ldots,119\}\cup\{125\}.
\]
Exact elimination gives
\[
\det N_{\mathcal I,:}
=\frac{9\bigl(981996834544536330001\sqrt2-1388752894942616432006\bigr)}{2^{166}}
\ne0.
\]
Nonvanishing also has a smaller independent arithmetic certificate: reduce in $\mathbb F_7$ under $\sqrt2\mapsto3$, which is valid since $3^2=2\pmod7$ and the denominators are powers of two.  The same minor has determinant $3\pmod7$.
Thus $N$ has column rank $99$ over $\mathbb Q(\sqrt2)$ and over its real embedding.  The only block-diagonal solution with $M_{11}=0$ is zero, proving that the full real commutant is scalar.
The trace $\operatorname{tr}(A^2SAS)=-1/2$ is obtained from the same seven-dimensional matrices.  Together these finite certificates supply both inputs to the density argument.

\section{Exact certificates for the terminal strata}
\label{app:terminal-certificates}
The normalized matrix below is the input to the elimination in
\Cref{thm:low-complete}; the following norms and Casimir spectra verify its
fixed-space assertions.  The final three Gram blocks certify the full
six-dimensional closure in \Cref{thm:rankfour-context-complete}.
All formulas use the Fano convention and the local basis fixed in the main text.

For $q_a=(a,1,0)$ and $q_b=(d,b,c)$, the local derivation is
\[
 Z=\begin{pmatrix}
 0&0&2&0&0&0&0\\
 0&0&-2a&0&0&0&0\\
 -2&2a&0&0&0&0&0\\
 0&0&0&0&a-d&1-b&-c\\
 0&0&0&-a+d&0&-c&b+1\\
 0&0&0&b-1&c&0&-a-d\\
 0&0&0&c&-b-1&a+d&0
 \end{pmatrix}.
\]
For the three low-dimensional normal forms $(a,d,b,c)=(0,0,\eta,0)$, the values
\[
 (\|A\|_{\rm tr}^2,\|B\|_{\rm tr}^2,\|[A,B]\|_{\rm tr}^2)
 =\begin{cases}(8,8,64),&\eta=1,\\(8,8,16),&\eta=-1,\\(24,24,192),&\eta=-3\end{cases}
\]
and the double-bracket equations prove three-dimensional closure exactly.  The axial form $q_a=0$, $q_b=(0,1,0)$ has corresponding norms $(2,2,1)$.  The Casimir of the trace-orthonormal basis $(A/\|A\|,B/\|B\|,[A,B]/\|[A,B]\|)$, with an overall minus sign, has eigenvalues
\[
\begin{array}{c|c}
\eta=1&0\text{ (multiplicity 3)},\;3/4\text{ (multiplicity 4)}\\
\eta=-1&0\text{ (multiplicity 1)},\;1/2\text{ (multiplicity 6)}\\
\eta=-3&2/3\text{ (multiplicity 3)},\;1/4\text{ (multiplicity 4)}\\
\text{axial}&0\text{ (multiplicity 1)},\;1/2\text{ (multiplicity 6)}.
\end{array}
\]
These give the representation and fixed-space distinctions used in \Cref{thm:low-complete}.
Equivalently, the four nonabelian low-dimensional families are obtained by residual-circle conjugation and nonzero rescaling of the following simple pulses:
\[
\begin{array}{c|c|c}
\text{family}&\text{representative}&\operatorname{rank}Z\\\hline
\eta=1&-2D_{13}+D_{46}&4\\
\eta=-1&-D_{13}-D_{46}&4\\
\eta=-3&-D_{46}&6\\
\text{axial}&-D_{13}+2D_{46}&4.
\end{array}
\]


For the six words $F$ in \Cref{thm:rankfour-context-complete}, trace invariance and symmetric parity split the Gram matrix into three $2\times2$ blocks, on the index pairs $(1,5)$, $(2,4)$, and $(3,6)$.  Set $D=\tau^2+3$ and $R=\tau^2-2\tau+3$.  The three determinants are
\begin{align*}
 B_{15}&=\frac{2^8 3^9(3-\tau)^4(1-\tau)^2(1+\tau)^2D^8}{\tau^2},\\
 B_{24}&=\frac{2^{10}3^9(3-\tau)^2(1-\tau)^2(1+\tau)^4D^6R^2}{\tau^2},\\
 B_{36}&=\frac{2^{14}3^{13}(3-\tau)^4(1-\tau)^2(1+\tau)^4D^{10}R^2}{\tau^2}.
\end{align*}
Multiplying them gives \eqref{eq:rankfour-new-gram}.  Each is positive for $0<\tau<1$.  The computation takes place exactly in
\[
 \mathbb Q(\tau)[\kappa]\big/
 \left(\kappa^2-\frac{(1-\tau)(\tau^2+3)}\tau\right),
\]
so neither the determinant nor its nonvanishing is inferred from sampled parameter values.

Every determinant and rank statement used above is an exact identity over the indicated rational, quadratic, or rational-function field.  Numerical rank tests are not used in the proofs.

\section{Matrix formulas for the universal criterion}
\label{app:universal-certificates}

Choose an orthogonal matrix $R$ whose first three columns span the positive eigenspace of $T$, and whose last four columns span its negative eigenspace.  In that frame,
\[
R^{\mathsf T}TR=\operatorname{diag}(I_3,-I_4),\qquad
Z'=R^{\mathsf T}ZR=\begin{pmatrix}A&B\\-B^{\mathsf T}&D\end{pmatrix},
\]
where $A\in\mathfrak{so}(3)$, $D\in\mathfrak{so}(4)$ and $B\in\mathbb R^{3\times4}$.  The symmetric commutant condition is the following linear system in the entries of two symmetric matrices $U\in\mathbb R^{3\times3}$, $W\in\mathbb R^{4\times4}$:
\begin{equation}
AU=UA,\qquad DW=WD,\qquad BW=UB,
\qquad \operatorname{tr}U+\operatorname{tr}W=0.
\label{eq:block-commutant}
\end{equation}
There are fifteen independent unknowns after the trace constraint.  A nonzero solution is equivalent to $\mathcal R_{30}(T,Z)=0$.

For a fixed basis of these unknowns, let $E_{ij}$ be the matrix units in this orthonormal frame.  Use the six diagonal matrices $E_{ii}-E_{77}$, $1\le i\le6$, followed by the nine matrices $E_{ij}+E_{ji}$ with $i<j$ in either of the blocks $\{1,2,3\}$ or $\{4,5,6,7\}$.  Call the resulting list $F_1,\ldots,F_{15}$.  Its Frobenius Gram determinant is $7\cdot2^9$: the diagonal part has Gram matrix $I_6+\boldsymbol1\boldsymbol1^{\mathsf T}$, where $\boldsymbol1=(1,1,1,1,1,1)^{\mathsf T}\in\RR^6$, and the other nine squared norms are two.  Thus
\begin{equation}
\mathcal R_{30}(T,Z)
=\frac{1}{7\cdot2^9}\det\bigl(\operatorname{tr}([Z',F_i][Z',F_j])\bigr)_{i,j=1}^{15}.
\label{eq:explicit15}
\end{equation}
The commutators in this formula are symmetric.  Formula \eqref{eq:R30-framefree} gives the same value in any fixed frame without constructing $R$.

For the original generating pair $(S,X)$, exact substitution gives
\[
\mathcal R_{30}(S,X)=2^{36}3^4\cdot7\cdot13\cdot265747,
\qquad
\mathcal H_8(S,X)=2^{11}3^3\cdot11.
\]
These provide a normalization check for both formulas.  At the rank-four and vector-stabilizer failures printed earlier, the first determinant vanishes; at the low-dimensional forms, the degree-eight factor vanishes as well.  Their interpretation is proved in \Cref{thm:universal-locus}, independently of these examples.

For \Cref{prop:principal-moving}, the reflection has the sparse matrix
\[
T_*=\begin{pmatrix}
0&0&0&0&-\sqrt6/4&-\sqrt{10}/4&0\\
0&0&0&0&\sqrt{10}/4&-\sqrt6/4&0\\
0&0&-1&0&0&0&0\\
0&0&0&1/4&0&0&-\sqrt{15}/4\\
-\sqrt6/4&\sqrt{10}/4&0&0&0&0&0\\
-\sqrt{10}/4&-\sqrt6/4&0&0&0&0&0\\
0&0&0&-\sqrt{15}/4&0&0&-1/4
\end{pmatrix}.
\]
It is symmetric, its square is the identity, its trace is $-1$, and multiplication gives $T_*K_{\rm p}T_*=P_{\rm p}$.  Its membership in $G_2$ follows from the polynomial-exponential identity in \eqref{eq:principal-reflection}; it can also be checked directly on the Fano products.  The identities in \eqref{eq:principal-triple} and the scalar Casimir certify the principal closure, without a numerical irreducibility test.

\paragraph{An azimuthal realizing reflection.}
For the example after \Cref{prop:principal-reflection-fiber}, set
\[
u_0=\exp(\pi K_{\rm p}/2)=\operatorname{diag}(1,-I_2,J,-J),
\qquad J=\begin{pmatrix}0&-1\\1&0\end{pmatrix}.
\]
The blocks use the ordered basis $(e_1; e_2,e_3; e_4,e_5; e_6,e_7)$.
Thus $u_0\in C_{\rm p}$, and direct conjugation of the displayed $T_*$ gives
\[
T_{\rm az}:=u_0T_*u_0^{-1}
=\frac14\begin{pmatrix}
0&0&0&\sqrt6&0&0&\sqrt{10}\\
0&0&0&\sqrt{10}&0&0&-\sqrt6\\
0&0&-4&0&0&0&0\\
\sqrt6&\sqrt{10}&0&0&0&0&0\\
0&0&0&0&1&-\sqrt{15}&0\\
0&0&0&0&-\sqrt{15}&-1&0\\
\sqrt{10}&-\sqrt6&0&0&0&0&0
\end{pmatrix}.
\]
This matrix is symmetric, has square $I_7$ and trace $-1$, and satisfies
\[
u_0P_{\rm p}u_0^{-1}=Q_{\rm p},\qquad
T_{\rm az}K_{\rm p}T_{\rm az}=Q_{\rm p}.
\]
Its membership in $G_2$ follows from $u_0,T_*\in G_2$, independently of the matrix check.
Since $K_{\rm p}$ and $Q_{\rm p}$ generate $\mathfrak s_{\rm p}$, this is another realizing quaternionic reflection, but it is not a member of the un-conjugated arc $T_\theta$.


\begin{thebibliography}{99}\small
\bibitem{Kuranishi1951} M. Kuranishi, \textit{On Everywhere Dense Imbedding of Free Groups in Lie Groups}, Nagoya Math. J. \textbf{2}, 63--71 (1951).
\bibitem{AlbuquerqueSilvaLeite1989} H. Albuquerque and F. Silva Leite, \textit{On the Generators of Semisimple Lie Algebras}, Linear Algebra Appl. \textbf{119}, 51--56 (1989).
\bibitem{Bois2009} J.-M. Bois, \textit{Generators of Simple Lie Algebras in Arbitrary Characteristics}, Math. Z. \textbf{262}, 715--741 (2009); arXiv:0708.1711.
\bibitem{DetinkoDeGraaf2020} A. S. Detinko and W. A. de Graaf, \textit{2-Generation of Simple Lie Algebras and Free Dense Subgroups of Algebraic Groups}, J. Algebra \textbf{545}, 159--173 (2020); doi:10.1016/j.jalgebra.2019.06.012.
\bibitem{BauerLevaillantFreedman2014} B. Bauer, C. Levaillant, and M. Freedman, \textit{Universality of Single Quantum Gates}, arXiv:1404.7822 (2014).
\bibitem{Chirvasitu2021} A. Chirvasitu, \textit{Large Sets of Generating Tuples for Lie Groups}, arXiv:2106.11955 (2021).
\bibitem{HarveyLawson1982} R. Harvey and H. B. Lawson, Jr., \textit{Calibrated Geometries}, Acta Math. \textbf{148}, 47--157 (1982).
\bibitem{Cacciatori2005} S. L. Cacciatori, B. L. Cerchiai, A. Della Vedova, G. Ortenzi, and A. Scotti, \textit{Euler Angles for $G_2$}, J. Math. Phys. \textbf{46}, 083512 (2005); arXiv:hep-th/0503106.
\bibitem{KnarrStroppel2025} N. Knarr and M. J. Stroppel, \textit{Subalgebras of Octonion Algebras}, J. Algebra \textbf{664}, 42--74 (2025); doi:10.1016/j.jalgebra.2024.10.004; arXiv:2303.00335.
\bibitem{DraperMartin2025} C. Draper and C. Mart\'in-Gonz\'alez, \textit{A Perspective on Totally Geodesic Submanifolds of the Symmetric Space $G_2/\SO(4)$}, Trans. Amer. Math. Soc. \textbf{378}(12), 8689--8721 (2025); doi:10.1090/tran/9479; arXiv:2504.07586.
\bibitem{Baez2002} J. C. Baez, \textit{The Octonions}, Bull. Amer. Math. Soc. \textbf{39}, 145--205 (2002); arXiv:math/0105155.
\bibitem{SpringerVeldkamp2000} T. A. Springer and F. D. Veldkamp, \textit{Octonions, Jordan Algebras and Exceptional Groups}, Springer, Berlin (2000).
\bibitem{ChemtovKarigiannis2022} M. Chemtov and S. Karigiannis, \textit{Observations about the Lie Algebra $\mathfrak g_2\subset\mathfrak{so}(7)$, Associative $3$-Planes, and $\mathfrak{so}(4)$ Subalgebras}, Expo. Math. \textbf{40}, 845--869 (2022); arXiv:2209.10613.
\bibitem{RauschSlupinski2022} M. Rausch de Traubenberg and M. J. Slupinski, \textit{Commutation Relations of $\mathfrak g_2$ and the Incidence Geometry of the Fano Plane}, arXiv:2207.13946 (2022).
\bibitem{Dynkin1952b} E. B. Dynkin, \textit{Semisimple Subalgebras of Semisimple Lie Algebras}, Mat. Sbornik N.S. \textbf{30(72)}, 349--462 (1952); English transl., Amer. Math. Soc. Transl. Ser. 2 \textbf{6}, 111--244 (1957).
\bibitem{JurdjevicSussmann1972} V. Jurdjevic and H. J. Sussmann, \textit{Control Systems on Lie Groups}, J. Differential Equations \textbf{12}, 313--329 (1972).
\bibitem{SilvaLeiteCrouch1988} F. Silva Leite and P. E. Crouch, \textit{Controllability on Classical Lie Groups}, Math. Control Signals Systems \textbf{1}, 31--42 (1988).
\bibitem{AgrachevSachkov2004} A. A. Agrachev and Y. L. Sachkov, \textit{Control Theory from the Geometric Viewpoint}, Springer, Berlin (2004).
\bibitem{ZeierSchulteHerbrueggen2011} R. Zeier and T. Schulte-Herbr\"uggen, \textit{Symmetry Principles in Quantum Systems Theory}, J. Math. Phys. \textbf{52}, 113510 (2011); doi:10.1063/1.3657939; arXiv:1012.5256.
\bibitem{BochnakCosteRoy1998} J. Bochnak, M. Coste, and M.-F. Roy, \textit{Real Algebraic Geometry}, Ergebnisse der Mathematik und ihrer Grenzgebiete (3), vol. 36, Springer, Berlin (1998); doi:10.1007/978-3-662-03718-8.
\bibitem{WeiNorman1964} J. Wei and E. Norman, \textit{On Global Representations of the Solutions of Linear Differential Equations as a Product of Exponentials}, Proc. Amer. Math. Soc. \textbf{15}, 327--334 (1964).
\bibitem{Altafini2003} C. Altafini, \textit{Parameter Differentiation and Quantum State Decomposition for Time Varying Schr\"odinger Equations}, Rep. Math. Phys. \textbf{52}(3), 381--400 (2003); doi:10.1016/S0034-4877(03)80037-X; arXiv:quant-ph/0201034.
\bibitem{Hall2000} B. C. Hall, \textit{An Elementary Introduction to Groups and Representations}, arXiv:math-ph/0005032 (2000).
\end{thebibliography}
\end{document}
